\section{Section graded ring infrastructure: tensor powers and graded sections}
\label{mk-chap:Picard_SectionGradedRing}

\providecommand{\PshMod}{\mathrm{PresheafOfModules}}
\providecommand{\ShMod}{\mathrm{SheafOfModules}}
\providecommand{\Modules}{\mathrm{Modules}}

The Hilbert-polynomial development of \ref{mk-chap:Picard_QuotScheme} represents the
degree-\(m\) graded component of a fibre by the global sections
\(\Gamma(X_s, \mathcal{F}_s \otimes L_s^{\otimes m})\) of a twist of a coherent
sheaf by a tensor power of a line bundle, and it organises these components into
a graded ring \(R(X_s, L_s) = \bigoplus_{m \ge 0} \Gamma(X_s, L_s^{\otimes m})\)
(\ref{mk-def:sectionGradedRing}) and a graded module
\(M(X_s, \mathcal{F}_s, L_s) = \bigoplus_{m \ge 0}
\Gamma(X_s, \mathcal{F}_s \otimes L_s^{\otimes m})\)
(\ref{mk-def:sectionGradedModule}). Three pieces of that picture have no direct
Mathlib counterpart and must be built here, each in its own section:

\begin{enumerate}
  \item \textbf{Tensor powers of a sheaf of modules}
    (\ref{mk-sec:sgr_tensor_powers}). Mathlib equips the category
    \(\PshMod\) of \emph{presheaves} of modules over a presheaf of commutative
    rings with a symmetric monoidal structure
    (\ref{mk-lem:presheafModule_monoidal_mathlib}) and provides the
    sheafification functor \(\PshMod \to \ShMod\)
    (\ref{mk-lem:presheafModule_sheafification_mathlib}); but the category
    \(\ShMod\) of \emph{sheaves} of modules carries \emph{no} monoidal
    structure. The tensor product of sheaves of modules
    \(\mathcal{F} \otimes_{\mathcal{O}_X} \mathcal{G}\) must therefore be built
    as the sheafification of the pointwise tensor presheaf, and the iterated
    tensor power \(\mathcal{L}^{\otimes m}\) recursively from it.
  \item \textbf{Lax-monoidal global sections} (\ref{mk-sec:sgr_lax_sections}).
    The global-sections functor \(\Gamma\) is only \emph{lax} monoidal: from the
    pointwise-strict monoidality of the presheaf tensor product and the
    sheafification unit one obtains a natural multiplication
    \(\Gamma(\mathcal{F}) \otimes_{\Gamma(\mathcal{O}_X)} \Gamma(\mathcal{G})
    \to \Gamma(\mathcal{F} \otimes_{\mathcal{O}_X} \mathcal{G})\), together with
    the associativity and unit coherence those products must satisfy.
  \item \textbf{Graded-ring / graded-module assembly}
    (\ref{mk-sec:sgr_graded_assembly}). The section multiplications and the
    tensor-power comparison isomorphisms assemble, through Mathlib's
    external-direct-sum graded API (\ref{mk-lem:directSum_gsemiring_mathlib},
    \ref{mk-lem:directSum_gcommSemiring_mathlib}, \ref{mk-lem:directSum_gmodule_mathlib}),
    into the graded \emph{ring} structure on
    \(\bigoplus_{m \ge 0} \Gamma(X_s, L_s^{\otimes m})\) and the graded
    \emph{module} structure on
    \(\bigoplus_{m \ge 0} \Gamma(X_s, \mathcal{F}_s \otimes L_s^{\otimes m})\).
    For a \emph{general} sheaf \(L_s\) the ring is only a non-commutative graded
    semiring (the free tensor algebra on \(\Gamma(L_s)\),
    \ref{mk-lem:sectionGradedRing_gsemiring}); it becomes \emph{commutative} exactly
    when \(L_s\) is an invertible sheaf / line bundle (\ref{mk-def:isInvertible}),
    which is the case relevant to the application
    (\ref{mk-lem:sectionGradedRing_gcommSemiring}). Stacks defines
    \(\Gamma_*(X,\mathcal{L})\) only for invertible \(\mathcal{L}\) for this reason.
\end{enumerate}

\section{Tensor powers of a sheaf of modules}
\label{mk-sec:sgr_tensor_powers}

We work over a fixed scheme \(X\) (in the application \(X = X_s\), a fibre over a
field \(\kappa(s)\)). Its category of sheaves of modules is
\(X.\Modules = \ShMod(\mathcal{O}_X)\), with global-sections notation
\(\Gamma(X, \mathcal{M}) = \mathcal{M}(X)\) for the value of the underlying
presheaf on the top open. The structure sheaf \(\mathcal{O}_X\) is a sheaf of
commutative rings, so its underlying presheaf factors through the category of
commutative rings; this is the hypothesis under which Mathlib's presheaf-level
monoidal structure applies.

\begin{lemma}[Monoidal structure on presheaves of modules]
  \label{mk-lem:presheafModule_monoidal_mathlib}
  \lean{PresheafOfModules.monoidalCategory}
  \mathlibok
  \dcref{custom}

  Let \(R : \mathcal{C}^{\mathrm{op}} \to \mathrm{CommRingCat}\) be a presheaf of
  commutative rings. Then the category \(\PshMod(R)\) of presheaves of modules
  over \(R\) carries a symmetric monoidal structure whose tensor product is
  computed objectwise: for presheaves of modules \(\mathcal{M}_1, \mathcal{M}_2\)
  and an object \(U\),
  \[
    (\mathcal{M}_1 \otimes \mathcal{M}_2)(U)
      \;=\; \mathcal{M}_1(U) \otimes_{R(U)} \mathcal{M}_2(U),
  \]
  with restriction maps the tensor products of the restriction maps. The unit is
  the presheaf \(R\) viewed as a module over itself, and the associator,
  unitors, and braiding are induced objectwise from those of
  \(\mathrm{ModuleCat}\). In particular evaluation at a fixed object \(U\) is a
  \emph{strict} monoidal functor
  \(\PshMod(R) \to \mathrm{ModuleCat}(R(U))\).
\end{lemma}

\begin{lemma}[Sheafification of presheaves of modules]
  \label{mk-lem:presheafModule_sheafification_mathlib}
  \lean{PresheafOfModules.sheafification}
  \mathlibok
  \dcref{custom}

  Let \(R\) be a sheaf of (commutative) rings on a site, with underlying
  presheaf \(R_0\). There is a sheafification functor
  \[
    (-)^{\#} \;:\; \PshMod(R_0) \longrightarrow \ShMod(R),
  \]
  left adjoint to the forgetful functor, together with the unit
  \(\eta_{\mathcal{M}_0} : \mathcal{M}_0 \to (\mathcal{M}_0)^{\#}\) of the
  adjunction. On underlying sheaves of abelian groups it agrees with the
  ordinary sheafification, and it sends a presheaf of \(R_0\)-modules to the
  associated sheaf of \(R\)-modules.
\end{lemma}

\begin{lemma}[The structure sheaf as the unit module]
  \label{mk-lem:moduleUnit_mathlib}
  \lean{SheafOfModules.unit}
  \mathlibok
  \dcref{custom}

  For a sheaf of rings \(R\) (in the application \(R = \mathcal{O}_X\)) the
  structure sheaf, viewed as a module over itself, is a distinguished object
  \(\mathbf{1} = \mathcal{O}_X \in \ShMod(\mathcal{O}_X)\). It is the value at
  the zeroth tensor power, \(\mathcal{L}^{\otimes 0} = \mathcal{O}_X\), and its
  global sections \(\Gamma(X, \mathcal{O}_X)\) form the degree-zero component of
  the section graded ring.
\end{lemma}


\begin{definition}[Unit module of a scheme]\label{mk-def:unitModule}
  \lean{AlgebraicGeometry.Scheme.Modules.unitModule}
  \leanok
  \dcref{custom}

  \uses{mk-lem:moduleUnit_mathlib}
  For a scheme \(X\), the \emph{unit module} \(\mathbf{1}_X \in X.\Modules\) is the
  structure sheaf viewed as a module over itself, i.e.\ Mathlib's
  \(\mathrm{SheafOfModules.unit}\) (\ref{mk-lem:moduleUnit_mathlib}) for
  \(\mathcal{O}_X\). It is the value of the zeroth tensor power,
  \(\mathcal{L}^{\otimes 0} = \mathbf{1}_X\) (\ref{mk-def:sheafTensorPow}).
\end{definition}

\begin{definition}[Tensor product of sheaves of modules]\label{mk-def:sheafTensorObj}
  \lean{AlgebraicGeometry.Scheme.Modules.sheafTensorObj}
  \leanok
  \dcref{01CA,custom}

  \uses{mk-lem:presheafModule_monoidal_mathlib,
    mk-lem:presheafModule_sheafification_mathlib}
  \textit{Source: \cite{stacks-project}, "Sheaves of Modules", Tag 01CA (tensor
  product).}
  Let \(\mathcal{F}, \mathcal{G}\) be sheaves of \(\mathcal{O}_X\)-modules. Their
  \emph{tensor product} is the sheaf of \(\mathcal{O}_X\)-modules
  \[
    \mathcal{F} \otimes_{\mathcal{O}_X} \mathcal{G}
      \;=\; \bigl(\mathcal{F} \otimes_{p,\mathcal{O}_X} \mathcal{G}\bigr)^{\#},
  \]
  the sheafification (\ref{mk-lem:presheafModule_sheafification_mathlib}) of the
  tensor product presheaf \(\mathcal{F} \otimes_{p,\mathcal{O}_X} \mathcal{G}\),
  whose value on an open \(U\) is the \(\mathcal{O}_X(U)\)-module
  \(\mathcal{F}(U) \otimes_{\mathcal{O}_X(U)} \mathcal{G}(U)\). Concretely, the
  tensor product presheaf is the Mathlib presheaf-level monoidal product
  (\ref{mk-lem:presheafModule_monoidal_mathlib}) of the underlying presheaves of
  modules of \(\mathcal{F}\) and \(\mathcal{G}\), and
  \(\mathcal{F} \otimes_{\mathcal{O}_X} \mathcal{G}\) is its sheafification. This
  construction is symmetric and associative because the presheaf tensor product
  is, and sheafification preserves the coherence isomorphisms.
\end{definition}

\begin{definition}[Tensor power of a line bundle]\label{mk-def:sheafTensorPow}
  \lean{AlgebraicGeometry.Scheme.Modules.tensorPow}
  \leanok
  \dcref{01CU,custom}

  \uses{mk-def:sheafTensorObj, mk-def:unitModule, mk-lem:moduleUnit_mathlib}
  \textit{Source: \cite{stacks-project}, "Sheaves of Modules", Tag 01CU (tensor
  powers).}
  Let \(\mathcal{L}\) be a sheaf of \(\mathcal{O}_X\)-modules (in the application
  an invertible sheaf / line bundle \(L_s\)). For \(m \in \mathbb{N}\) the
  \emph{\(m\)th tensor power} \(\mathcal{L}^{\otimes m}\) is defined by recursion
  on \(m\):
  \[
    \mathcal{L}^{\otimes 0} = \mathcal{O}_X
      \quad(\text{the unit module, \ref{mk-lem:moduleUnit_mathlib}}),
    \qquad
    \mathcal{L}^{\otimes (m+1)}
      = \mathcal{L}^{\otimes m} \otimes_{\mathcal{O}_X} \mathcal{L}
        \quad(\ref{mk-def:sheafTensorObj}).
  \]
  We restrict to non-negative powers, which is all the section graded ring uses;
  the Stacks definition extends \(\mathcal{L}^{\otimes n}\) to all
  \(n \in \mathbf{Z}\) using invertibility, but negative powers are not needed
  here. For \(\mathcal{L}\) invertible each \(\mathcal{L}^{\otimes m}\) is again
  invertible.
\end{definition}

\begin{definition}[Invertible sheaf of modules]\label{mk-def:isInvertible}
  \lean{AlgebraicGeometry.Scheme.Modules.IsInvertibleGr}
  \leanok
  \dcref{01CR,custom}

  \uses{mk-lem:tensorProductComm_eq_refl_mathlib}
  \textit{Source: \cite{stacks-project}, "Sheaves of Modules", Tag 01CR (invertible
  modules).}
  A sheaf of \(\mathcal{O}_X\)-modules \(\mathcal{L} \in X.\Modules\) is
  \emph{invertible} when it admits a \emph{trivializing basis}: an indexed family
  \(\{U_i\}_{i \in \iota}\) of opens of \(X\) whose range is a basis of the
  topology (\(\mathrm{Opens.IsBasis}\) of \(\{U_i\}\)) such that on each \(U_i\)
  the sections \(\Gamma(\mathcal{L}, U_i)\) form an \emph{invertible} module over
  the ring \(\Gamma(X, U_i)\), i.e.
  \(\mathrm{Module.Invertible}\,\Gamma(X,U_i)\,\Gamma(\mathcal{L},U_i)\)
  (\ref{mk-lem:tensorProductComm_eq_refl_mathlib}; equivalently
  \(\Gamma(\mathcal{L},U_i)\) is finite projective of constant rank one). This is
  the trivializing-cover / locally-free-of-rank-one form of invertibility: on each
  basis open \(\mathcal{L}\) restricts to a free rank-one module, hence
  \(\mathcal{L}|_{U_i} \cong \mathcal{O}_X|_{U_i}\). Over a scheme --- where every
  stalk \(\mathcal{O}_{X,x}\) is a local ring --- this is equivalent to the
  \(\otimes\)-invertible form (the existence of a sheaf \(\mathcal{N}\) with
  \(\mathcal{L} \otimes_{\mathcal{O}_X} \mathcal{N} \cong \mathbf{1}_X\),
  equivalently \(\mathcal{F} \mapsto \mathcal{L} \otimes_{\mathcal{O}_X}
  \mathcal{F}\) an equivalence of \(X.\Modules\)), Stacks Tag 01CR. The
  trivializing-basis form implies the triviality of the self-braiding
  \(\beta_{\mathcal{L},\mathcal{L}} = \mathrm{id}\)
  (\ref{mk-lem:braiding_eq_id_of_invertible}) by a basis-local argument; in
  particular \(\Gamma(X,\mathcal{L})\) itself need not be an invertible
  \(\Gamma(X,\mathcal{O}_X)\)-module. This is the line-bundle hypothesis under which
  the section graded ring becomes commutative
  (\ref{mk-lem:sectionGradedRing_gcommSemiring}). The structure sheaf \(\mathbf{1}_X\)
  is invertible (it trivializes on the basis of all opens), and any tensor power
  \(\mathcal{L}^{\otimes m}\) of an invertible sheaf is again invertible.
\end{definition}


\begin{definition}[Twist of a coherent sheaf by a tensor power]\label{mk-def:sheafModuleTwist}
  \lean{AlgebraicGeometry.Scheme.Modules.moduleTensorPow}
  \leanok
  \dcref{custom}

  \uses{mk-def:sheafTensorObj, mk-def:sheafTensorPow}
  Let \(\mathcal{F}\) be a sheaf of \(\mathcal{O}_X\)-modules (in the application
  a coherent sheaf \(\mathcal{F}_s\)) and \(\mathcal{L}\) a line bundle. For
  \(m \in \mathbb{N}\) the \emph{\(m\)-twist} of \(\mathcal{F}\) by
  \(\mathcal{L}\) is
  \[
    \mathcal{F}(m) \;:=\;
      \mathcal{F} \otimes_{\mathcal{O}_X} \mathcal{L}^{\otimes m}
      \qquad(\ref{mk-def:sheafTensorObj}, \ref{mk-def:sheafTensorPow}),
  \]
  the degree-\(m\) carrier of the section graded module. By construction
  \(\mathcal{F}(0) = \mathcal{F} \otimes_{\mathcal{O}_X} \mathcal{O}_X \cong
  \mathcal{F}\).
\end{definition}


The next four isomorphisms are the parts of the (would-be) symmetric monoidal
structure on \(X.\Modules\) that descend through sheafification from
\ref{mk-lem:presheafModule_monoidal_mathlib} using only \emph{functoriality} of the
sheafification functor (\ref{mk-def:schemeModuleSheafification}) and, for the
unitors, the reflective counit isomorphism --- no strong monoidality of
sheafification is required. They are the structural inputs of the tensor-power
comparison \ref{mk-lem:sheafTensorPow_add}.


The strong-monoidality comparison \ref{mk-lem:isIso_sheafification_whiskerRight_unit}
below is assembled from three reduction lemmas --- a localization criterion, the
local-isomorphism property of the unit, and their corollary --- together with the
coequalizer presentation of the relative module-presheaf tensor product, the single
Mathlib-absent brick on which the comparison rests.


\begin{definition}[Scheme-level module sheafification functor]\label{mk-def:schemeModuleSheafification}
  \lean{AlgebraicGeometry.Scheme.Modules.sheafification}
  \leanok
  \dcref{custom}

  \uses{mk-lem:presheafModule_sheafification_mathlib}
  For a scheme \(X\), the \emph{module sheafification functor}
  \[
    (-)^{\#} \;:\; X.\PshMod \longrightarrow X.\Modules
  \]
  is Mathlib's presheaf-of-modules sheafification
  (\ref{mk-lem:presheafModule_sheafification_mathlib}) for the identity morphism of
  the underlying presheaf of (commutative) rings of \(\mathcal{O}_X\), which is
  locally bijective. It is left adjoint to the forgetful functor
  \(X.\Modules \to X.\PshMod\), with adjunction unit
  \(\eta_{\mathcal{M}_0} : \mathcal{M}_0 \to (\mathcal{M}_0)^{\#}\) and an
  invertible counit (the forgetful functor is fully faithful). Being a functor it
  carries isomorphisms to isomorphisms, which is precisely what lets the
  presheaf-level coherence isomorphisms of
  \ref{mk-lem:presheafModule_monoidal_mathlib} descend to \(X.\Modules\).
\end{definition}

\begin{definition}[Left unitor of the sheaf tensor product]\label{mk-def:tensorObjUnitIso}
  \lean{AlgebraicGeometry.Scheme.Modules.tensorObjUnitIso}
  \leanok
  \dcref{custom}

  \uses{mk-def:sheafTensorObj, mk-def:unitModule, mk-def:schemeModuleSheafification,
    mk-lem:presheafModule_monoidal_mathlib}
  For a sheaf of \(\mathcal{O}_X\)-modules \(G\), the \emph{left unitor}
  \[
    \mathbf{1}_X \otimes_{\mathcal{O}_X} G \;\xrightarrow{\ \sim\ }\; G
  \]
  is obtained by descending the presheaf-level left unitor of
  \ref{mk-lem:presheafModule_monoidal_mathlib} through the sheafification functor
  (\ref{mk-def:schemeModuleSheafification}) and composing with the reflective counit
  isomorphism \((\mathcal{M}_0)^{\#} \cong \mathcal{M}\) of the sheafification
  adjunction. It is the base case \(m = 0\) of the tensor-power comparison
  \ref{mk-lem:sheafTensorPow_add}, and uses no monoidal structure on
  \(X.\Modules\) --- only functoriality of sheafification and the invertible
  counit.
\end{definition}

\begin{definition}[Braiding of the sheaf tensor product]\label{mk-def:tensorBraiding}
  \lean{AlgebraicGeometry.Scheme.Modules.tensorBraiding}
  \leanok
  \dcref{custom}

  \uses{mk-def:sheafTensorObj, mk-def:schemeModuleSheafification,
    mk-lem:presheafModule_monoidal_mathlib}
  For sheaves of \(\mathcal{O}_X\)-modules \(F, G\), the \emph{braiding}
  \[
    F \otimes_{\mathcal{O}_X} G \;\xrightarrow{\ \sim\ }\;
      G \otimes_{\mathcal{O}_X} F
  \]
  is the image under the sheafification functor
  (\ref{mk-def:schemeModuleSheafification}) of the symmetric braiding of the
  presheaf symmetric monoidal structure (\ref{mk-lem:presheafModule_monoidal_mathlib})
  on the underlying presheaves of modules. As pure sheafification-functoriality of
  an isomorphism it is again an isomorphism, requiring no monoidal structure on
  \(X.\Modules\). It is the symmetry used in the inductive step of
  \ref{mk-lem:sheafTensorPow_add}.
\end{definition}

\begin{lemma}[Sheafification inverts the whiskered localization unit]\label{mk-lem:isIso_sheafification_whiskerRight_unit}
  \lean{AlgebraicGeometry.Scheme.Modules.isIso_sheafification_whiskerRight_unit}
  \leanok
  \dcref{custom}

  \uses{mk-def:schemeModuleSheafification, mk-lem:presheafModule_monoidal_mathlib,
    mk-lem:presheafModule_sheafification_mathlib}
  Let \(P, Q\) be presheaves of \(\mathcal{O}_X\)-modules and let
  \(\eta_P : P \to P^{\#}\) be the unit of the sheafification adjunction
  (\ref{mk-def:schemeModuleSheafification},
  \ref{mk-lem:presheafModule_sheafification_mathlib}). Then the sheafification of the
  right-whiskered unit
  \[
    (\eta_P \rhd Q)^{\#} :
      (P \otimes_{p,\mathcal{O}_X} Q)^{\#} \longrightarrow
      (P^{\#} \otimes_{p,\mathcal{O}_X} Q)^{\#},
  \]
  where \(\otimes_{p,\mathcal{O}_X}\) is the presheaf monoidal product
  (\ref{mk-lem:presheafModule_monoidal_mathlib}), is an isomorphism of sheaves of
  modules. This is the strong-monoidality comparison of the module sheafification
  functor on a whiskered unit, the single Mathlib-absent brick on which the
  sheaf-level associator \ref{mk-cor:sheafTensorObjAssoc} and the tensor-power
  comparison \ref{mk-lem:sheafTensorPow_add} rest.
\end{lemma}

\begin{proof}
  \leanok
  isomorphism. The unit \(\eta_P\) is, on underlying abelian presheaves, the
  ordinary sheafification unit (the canonical map from a presheaf to its associated sheaf), which is a local
  isomorphism. Right-whiskering by \(Q\) preserves this: the relative
  module-presheaf tensor product is presented as a coequalizer, and tensoring with
  the fixed presheaf \(Q\) carries the local isomorphism \(\eta_P\) to a local
  isomorphism. Feeding this through the localization criterion gives the claim.
\end{proof}

\subsection{Monoidal structure by localization transport}
\label{mk-subsec:sgr_monoidal_localization}

The preceding isomorphisms (unitors, braiding, whiskered-unit comparison) are
exactly the data Mathlib needs to \emph{transport} the entire symmetric monoidal
structure from presheaves of modules onto \(X.\Modules\) along the sheafification
functor, rather than reassembling associator, pentagon, triangle and hexagon by
hand. The mechanism is \emph{monoidal localization}: a localization functor whose
inverted class of morphisms is compatible with the tensor product carries the
monoidal structure of its source to its target, with all coherence laws inherited.
The two Mathlib anchors below supply the general transport theorem and the fact
that module sheafification is such a localization; the two project blocks assemble
the compatibility hypothesis and read off the resulting structure.

\begin{lemma}[Monoidal localization transport]
  \label{mk-lem:monoidalLocalization_mathlib}
  \lean{CategoryTheory.LocalizedMonoidal}
  \mathlibok
  \dcref{custom}

  Let \(\mathcal{C}\) be a monoidal category and let \(W\) be a class of morphisms
  of \(\mathcal{C}\) satisfying \(\mathrm{MorphismProperty.IsMonoidal}\): \(W\) is
  multiplicative and stable under left- and right-whiskering by arbitrary objects.
  Let \(L : \mathcal{C} \to \mathcal{D}\) be a
  localization functor for \(W\). Then \(\mathcal{D}\) acquires a full monoidal
  category structure --- tensor product, associator and unitors --- for which the
  pentagon and triangle coherence identities hold, and \(L\) is strong monoidal,
  with comparison isomorphisms
  \[
    \mu_{X,Y} : L(X) \otimes L(Y) \;\xrightarrow{\ \sim\ }\; L(X \otimes Y).
  \]
  All of this structure is inherited from \(\mathcal{C}\) through \(L\). If moreover
  \(\mathcal{C}\) is braided (resp.\ symmetric), then \(\mathcal{D}\) inherits a
  braided (resp.\ symmetric) structure for which the hexagon identities hold and
  \(L\) is braided monoidal. The sheafification of presheaves over a site is the
  archetypal instance of this transport.
\end{lemma}

\begin{lemma}[Module sheafification is a monoidal localization input]
  \label{mk-lem:moduleSheafification_isLocalization_mathlib}
  \lean{PresheafOfModules.sheafification}
  \mathlibok
  \dcref{custom}

  The sheafification functor for presheaves of modules
  (\ref{mk-lem:presheafModule_sheafification_mathlib}) is a \emph{localization
  functor} for the morphism class
  \[
    W' \;=\; J.W.\mathrm{inverseImage}\,(\mathrm{toPresheaf}\,R_0),
  \]
  the class of morphisms of presheaves of \(R_0\)-modules whose underlying
  morphism of abelian-group presheaves lies in the local-isomorphism class
  \(J.W\) of the topology --- equivalently, the morphisms that sheafification
  sends to isomorphisms. Moreover the source category \(\PshMod(R_0)\) is
  symmetric monoidal (\ref{mk-lem:presheafModule_monoidal_mathlib}). Thus
  sheafification supplies precisely a localization functor \(L\) on a symmetric
  monoidal source, the two standing inputs of the transport theorem
  \ref{mk-lem:monoidalLocalization_mathlib}; the only remaining hypothesis is that
  \(W'\) be compatible with the tensor product. The \(\mathrm{IsLocalization}\)
  instance quoted above is for the bare Mathlib functor; the corresponding
  instance for the scheme-level sheafification functor
  \(\mathbf{(-)}^{\#}\), re-typed onto the monoidal carrier, is
  \ref{mk-def:sheafificationMon_isLocalization}.
\end{lemma}

\begin{definition}[Monoidal-carrier sheafification functor]\label{mk-def:sheafificationMon}
  \lean{AlgebraicGeometry.Scheme.Modules.sheafificationMon}
  \leanok
  \dcref{custom}

  \uses{mk-def:schemeModuleSheafification}
  For a scheme \(X\), \(\mathrm{sheafificationMon}\) is the module sheafification
  functor (\ref{mk-def:schemeModuleSheafification}) re-typed so that its source is
  the carrier \(\mathrm{MonoidalPresheaf}\,X\) of the presheaf-level symmetric
  monoidal structure (\ref{mk-lem:presheafModule_monoidal_mathlib}):
  \[
    \mathrm{sheafificationMon} \;:\;
      \mathrm{MonoidalPresheaf}\,X \longrightarrow X.\Modules .
  \]
  It is definitionally the same functor as \((-)^{\#}\); the only change is that
  its domain is presented as the carrier equipped with the symmetric monoidal
  structure of \ref{mk-lem:presheafModule_monoidal_mathlib}, so that the localization
  transport \ref{mk-lem:monoidalLocalization_mathlib} applies to it. This re-typing
  carries no mathematical content.
\end{definition}

\begin{definition}[The sheafification localization class \(W'\)]\label{mk-def:sheafificationW}
  \lean{AlgebraicGeometry.Scheme.Modules.sheafificationW}
  \leanok
  \dcref{custom}

  \uses{mk-lem:moduleSheafification_isLocalization_mathlib}
  For a scheme \(X\), \(\mathrm{sheafificationW}\) is the morphism class
  \[
    W' \;=\; J.W.\mathrm{inverseImage}\,(\mathrm{toPresheaf}\,R_0)
      \;\subseteq\; \mathrm{Mor}(\mathrm{MonoidalPresheaf}\,X),
  \]
  for \(R_0\) the underlying presheaf of \(\mathcal{O}_X\) and \(J\) the opens
  topology of \(X\): the morphisms of presheaves of modules whose underlying
  morphism of abelian-group presheaves is a local isomorphism, equivalently the
  morphisms inverted by sheafification
  (\ref{mk-lem:moduleSheafification_isLocalization_mathlib}). This is the class
  \(W'\) referred to in \ref{mk-def:sheafModule_W_isMonoidal} and
  \ref{mk-lem:moduleSheafification_isLocalization_mathlib}.
\end{definition}

\begin{lemma}[Sheafification is the localization at \(W'\) on the monoidal carrier]\label{mk-def:sheafificationMon_isLocalization}
  \lean{AlgebraicGeometry.Scheme.Modules.sheafificationMon_isLocalization}
  \leanok
  \dcref{custom}

  \uses{mk-def:sheafificationMon, mk-def:sheafificationW,
    mk-lem:moduleSheafification_isLocalization_mathlib}
  For a scheme \(X\), the functor \(\mathrm{sheafificationMon}\)
  (\ref{mk-def:sheafificationMon}) is a localization functor for the class
  \(W' = \mathrm{sheafificationW}\) (\ref{mk-def:sheafificationW}). The
  \(\mathrm{IsLocalization}\) property of \(\mathrm{sheafificationMon}\) is
  obtained by transporting the localization fact
  \ref{mk-lem:moduleSheafification_isLocalization_mathlib} across the definitional
  identification of \((-)^{\#}\) with the Mathlib sheafification functor and the
  re-typing of its domain onto the monoidal carrier \(\mathrm{MonoidalPresheaf}\,X\).
  It is the instance that allows the transport theorem
  \ref{mk-lem:monoidalLocalization_mathlib} to be applied to
  \(\mathrm{sheafificationMon}\).
\end{lemma}

\begin{definition}[Transport unit of the monoidal localization]\label{mk-def:localizationUnitIso}
  \lean{AlgebraicGeometry.Scheme.Modules.localizationUnitIso}
  \leanok
  \dcref{custom}

  \uses{mk-def:unitModule, mk-def:sheafificationMon, mk-def:schemeModuleSheafification}
  For a scheme \(X\), \(\mathrm{localizationUnitIso}\) is the isomorphism
  \[
    \varepsilon \;:\;
      \mathrm{sheafificationMon}(\mathbf{1}_{\mathrm{MonoidalPresheaf}\,X})
      \;\xrightarrow{\ \sim\ }\; \mathbf{1}_X
  \]
  identifying the image under \(\mathrm{sheafificationMon}\)
  (\ref{mk-def:sheafificationMon}) of the presheaf-level monoidal unit with the
  unit module \(\mathbf{1}_X\) (\ref{mk-def:unitModule}). Since the unit module is
  already a sheaf, this is the reflective counit isomorphism of the
  sheafification adjunction (\ref{mk-def:schemeModuleSheafification}) at the unit.
  It is the datum \(\varepsilon\) the transport theorem
  \ref{mk-lem:monoidalLocalization_mathlib} requires in order to fix the unit
  object of the transported monoidal structure on \(X.\Modules\) to be
  \(\mathbf{1}_X\).
\end{definition}

\begin{definition}[Tensor-compatibility of the sheafification localization class]\label{mk-def:sheafModule_W_isMonoidal}
  \lean{AlgebraicGeometry.Scheme.Modules.W_isMonoidal}
  \leanok
  \dcref{custom}

  \uses{mk-lem:isIso_sheafification_whiskerRight_unit,
    mk-lem:presheafModule_monoidal_mathlib,
    mk-lem:moduleSheafification_isLocalization_mathlib,
    mk-def:sheafificationW}
  For a scheme \(X\), write \(R_0\) for the underlying presheaf of the structure
  sheaf \(\mathcal{O}_X\) and
  \(W' = J.W.\mathrm{inverseImage}\,(\mathrm{toPresheaf}\,R_0)\) for the
  sheafification localization class of
  \ref{mk-lem:moduleSheafification_isLocalization_mathlib}. Then \(W'\) satisfies
  \(\mathrm{MorphismProperty.IsMonoidal}\): it is multiplicative and stable under
  left- and right-whiskering. This is the compatibility hypothesis that, together
  with \ref{mk-lem:moduleSheafification_isLocalization_mathlib}, feeds the transport
  theorem \ref{mk-lem:monoidalLocalization_mathlib}.
\end{definition}

\begin{proof}
  \leanok
  \ref{mk-lem:isIso_sheafification_whiskerRight_unit} --- stated there for the unit,
  but established for an arbitrary morphism inverted by sheafification --- shows
  that \(f \rhd Q\) is again inverted by sheafification, hence lies in \(W'\); the
  bridge between ``lies in \(W'\)'' and ``sheafification inverts'' is the
  localization criterion for the sheafification functor. Left-whiskering stability
  follows by conjugating the right-whiskering argument with the symmetric braiding of
  the presheaf monoidal structure (\ref{mk-lem:presheafModule_monoidal_mathlib}):
  since \(Q \lhd f\) is carried to \(f \rhd Q\) by the braiding isomorphism on
  both ends, and membership in \(W'\) is invariant under isomorphism of arrows,
  stability under left-whiskering reduces to stability under right-whiskering.
\end{proof}

\begin{definition}[Symmetric monoidal structure on sheaves of modules by transport]\label{mk-def:sheafModule_monoidalStructure}
  \lean{AlgebraicGeometry.Scheme.Modules.monoidalCategory}
  \leanok
  \dcref{custom}

  \uses{mk-def:sheafModule_W_isMonoidal, mk-lem:monoidalLocalization_mathlib, mk-def:unitModule,
    mk-def:sheafificationMon_isLocalization, mk-def:localizationUnitIso}
  For a scheme \(X\), the category \(X.\Modules\) of sheaves of
  \(\mathcal{O}_X\)-modules carries a symmetric monoidal structure obtained by
  transporting the symmetric monoidal structure of \(X.\PshMod\) along the
  sheafification functor. Concretely, one instantiates the monoidal localization
  transport \ref{mk-lem:monoidalLocalization_mathlib} with source the symmetric
  monoidal category \(X.\PshMod\), localization functor the module sheafification
  \ref{mk-lem:moduleSheafification_isLocalization_mathlib}, and the
  tensor-compatibility of its inverted class \ref{mk-def:sheafModule_W_isMonoidal};
  the unit is chosen to be the unit module \(\mathbf{1}_X\) (\ref{mk-def:unitModule}).
  The resulting tensor product agrees objectwise with the sheaf tensor product
  \ref{mk-def:sheafTensorObj}, via the strong-monoidal comparison
  \(\mu\) of \ref{mk-lem:monoidalLocalization_mathlib}. In particular the associator
  is the canonical (Mac Lane) associator, the pentagon and triangle laws hold, and
  --- the source being symmetric --- the braiding and its hexagon laws are
  inherited as well, giving \(X.\Modules\) the structure of a
  \emph{symmetric monoidal category}.
\end{definition}

\begin{proof}
  \leanok
  the source being symmetric --- a symmetric structure with the hexagon
  identities, all inherited through sheafification. What is inherited in this way
  is precisely the \emph{ambient} coherence of the monoidal and braided structure:
  the associator, the unitors, the braiding, and their pentagon, triangle, and
  hexagon laws, which hold for every pair of sheaves and require no separate
  induction. This ambient coherence is the raw material for the section-graded-ring
  assembly, but it is not the same as the multiplication identities of that ring:
  the \(\mu\)-comparison coherences \ref{mk-lem:sheafTensorPow_add} (associativity)
  and \ref{mk-lem:tensorPowAdd_comm} (commutativity) are established by induction on
  the tensor-power index --- the associativity pentagon and the commutativity
  hexagon respectively --- with the inherited coherence supplying only the
  per-step canonical isomorphism. The downstream section-ring identities are thus
  assembled inductively \emph{from} the inherited coherence, not read off it as a
  free consequence.
\end{proof}

\begin{corollary}
[Inherited braided and symmetric structure on sheaves of modules]
  \label{mk-cor:sheafModule_braided_symmetric}
  \lean{AlgebraicGeometry.Scheme.Modules.braidedCategory,
    AlgebraicGeometry.Scheme.Modules.symmetricCategory}
  \dcref{custom}

  \uses{mk-def:sheafModule_monoidalStructure, mk-lem:monoidalLocalization_mathlib,
    mk-lem:presheafModule_monoidal_mathlib}
  For a scheme \(X\), the transported monoidal structure on \(X.\Modules\)
  (\ref{mk-def:sheafModule_monoidalStructure}) refines to a braided and a symmetric
  monoidal structure. Because the presheaf source \(X.\PshMod\) is symmetric
  monoidal (\ref{mk-lem:presheafModule_monoidal_mathlib}), the braided- and
  symmetric-transport clauses of \ref{mk-lem:monoidalLocalization_mathlib} apply:
  the braiding of \(X.\Modules\) is the localization transport of the presheaf
  braiding, its hexagon identities are inherited, and the symmetry
  \(\beta_{F,G} \circ \beta_{G,F} = \mathrm{id}\) holds by transport. These are
  the intermediate \(\mathrm{BraidedCategory}\) and the final
  \(\mathrm{SymmetricCategory}\) instances on \(X.\Modules\) whose existence the
  prose of \ref{mk-def:sheafModule_monoidalStructure} asserts. What the inherited
  symmetric structure supplies for free --- for \emph{every} sheaf, with no
  hypothesis on \(\mathcal{L}\) --- is the braiding \(\beta\) together with its
  forward and reverse hexagon identities; this is exactly the input that
  \ref{mk-lem:tensorBraiding_hexagon_forward} descends to the tensor-power braiding
  and consumes.

  It does \emph{not}, by itself, make the commutativity identity of the section
  graded ring \ref{mk-lem:sectionMul_coherent} --- equivalently the comparison law
  \(\mu_{m,m'} = \mu_{m',m} \circ \beta\) of \ref{mk-lem:tensorPowAdd_comm} --- a
  free consequence. That identity additionally requires the self-braiding to be
  trivial, \(\beta_{\mathcal{L},\mathcal{L}} = \mathrm{id}\)
  (\ref{mk-lem:braiding_eq_id_of_invertible}), which holds precisely when
  \(\mathcal{L}\) is invertible (\ref{mk-def:isInvertible}). For a general
  \(\mathcal{L}\) the section ring is the free tensor algebra and commutativity
  \emph{fails}. So commutativity is invertibility-gated and hand-proved --- from
  the inherited hexagon together with \(\beta_{\mathcal{L},\mathcal{L}} =
  \mathrm{id}\) --- rather than inherited outright.
\end{corollary}

\begin{definition}[Bridge between the inherited and the project tensor product]\label{mk-def:tensorObjIso}
  \lean{AlgebraicGeometry.Scheme.Modules.tensorObjIso}
  \leanok
  \dcref{custom}

  \uses{mk-def:sheafTensorObj, mk-def:sheafModule_monoidalStructure,
    mk-lem:monoidalLocalization_mathlib}
  For sheaves of \(\mathcal{O}_X\)-modules \(F, G\), the \emph{tensor-agreement
  bridge} is the isomorphism
  \[
    F \otimes G \;\xrightarrow{\ \sim\ }\;
      F \otimes_{\mathcal{O}_X} G
  \]
  identifying the inherited tensor product \(F \otimes G\) of the transported
  monoidal structure (\ref{mk-def:sheafModule_monoidalStructure}) with the project's
  sheaf tensor product \(F \otimes_{\mathcal{O}_X} G\) (\ref{mk-def:sheafTensorObj}).
  It is built from the strong-monoidal comparison \(\mu\) of
  \ref{mk-lem:monoidalLocalization_mathlib} --- whose codomain
  \(L(a \otimes_{p} b)\) is exactly the sheafification of the presheaf tensor, i.e.\
  \(F \otimes_{\mathcal{O}_X} G\) --- precomposed with the reflective counit
  isomorphisms identifying \(F, G\) with the sheafifications of their underlying
  presheaves. This bridge is the structural device along which the canonical
  associator \(\alpha\), braiding \(\beta\) and unitors \(\lambda, \rho\) of the
  inherited symmetric monoidal structure are transported onto the project's
  \(\mathrm{tensorObj}\) family: each project-level coherence isomorphism is the
  conjugate, by \(\mathrm{tensorObjIso}\), of the corresponding canonical map of
  \(X.\Modules\). It is the launching point for the canonical reconstruction of
  the associator \ref{mk-cor:sheafTensorObjAssoc} and the tensor-power comparison
  \ref{mk-lem:sheafTensorPow_add}.
\end{definition}

\begin{corollary}[Associator of the sheaf tensor product]\label{mk-cor:sheafTensorObjAssoc}
  \lean{AlgebraicGeometry.Scheme.Modules.tensorObjAssoc}
  \leanok
  \dcref{custom}

  \uses{mk-def:sheafTensorObj, mk-def:sheafModule_monoidalStructure, mk-def:tensorObjIso}
  For sheaves of \(\mathcal{O}_X\)-modules \(A, B, C\) there is a canonical
  associativity isomorphism
  \[
    (A \otimes_{\mathcal{O}_X} B) \otimes_{\mathcal{O}_X} C
      \;\xrightarrow{\ \sim\ }\;
      A \otimes_{\mathcal{O}_X} (B \otimes_{\mathcal{O}_X} C)
  \]
  of sheaves of \(\mathcal{O}_X\)-modules.
\end{corollary}

\begin{proof}
  \leanok
  the bridge. The required isomorphism
  \[
    (A \otimes_{\mathcal{O}_X} B) \otimes_{\mathcal{O}_X} C
      \;\xrightarrow{\ \sim\ }\;
      A \otimes_{\mathcal{O}_X} (B \otimes_{\mathcal{O}_X} C)
  \]
  is obtained by conjugating the canonical associator
  \(\alpha_{A,B,C} : (A \otimes B) \otimes C \xrightarrow{\sim}
  A \otimes (B \otimes C)\) of \(X.\Modules\) by the bridge on each bracketing:
  one transports both iterated project tensors to the corresponding iterated
  inherited tensors using \(\Phi\) (whiskered appropriately on the outer and inner
  factors), applies \(\alpha\), and transports back. As a composite of
  isomorphisms it is an isomorphism, and being the conjugate of the canonical
  associator it is itself canonical --- no hand-rolled double-braiding detour and
  no separate strong-monoidality bricks are involved. Its pentagon and triangle
  coherence are then inherited from those of \(\alpha\)
  (\ref{mk-def:sheafModule_monoidalStructure}), which is exactly what
  \ref{mk-lem:sheafTensorPow_add} consumes.
\end{proof}


\subsection{Bridge lemmas to the canonical coherence maps}

The associator of \ref{mk-cor:sheafTensorObjAssoc} is obtained by conjugating the
canonical Mac Lane associator of \(X.\Modules\) along the tensor-agreement bridge
\(\mathrm{tensorObjIso}\) (\ref{mk-def:tensorObjIso}). The remaining coherence data
of the project's \(\mathrm{tensorObj}\) family --- the unitors, the braiding, and
the two whiskerings --- admit the same description. The bridge lemmas of this
subsection make that description explicit: each hand-built project-level
coherence isomorphism is shown to equal the conjugate, by
\(\mathrm{tensorObjIso}\), of the corresponding canonical map of the inherited
symmetric monoidal structure. They are the coherence layer along
which the inherited triangle, pentagon and hexagon are transported onto
\(\otimes_{\mathcal{O}_X}\); downstream, the \(\mathrm{tensorObjIso}\) bridges
introduced by adjacent factors cancel, leaving the canonical coherence laws to do
the work. Throughout, write \(\Phi_{F,G} : F \otimes G \xrightarrow{\sim}
F \otimes_{\mathcal{O}_X} G\) for the bridge \(\mathrm{tensorObjIso}\)
(\ref{mk-def:tensorObjIso}).

\begin{definition}[Right unitor of the sheaf tensor product]\label{mk-def:tensorObjRightUnitor}
  \lean{AlgebraicGeometry.Scheme.Modules.tensorObjRightUnitor}
  \leanok
  \dcref{custom}

  \uses{mk-def:sheafTensorObj, mk-def:unitModule, mk-def:schemeModuleSheafification,
    mk-lem:presheafModule_monoidal_mathlib}
  For a sheaf of \(\mathcal{O}_X\)-modules \(G\), the \emph{right unitor}
  \[
    G \otimes_{\mathcal{O}_X} \mathbf{1}_X \;\xrightarrow{\ \sim\ }\; G
  \]
  is obtained by descending the presheaf-level right unitor of
  \ref{mk-lem:presheafModule_monoidal_mathlib} through the sheafification functor
  (\ref{mk-def:schemeModuleSheafification}) and composing with the reflective counit
  isomorphism \((\mathcal{M}_0)^{\#} \cong \mathcal{M}\) of the sheafification
  adjunction. Like the left unitor (\ref{mk-def:tensorObjUnitIso}) it uses no
  monoidal structure on \(X.\Modules\) --- only functoriality of sheafification
  and the invertible counit --- and it is the right-unitality datum consumed by
  the section-level coherence of \ref{mk-lem:sectionMul_coherent}.
\end{definition}

\begin{lemma}[The right unitor is the bridge-conjugate of the canonical right unitor]\label{mk-lem:tensorObjRightUnitor_eq}
  \lean{AlgebraicGeometry.Scheme.Modules.tensorObjRightUnitor_eq}
  \leanok
  \dcref{custom}

  \uses{mk-def:tensorObjRightUnitor, mk-def:tensorObjIso,
    mk-def:sheafModule_monoidalStructure}
  For a sheaf of \(\mathcal{O}_X\)-modules \(G\), the hand-built right unitor
  (\ref{mk-def:tensorObjRightUnitor}) equals the conjugate of the canonical right
  unitor \(\rho\) of the inherited monoidal structure
  (\ref{mk-def:sheafModule_monoidalStructure}) by the bridge:
  \[
    \mathrm{tensorObjRightUnitor}_G
      \;=\; \rho_G \circ \Phi_{G,\mathbf{1}}^{-1},
  \]
  where \(\rho_G : G \otimes \mathbf{1} \xrightarrow{\sim} G\) and
  \(\Phi_{G,\mathbf{1}}^{-1} : G \otimes_{\mathcal{O}_X} \mathbf{1}_X
  \xrightarrow{\sim} G \otimes \mathbf{1}\).
\end{lemma}

\begin{proof}
  \(\rho\) exhibits exactly the descended presheaf right unitor and counit that
  define \(\mathrm{tensorObjRightUnitor}_G\), after the bridge
  \(\Phi_{G,\mathbf{1}}\) is moved across by naturality of \(\rho\). Cancelling the
  bridge on the unit factor gives the stated equality.
\end{proof}

\begin{lemma}[The left unitor is the bridge-conjugate of the canonical left unitor]\label{mk-lem:tensorObjUnitIso_eq}
  \lean{AlgebraicGeometry.Scheme.Modules.tensorObjUnitIso_eq}
  \leanok
  \dcref{custom}

  \uses{mk-def:tensorObjUnitIso, mk-def:tensorObjIso,
    mk-def:sheafModule_monoidalStructure}
  For a sheaf of \(\mathcal{O}_X\)-modules \(G\), the hand-built left unitor
  (\ref{mk-def:tensorObjUnitIso}) equals the conjugate of the canonical left unitor
  \(\lambda\) of the inherited monoidal structure
  (\ref{mk-def:sheafModule_monoidalStructure}) by the bridge:
  \[
    \mathrm{tensorObjUnitIso}_G
      \;=\; \lambda_G \circ \Phi_{\mathbf{1},G}^{-1},
  \]
  with \(\lambda_G : \mathbf{1} \otimes G \xrightarrow{\sim} G\) and
  \(\Phi_{\mathbf{1},G}^{-1} : \mathbf{1}_X \otimes_{\mathcal{O}_X} G
  \xrightarrow{\sim} \mathbf{1} \otimes G\). It is an auxiliary lemma used to put
  the \(m = 0\) base case of \ref{mk-lem:sheafTensorPow_add} on canonical footing.
\end{lemma}

\begin{proof}
  transported across by naturality of \(\lambda\) and cancelled on the unit
  factor.
\end{proof}

\begin{lemma}[The braiding is the bridge-conjugate of the canonical braiding]\label{mk-lem:tensorBraiding_eq}
  \lean{AlgebraicGeometry.Scheme.Modules.tensorBraiding_eq}
  \leanok
  \dcref{custom}

  \uses{mk-def:tensorBraiding, mk-def:tensorObjIso,
    mk-cor:sheafModule_braided_symmetric, mk-def:sheafModule_monoidalStructure}
  For sheaves of \(\mathcal{O}_X\)-modules \(F, G\), the hand-built braiding
  (\ref{mk-def:tensorBraiding}) equals the conjugate of the canonical braiding
  \(\beta\) of the inherited symmetric monoidal structure
  (\ref{mk-cor:sheafModule_braided_symmetric}) by the bridge on each side:
  \[
    \mathrm{tensorBraiding}_{F,G}
      \;=\; \Phi_{G,F} \circ \beta_{F,G} \circ \Phi_{F,G}^{-1},
  \]
  where \(\beta_{F,G} : F \otimes G \xrightarrow{\sim} G \otimes F\),
  \(\Phi_{F,G}^{-1} : F \otimes_{\mathcal{O}_X} G \xrightarrow{\sim} F \otimes G\)
  and \(\Phi_{G,F} : G \otimes F \xrightarrow{\sim} G \otimes_{\mathcal{O}_X} F\).
  It is an auxiliary lemma and the template for the two whiskering bridges
  below.
\end{lemma}

\begin{proof}
  rewrites \(\Phi_{G,F} \circ \beta_{F,G} \circ \Phi_{F,G}^{-1}\) as the descended
  presheaf braiding, which is \(\mathrm{tensorBraiding}_{F,G}\) by definition
  (\ref{mk-def:tensorBraiding}). The two bridges sit on opposite ends and are
  absorbed by this naturality rather than cancelling against each other.
\end{proof}

\begin{lemma}[Left whiskering is the bridge-conjugate of the canonical left whiskering]\label{mk-lem:tensorObjWhiskerLeftIso_eq}
  \lean{AlgebraicGeometry.Scheme.Modules.tensorObjWhiskerLeftIso_eq}
  \leanok
  \dcref{custom}

  \uses{mk-def:tensorObjIso, mk-def:sheafModule_monoidalStructure}
  Let \(F\) be a sheaf of \(\mathcal{O}_X\)-modules and \(e : G \xrightarrow{\sim}
  H\) an isomorphism of such sheaves. The hand-built left-whiskered isomorphism
  \(\mathrm{tensorObjWhiskerLeftIso}\,F\,e : F \otimes_{\mathcal{O}_X} G
  \xrightarrow{\sim} F \otimes_{\mathcal{O}_X} H\) equals the conjugate of the
  canonical left whiskering by the bridge:
  \[
    \mathrm{tensorObjWhiskerLeftIso}\,F\,e
      \;=\; \Phi_{F,H} \circ (F \lhd e) \circ \Phi_{F,G}^{-1},
  \]
  where \(F \lhd e : F \otimes G \xrightarrow{\sim} F \otimes H\) is the canonical
  left whiskering of \(X.\Modules\) (\ref{mk-def:sheafModule_monoidalStructure}).
\end{lemma}

\begin{proof}
  \(\Phi_{F,H} \circ (F \lhd e) \circ \Phi_{F,G}^{-1}\) into the descended
  whiskering, which is \(\mathrm{tensorObjWhiskerLeftIso}\,F\,e\).
\end{proof}

\begin{lemma}[Right whiskering is the bridge-conjugate of the canonical right whiskering]\label{mk-lem:tensorObjWhiskerRightIso_eq}
  \lean{AlgebraicGeometry.Scheme.Modules.tensorObjWhiskerRightIso_eq}
  \leanok
  \dcref{custom}

  \uses{mk-def:tensorObjIso, mk-def:sheafModule_monoidalStructure}
  Let \(e : G \xrightarrow{\sim} H\) be an isomorphism of sheaves of
  \(\mathcal{O}_X\)-modules and \(F\) a further such sheaf. The hand-built
  right-whiskered isomorphism \(\mathrm{tensorObjWhiskerRightIso}\,e\,F :
  G \otimes_{\mathcal{O}_X} F \xrightarrow{\sim} H \otimes_{\mathcal{O}_X} F\)
  equals the conjugate of the canonical right whiskering by the bridge:
  \[
    \mathrm{tensorObjWhiskerRightIso}\,e\,F
      \;=\; \Phi_{H,F} \circ (e \rhd F) \circ \Phi_{G,F}^{-1},
  \]
  where \(e \rhd F : G \otimes F \xrightarrow{\sim} H \otimes F\) is the canonical
  right whiskering of \(X.\Modules\) (\ref{mk-def:sheafModule_monoidalStructure}).
\end{lemma}

\begin{proof}
  whiskering \(\mathrm{tensorObjWhiskerRightIso}\,e\,F\).
\end{proof}


\begin{lemma}[Tensor-power comparison isomorphism]\label{mk-lem:sheafTensorPow_add}
  \lean{AlgebraicGeometry.Scheme.Modules.tensorPowAdd}
  \leanok
  \dcref{01CU,custom}

  \uses{mk-def:sheafTensorPow, mk-def:sheafTensorObj}
  \textit{Source: \cite{stacks-project}, "Sheaves of Modules", Tag 01CU (canonical
  power isomorphisms).}
  For a line bundle \(\mathcal{L}\) and \(m, m' \in \mathbb{N}\) there is a
  canonical isomorphism of sheaves of \(\mathcal{O}_X\)-modules
  \[
    \mu_{m,m'} :
    \mathcal{L}^{\otimes m} \otimes_{\mathcal{O}_X} \mathcal{L}^{\otimes m'}
    \;\xrightarrow{\ \sim\ }\; \mathcal{L}^{\otimes (m+m')},
  \]
  natural in \(\mathcal{L}\), and these isomorphisms satisfy a commutativity
  constraint (\(\mu_{m,m'}\) agrees with \(\mu_{m',m}\) after the braiding) and
  an associativity constraint (the two ways of combining
  \(\mathcal{L}^{\otimes m} \otimes \mathcal{L}^{\otimes m'} \otimes
  \mathcal{L}^{\otimes m''}\) into \(\mathcal{L}^{\otimes (m+m'+m'')}\) agree).
\end{lemma}

\begin{proof}
  \leanok
  All the structural isomorphisms used to assemble \(\mu_{m,m'}\) are the
  \emph{canonical} associator, braiding and unitors of the inherited symmetric
  monoidal structure on \(X.\Modules\) (\ref{mk-def:sheafModule_monoidalStructure},
  \ref{mk-cor:sheafModule_braided_symmetric}), transported to the project's tensor
  product \(\otimes_{\mathcal{O}_X}\) along the tensor-agreement bridge
  \(\mathrm{tensorObjIso}\) (\ref{mk-def:tensorObjIso}). In particular the associator
  is \ref{mk-cor:sheafTensorObjAssoc} --- itself the transported Mac Lane associator
  \(\alpha\) --- the braiding is the transport of the canonical \(\beta\)
  (\ref{mk-def:tensorBraiding}), and the left unitor is the transport of \(\lambda\)
  (\ref{mk-def:tensorObjUnitIso}). No hand-rolled associator and no separate
  strong-monoidality brick enters.

  \textbf{Construction of \(\mu_{m,m'}\).} Fix \(m\) and recurse on \(m'\),
  matching the right-growth of \(\mathrm{tensorPow}\) (\ref{mk-def:sheafTensorPow},
  \(\mathcal{L}^{\otimes(n+1)} = \mathcal{L}^{\otimes n} \otimes_{\mathcal{O}_X}
  \mathcal{L}\)) and the recursion of \(\mathbb{N}\)-addition on its \emph{second}
  argument. For \(m' = 0\), using \(\mathcal{L}^{\otimes 0} = \mathbf{1}_X\) and
  \(m + 0 = m\) (both \emph{definitional}),
  \[
    \mathcal{L}^{\otimes m} \otimes_{\mathcal{O}_X} \mathcal{L}^{\otimes 0}
      = \mathcal{L}^{\otimes m} \otimes_{\mathcal{O}_X} \mathbf{1}_X
      \xrightarrow{\ \rho\ } \mathcal{L}^{\otimes m}
  \]
  is the (transported) right unitor \(\rho\) of \ref{mk-def:tensorObjRightUnitor}.
  For the inductive step the recursive definition \(\mathcal{L}^{\otimes (c+1)} =
  \mathcal{L}^{\otimes c} \otimes_{\mathcal{O}_X} \mathcal{L}\)
  (\ref{mk-def:sheafTensorPow}) gives, with \(m + (c+1) = (m+c)+1\) again
  \emph{definitional},
  \[
    \mathcal{L}^{\otimes m} \otimes_{\mathcal{O}_X} \mathcal{L}^{\otimes (c+1)}
    = \mathcal{L}^{\otimes m} \otimes_{\mathcal{O}_X}
      (\mathcal{L}^{\otimes c} \otimes_{\mathcal{O}_X} \mathcal{L})
    \xrightarrow{\ \alpha^{-1}\ }
    (\mathcal{L}^{\otimes m} \otimes_{\mathcal{O}_X} \mathcal{L}^{\otimes c})
      \otimes_{\mathcal{O}_X} \mathcal{L}
    \xrightarrow{\ \mu_{m,c} \rhd \mathcal{L}\ }
    \mathcal{L}^{\otimes (m+c)} \otimes_{\mathcal{O}_X} \mathcal{L}
    = \mathcal{L}^{\otimes (m + (c+1))},
  \]
  using only the inverse associator \(\alpha^{-1}\) (\ref{mk-cor:sheafTensorObjAssoc})
  and the inductive comparison \(\mu_{m,c}\) whiskered by \(\mathcal{L}\) on the
  right; the final equality is the recursive definition of \(\mathrm{tensorPow}\)
  together with \(m + (c+1) = (m+c)+1\). The composite is \(\mu_{m,c+1}\), an
  isomorphism as a composite of isomorphisms. \emph{No braiding and no \(eqToIso\)
  reindexer enter the construction}: because both \(\mathrm{tensorPow}\) and
  \(\mathbb{N}\)-addition grow on the right, the new \(\mathcal{L}\)-factor already
  sits at the right edge of both source and target, so it never has to be
  transported across \(\mathcal{L}^{\otimes c}\). (This is the categorification of
  \(\mathrm{pow\_add}\), proved in Mathlib by induction on the second exponent with
  \(\mathrm{pow\_succ}\) and associativity only, no commutativity.) Naturality in
  \(\mathcal{L}\) holds because each constituent map is natural.

  \textbf{Associativity constraint.} Because \(\alpha\) and \(\lambda\) are the
  canonical coherence data of a genuine symmetric monoidal category on
  \(X.\Modules\) (\ref{mk-def:sheafModule_monoidalStructure},
  \ref{mk-cor:sheafModule_braided_symmetric}), the associativity diagram (the two
  bracketings of \(\mathcal{L}^{\otimes m} \otimes_{\mathcal{O}_X}
  \mathcal{L}^{\otimes m'} \otimes_{\mathcal{O}_X} \mathcal{L}^{\otimes m''} \to
  \mathcal{L}^{\otimes (m+m'+m'')}\)) is an instance of the pentagon identity of
  that inherited structure. It therefore holds by \emph{inheritance} --- Mac Lane
  coherence applied once, at the level of the symmetric monoidal category
  \(X.\Modules\) --- rather than by a separate induction over a hand-rolled
  associator. The bridge \(\mathrm{tensorObjIso}\) (\ref{mk-def:tensorObjIso})
  transports the coherence diagram from the inherited tensor \(\otimes\) to the
  project tensor \(\otimes_{\mathcal{O}_X}\) intact, so the associativity
  constraint passes to the family \(\{\mu_{m,m'}\}\) verbatim; this holds for an
  \emph{arbitrary} sheaf of modules \(\mathcal{L}\) and is \ref{mk-lem:tensorPowAdd_assoc}.

  \textbf{Commutativity constraint (for invertible \(\mathcal{L}\) only).} The
  commutativity square (\(\mu_{m,m'}\) versus \(\mu_{m',m}\) postcomposed with the
  braiding) does \emph{not} follow from Mac Lane coherence alone and requires
  \(\mathcal{L}\) to be invertible (\ref{mk-def:isInvertible}). After the hexagon and
  bridge reductions, the per-step inductive obligation reduces to a self-braiding
  \(\beta_{\mathcal{L},\mathcal{L}}\), which equals the identity precisely when
  \(\mathcal{L}\) is invertible (\ref{mk-lem:braiding_eq_id_of_invertible}); for a
  general sheaf of modules this identity fails. The full argument is
  \ref{mk-lem:tensorPowAdd_comm}.

  Concretely, the constituents \(\alpha^{-1}\) and \(\mu_{m,c} \rhd \mathcal{L}\)
  of the construction are rewritten as bridge-conjugates of the canonical associator
  and right-whiskering of \(X.\Modules\) by the whiskering bridge
  \ref{mk-lem:tensorObjWhiskerRightIso_eq} (the associator already canonical by
  \ref{mk-cor:sheafTensorObjAssoc}), so that the associativity diagram collapses to the
  canonical \emph{pentagon} of \ref{mk-def:sheafModule_monoidalStructure} --- with the
  \(\mathrm{tensorObjIso}\) bridges introduced by adjacent factors cancelling in
  pairs --- and is then dischargeable by the \(\mathrm{monoidal}\) coherence tactic.
  The commutativity diagram (for invertible \(\mathcal{L}\) only) additionally
  introduces the braiding via \ref{mk-lem:tensorBraiding_eq} and
  \ref{mk-lem:tensorObjWhiskerLeftIso_eq} and collapses to the canonical hexagon.
\end{proof}

\subsection{Coherence constraints of the tensor-power comparison}
\label{mk-subsec:sgr_tensorPowAdd_constraints}

The commutativity and associativity constraints of the comparison family
\(\{\mu_{m,m'}\}\) --- the content Stacks Tag 01CU leaves with ``formulation
omitted'' in \ref{mk-lem:sheafTensorPow_add} --- are decomposed here into named
iso-level identities, so each can be discharged against a single canonical
coherence law of the inherited symmetric monoidal structure on \(X.\Modules\)
(\ref{mk-def:sheafModule_monoidalStructure},
\ref{mk-cor:sheafModule_braided_symmetric}). After re-orienting \(\mu\) to recurse on
its \emph{second} index, the right-unit constraint
\ref{mk-lem:tensorPowAdd_zero_right} holds \emph{definitionally} --- it is the base
clause of that recursion --- so the worked templates are now the left-unit
constraint \ref{mk-lem:tensorPowAdd_zero_left} and the associativity constraint
\ref{mk-lem:tensorPowAdd_assoc}: each is an induction mirroring the recursion of
\(\mu\), whose per-step obligation collapses --- after the bridge lemmas of the
previous subsection rewrite every hand-built constituent to its canonical
counterpart and the \(\mathrm{tensorObjIso}\) bridges cancel in adjacent pairs ---
to a single canonical coherence law (the unitor coherence, resp.\ the pentagon) of
the inherited symmetric monoidal structure on \(X.\Modules\).

\begin{lemma}[Right-unit constraint for the tensor-power comparison]\label{mk-lem:tensorPowAdd_zero_right}
  \lean{AlgebraicGeometry.Scheme.Modules.tensorPowAdd_zero_right}
  \leanok
  \dcref{custom}

  \uses{mk-lem:sheafTensorPow_add, mk-def:tensorObjRightUnitor}
  For a sheaf of modules \(\mathcal{L}\) and \(n \in \mathbb{N}\), the degree-\((n,0)\)
  comparison isomorphism \(\mu_{n,0}\) (\ref{mk-lem:sheafTensorPow_add}) is the right
  unitor:
  \[
    \mu_{n,0} \;=\;
      \mathrm{tensorObjRightUnitor}_{\mathcal{L}^{\otimes n}}
    \;:\; \mathcal{L}^{\otimes n} \otimes_{\mathcal{O}_X} \mathbf{1}_X
    \;\xrightarrow{\ \sim\ }\; \mathcal{L}^{\otimes n}
  \]
  (\ref{mk-def:tensorObjRightUnitor}). Here \(\mathcal{L}^{\otimes 0} = \mathbf{1}_X\)
  and \(n + 0 = n\) hold on the nose, so the two sides inhabit a common type.
\end{lemma}

\begin{proof}
  \leanok
  holds definitionally, so no reindexing intervenes. Hence \(\mu_{n,0} = \rho\) for
  every \(n\), without induction.
\end{proof}

\begin{lemma}[Left-unit constraint for the tensor-power comparison]\label{mk-lem:tensorPowAdd_zero_left}
  \lean{AlgebraicGeometry.Scheme.Modules.tensorPowAdd_zero_left}
  \leanok
  \dcref{custom}

  \uses{mk-lem:sheafTensorPow_add, mk-def:tensorObjUnitIso, mk-lem:tensorObjUnitIso_eq,
    mk-lem:tensorObjRightUnitor_eq, mk-lem:tensorObjWhiskerRightIso_eq,
    mk-def:sheafModule_monoidalStructure}
  For a sheaf of modules \(\mathcal{L}\) and \(n \in \mathbb{N}\), the degree-\((0,n)\)
  comparison isomorphism \(\mu_{0,n}\) (\ref{mk-lem:sheafTensorPow_add}) is the left
  unitor, reindexed along \(0 + n = n\):
  \[
    \mu_{0,n} \;=\;
      \mathrm{tensorObjUnitIso}_{\mathcal{L}^{\otimes n}}
      \;:\; \mathbf{1}_X \otimes_{\mathcal{O}_X} \mathcal{L}^{\otimes n}
      \;\xrightarrow{\ \sim\ }\; \mathcal{L}^{\otimes n}
  \]
  (\ref{mk-def:tensorObjUnitIso}, transport of the left unitor \(\lambda\)). Unlike
  the right-unit case, this is \emph{not} the base clause of the second-index
  recursion, so it requires an induction on \(n\).
\end{lemma}

\begin{proof}
  \leanok
  \uses{mk-lem:sheafTensorPow_add, mk-def:tensorObjUnitIso, mk-lem:tensorObjUnitIso_eq,
    mk-lem:tensorObjRightUnitor_eq, mk-lem:tensorObjWhiskerRightIso_eq,
    mk-def:sheafModule_monoidalStructure}
  Induction on \(n\). For \(n = 0\), \(\mu_{0,0}\) is the base clause
  \(\mathrm{tensorObjRightUnitor}_{\mathbf{1}_X}\) (right unitor of the unit), and on
  the unit object the left and right unitors coincide
  (\(\mathrm{MonoidalCategory.unitors\_equal}\)); descending that equality through
  sheafification gives \(\mu_{0,0} = \mathrm{tensorObjUnitIso}\). For the inductive
  step \(n = c+1\), unfold \(\mu_{0,c+1}\) by the succ-clause
  (\ref{mk-lem:sheafTensorPow_add}), \(\alpha^{-1} \circ (\mu_{0,c} \rhd \mathcal{L})\),
  and substitute the inductive hypothesis \(\mu_{0,c} = \lambda\). After the bridges
  \ref{mk-lem:tensorObjUnitIso_eq}/\ref{mk-lem:tensorObjWhiskerRightIso_eq} rewrite to
  canonical maps and the \(\mathrm{tensorObjIso}\) bridges cancel in pairs, the step
  is the canonical \emph{triangle} identity \(\lambda_{A \otimes \mathcal{L}} =
  (\lambda_A \rhd \mathcal{L}) \circ \alpha^{-1}_{\mathbf{1}, A, \mathcal{L}}\) of the
  inherited monoidal structure (\ref{mk-def:sheafModule_monoidalStructure}), closed by
  the \(\mathrm{monoidal}\) coherence tactic.
\end{proof}

\begin{lemma}[Self-braiding of an invertible module is trivial (affine)]
  \label{mk-lem:tensorProductComm_eq_refl_mathlib}
  \lean{Module.Invertible.tensorProductComm_eq_refl}
  \mathlibok
  \dcref{custom}

  Let \(R\) be a commutative ring and \(M\) an invertible \(R\)-module
  (\(\mathrm{Module.Invertible}\,R\,M\): the evaluation map
  \(M \otimes_R M^{\ast} \to R\) is bijective, equivalently \(M\) is finite
  projective of constant rank one). Then the swap involution of the tensor square
  is the identity:
  \[
    \mathrm{TensorProduct.comm}_R\,M\,M \;=\;
      \mathrm{id}_{M \otimes_R M},
  \]
  i.e.\ \(x \otimes y = y \otimes x\) in \(M \otimes_R M\) for all
  \(x, y \in M\). This is the affine kernel of self-braiding triviality. Note the
  abstract monoidal statement ``a \(\otimes\)-invertible object has
  \(\beta = \mathrm{id}\)'' is \emph{false} --- an odd line in super vector spaces
  is \(\otimes\)-invertible yet has \(\beta = -\mathrm{id}\) --- so the result
  genuinely uses the concrete swap on a module that is invertible in the
  \(\mathrm{Module.Invertible}\) sense, not mere categorical invertibility.
\end{lemma}

\begin{lemma}[Trivial self-braiding of an invertible sheaf]\label{mk-lem:braiding_eq_id_of_invertible}
  \lean{AlgebraicGeometry.Scheme.Modules.tensorBraiding_self_eq_id_of_isInvertible}
  \leanok
  \dcref{custom}

  \uses{mk-def:isInvertible, mk-def:tensorBraiding, mk-def:sheafTensorObj,
    mk-def:sheafTensorPow, mk-def:schemeModuleSheafification,
    mk-lem:tensorProductComm_eq_refl_mathlib,
    mk-lem:presheafModule_monoidal_mathlib}
  Let \(\mathcal{L}\) be an invertible sheaf (\ref{mk-def:isInvertible}) on a scheme
  \(X\). Then the braiding of \(\mathcal{L}\) with itself is the identity:
  \[
    \beta_{\mathcal{L},\mathcal{L}} \;=\;
      \mathrm{id}_{\mathcal{L} \otimes_{\mathcal{O}_X} \mathcal{L}}
    \qquad(\ref{mk-def:tensorBraiding}).
  \]
  More generally each tensor power \(\mathcal{L}^{\otimes m}\) is again invertible
  (\ref{mk-def:sheafTensorPow}, \ref{mk-def:isInvertible}), so its self-braiding is
  trivial as well, \(\beta_{\mathcal{L}^{\otimes m},\mathcal{L}^{\otimes m}} =
  \mathrm{id}\). This is the single arithmetic input distinguishing the invertible
  case: it holds precisely because \(\mathcal{L}\) is locally free of rank one, and
  fails for a general sheaf of modules --- already for
  \(\mathcal{O}_X^{\oplus 2}\) over a point, where
  \(e_1 \otimes e_2 \neq e_2 \otimes e_1\).
\end{lemma}

\begin{proof}
  \leanok
  \uses{mk-def:isInvertible, mk-def:tensorBraiding, mk-def:sheafTensorPow,
    mk-def:schemeModuleSheafification, mk-lem:tensorProductComm_eq_refl_mathlib,
    mk-lem:presheafModule_monoidal_mathlib}
  Equality of morphisms of sheaves of modules is a \emph{local} condition that can
  be tested on a basis of opens; we obtain
  \(\beta_{\mathcal{L},\mathcal{L}} = \mathrm{id}\) by descent from the trivializing
  basis carried by the invertibility datum.

  \emph{The trivializing basis.} By \ref{mk-def:isInvertible}, invertibility of
  \(\mathcal{L}\) supplies an indexed family \(\{U_i\}\) of opens whose range is a
  basis of the topology of \(X\), together with, on each \(U_i\), an
  invertible-module structure
  \(\mathrm{Module.Invertible}\,\Gamma(X,U_i)\,\Gamma(\mathcal{L},U_i)\)
  (\ref{mk-lem:tensorProductComm_eq_refl_mathlib}).

  \emph{The presheaf braiding is the identity on each basis open.} By construction
  (\ref{mk-def:tensorBraiding}) the braiding \(\beta_{\mathcal{L},\mathcal{L}}\) is
  the image under sheafification (\ref{mk-def:schemeModuleSheafification}) of the
  presheaf braiding \(\beta^{\mathrm{pre}}\), whose component on an open \(U\) is the
  swap \(\mathrm{TensorProduct.comm}_{\Gamma(X,U)}\,\Gamma(\mathcal{L},U)\,
  \Gamma(\mathcal{L},U)\) (\ref{mk-lem:presheafModule_monoidal_mathlib}). On a basis
  open \(U_i\) the module \(\Gamma(\mathcal{L},U_i)\) is invertible over
  \(\Gamma(X,U_i)\), so that swap is the identity by
  \ref{mk-lem:tensorProductComm_eq_refl_mathlib}; hence \(\beta^{\mathrm{pre}}\)
  agrees with \(\mathrm{id}\) on every basis open \(U_i\).

  \emph{Descent to the sheafification.} Write \(T = \mathcal{L} \otimes
  \mathcal{L}\) for the underlying presheaf tensor product and
  \(\eta : T \to T^{\#}\) for the sheafification unit, so that
  \(\beta_{\mathcal{L},\mathcal{L}}\) is the sheafified endomorphism of \(T^{\#}\)
  induced by \(\beta^{\mathrm{pre}}\). Its target \(T^{\#}\) is a sheaf, so two
  morphisms into it that agree on a basis of opens coincide (separatedness, the
  basis form of sheaf-morphism extensionality, applied to the underlying sheaf of
  abelian groups). The two presheaf morphisms
  \(\beta^{\mathrm{pre}} \circ \eta\) and \(\eta\) (both \(T \to T^{\#}\)) agree on
  every basis open \(U_i\), because there \(\beta^{\mathrm{pre}}\) is the identity;
  therefore \(\beta^{\mathrm{pre}} \circ \eta = \eta\). By naturality of \(\eta\),
  \[
    \eta \circ (\text{sheafification of } \beta^{\mathrm{pre}})
      = \beta^{\mathrm{pre}} \circ \eta = \eta = \eta \circ \mathrm{id}_{T^{\#}}.
  \]
  Precomposition with the unit \(\eta\) is injective on morphisms into a sheaf (the
  universal property of the sheafification adjunction,
  \ref{mk-def:schemeModuleSheafification}), so the sheafified map equals
  \(\mathrm{id}_{T^{\#}}\). An isomorphism whose forward map is the identity is the
  identity isomorphism, whence \(\beta_{\mathcal{L},\mathcal{L}} = \mathrm{id}\).
  Crucially the check is performed on the trivializing basis and \emph{never} at the
  global open \(\top\), where \(\Gamma(X,\mathcal{L})\) need not be invertible.

  The generalization to \(\mathcal{L}^{\otimes m}\) is the identical argument: a
  tensor power of an invertible sheaf is again invertible
  (\ref{mk-def:sheafTensorPow}), so it carries a trivializing basis on which the
  presheaf braiding is the identity, and the same descent gives
  \(\beta_{\mathcal{L}^{\otimes m},\mathcal{L}^{\otimes m}} = \mathrm{id}\).
\end{proof}

\begin{lemma}[First-index successor recursion of the tensor-power comparison]\label{mk-lem:tensorPowAdd_succ_left}
  \lean{AlgebraicGeometry.Scheme.Modules.tensorPowAdd_succ_left}
  \leanok
  \dcref{custom}

  \uses{mk-lem:sheafTensorPow_add, mk-lem:tensorPowAdd_assoc, mk-cor:sheafTensorObjAssoc,
    mk-lem:tensorObjWhiskerLeftIso_eq, mk-lem:tensorObjWhiskerRightIso_eq,
    mk-lem:tensorObjUnitIso_eq, mk-def:tensorObjUnitIso, mk-def:sheafTensorPow,
    mk-lem:tensorPowAdd_zero_right}
  For an \emph{arbitrary} sheaf of \(\mathcal{O}_X\)-modules \(\mathcal{L}\) (no
  invertibility needed) and \(c, m \in \mathbb{N}\), the comparison
  \[
    \mu_{c+1,m} : \mathcal{L}^{\otimes(c+1)} \otimes_{\mathcal{O}_X}
      \mathcal{L}^{\otimes m} \;\xrightarrow{\ \sim\ }\;
      \mathcal{L}^{\otimes((c+1)+m)}
  \]
  is recovered from the lower comparison \(\mu_{c,\,1+m}\) by peeling the
  freshly-added factor to the left:
  \[
    \mu_{c+1,m}
      \;=\;
      \mu_{c,\,1+m}
      \;\circ\;
      \bigl(\mathcal{L}^{\otimes c} \lhd \nu_m\bigr)
      \;\circ\;
      \alpha_{\mathcal{L}^{\otimes c},\,\mathcal{L},\,\mathcal{L}^{\otimes m}},
  \]
  where \(\alpha\) is the associator (\ref{mk-cor:sheafTensorObjAssoc}), \(\lhd\)
  the canonical left whiskering, and
  \(\nu_m : \mathcal{L} \otimes_{\mathcal{O}_X} \mathcal{L}^{\otimes m}
  \xrightarrow{\sim} \mathcal{L}^{\otimes(1+m)}\) is \(\mu_{1,m}\) precomposed
  with the left-unit identification \(\mathcal{L} \cong
  \mathbf{1}_X \otimes_{\mathcal{O}_X}\mathcal{L} = \mathcal{L}^{\otimes 1}
  \otimes_{\mathcal{O}_X} \cdots\) (\ref{mk-def:tensorObjUnitIso}), modulo the
  reindexing \((c+1)+m = c+(1+m)\). This is the first-index dual of the
  definitional second-index successor clause of \(\mu\)
  (\ref{mk-lem:sheafTensorPow_add}); unlike the commutativity hexagon it is
  \emph{braiding-free} and holds for every \(\mathcal{L}\), both sides realising
  the identity permutation of the \((c+1)+m\) identical \(\mathcal{L}\)-factors.
  This order-preserving form is correct, reusable infrastructure, but it is
  \emph{not} the ingredient consumed by the commutativity proof: the glue of
  \ref{mk-lem:tensorPowAdd_comm} requires the order-\emph{reversing} recursion
  \ref{mk-lem:tensorPowAdd_succ_left_braided} (brick~1\('\)), which exposes the lower
  comparison \(\mu_{c,m}\) rather than the unreachable \(\mu_{c,1+m}\).
\end{lemma}

\begin{proof}
  \leanok
  \uses{mk-lem:tensorPowAdd_assoc, mk-cor:sheafTensorObjAssoc, mk-lem:sheafTensorPow_add,
    mk-lem:tensorObjWhiskerLeftIso_eq, mk-lem:tensorObjWhiskerRightIso_eq,
    mk-lem:tensorObjUnitIso_eq, mk-lem:tensorPowAdd_zero_right}
  Instantiate the associativity constraint \ref{mk-lem:tensorPowAdd_assoc} at the
  indices \((m,m',m'') = (c,1,m)\), which relates the comparison
  \(\mu_{(c+1),m} = \mu_{c+1,m}\) to lower data:
  \[
    (\mu_{c,1} \rhd \mathcal{L}^{\otimes m}) \circ \mu_{c+1,m}
      \circ (\text{reindex})
    \;=\;
    \alpha_{\mathcal{L}^{\otimes c},\mathcal{L}^{\otimes 1},
        \mathcal{L}^{\otimes m}} \circ
      (\mathcal{L}^{\otimes c} \lhd \mu_{1,m}) \circ \mu_{c,\,1+m}.
  \]
  The factor \(\mu_{c,1}\) is the second-index successor of \(\mu_{c,0}\), and by
  \ref{mk-lem:sheafTensorPow_add} together with \ref{mk-lem:tensorPowAdd_zero_right}
  (which makes \(\mu_{c,0}\) the right unitor) it is the canonical right-unit
  isomorphism \(\mathcal{L}^{\otimes c}\otimes_{\mathcal{O}_X}\mathcal{L}^{
  \otimes 1}\cong\mathcal{L}^{\otimes(c+1)}\); hence its right-whiskering
  \(\mu_{c,1}\rhd\mathcal{L}^{\otimes m}\) is invertible with a structural
  inverse. Solving the displayed identity for \(\mu_{c+1,m}\) and absorbing the
  unit isomorphism \(\mu_{c,1}\) and the \(\mathcal{L}^{\otimes 1}=\mathbf{1}_X
  \otimes_{\mathcal{O}_X}\mathcal{L}\) identification into the left-unit factor
  \(\nu_m\) --- using \ref{mk-lem:tensorObjUnitIso_eq} and the whisker bridges
  \ref{mk-lem:tensorObjWhiskerLeftIso_eq}, \ref{mk-lem:tensorObjWhiskerRightIso_eq} to
  pass to canonical whiskerings --- yields the stated formula. No braiding occurs:
  the entire computation stays inside the associator/unitor (pentagon--triangle)
  fragment, so once the hand-built whiskerings are rewritten to canonical ones it
  is closed by the canonical coherence tactic \(\mathrm{monoidal}\).
\end{proof}

\begin{lemma}[Descended forward hexagon for the sheaf braiding]\label{mk-lem:tensorBraiding_hexagon_forward}
  \lean{AlgebraicGeometry.Scheme.Modules.tensorBraiding_hexagon_forward}
  \leanok
  \dcref{custom}

  \uses{mk-def:tensorBraiding, mk-lem:tensorBraiding_eq, mk-cor:sheafTensorObjAssoc,
    mk-lem:tensorObjWhiskerLeftIso_eq, mk-lem:tensorObjWhiskerRightIso_eq,
    mk-cor:sheafModule_braided_symmetric, mk-def:tensorObjIso}
  For sheaves of \(\mathcal{O}_X\)-modules \(F, A, B\), the hand-built braiding
  (\ref{mk-def:tensorBraiding}) and associator (\ref{mk-cor:sheafTensorObjAssoc})
  satisfy the \emph{forward hexagon} identity
  \[
    \alpha_{A,B,F} \circ
    \beta_{F,\,A\otimes_{\mathcal{O}_X}B} \circ
    \alpha_{F,A,B}
    \;=\;
    (A \lhd \beta_{F,B}) \circ
    \alpha_{A,F,B} \circ
    (\beta_{F,A} \rhd B),
  \]
  where \(\beta\) is the braiding (\ref{mk-def:tensorBraiding}), \(\alpha\) the
  associator (\ref{mk-cor:sheafTensorObjAssoc}), and \(\lhd,\rhd\) the canonical
  whiskerings. In words: braiding \(F\) past the tensor \(A\otimes_{\mathcal{O}_X}
  B\) factors, up to associators, as braiding \(F\) past \(A\) and then past
  \(B\) separately. This is the symmetric-monoidal hexagon of the inherited
  structure (\ref{mk-cor:sheafModule_braided_symmetric}) read off at the project's
  tensor product, the braided analogue of the right-unit/braiding coherence used
  in the base case of \ref{mk-lem:tensorPowAdd_comm}.
\end{lemma}

\begin{proof}
  \leanok
  \uses{mk-lem:tensorBraiding_eq, mk-cor:sheafTensorObjAssoc,
    mk-lem:tensorObjWhiskerLeftIso_eq, mk-lem:tensorObjWhiskerRightIso_eq,
    mk-cor:sheafModule_braided_symmetric, mk-def:tensorObjIso}
  This is the descent of the canonical forward hexagon identity
  through the tensor-agreement bridge, exactly mirroring how the
  base-case coherence \(\rho = \beta_{G,\mathbf{1}}\circ\lambda\) was descended
  from the canonical braiding--left-unitor coherence. The inherited symmetric
  monoidal structure on \(X.\Modules\) (\ref{mk-cor:sheafModule_braided_symmetric})
  satisfies the canonical forward hexagon
  \[
    \alpha'_{A,B,F}\circ\beta'_{F,A\otimes B}\circ\alpha'_{F,A,B}
    = (A\lhd\beta'_{F,B})\circ\alpha'_{A,F,B}\circ(\beta'_{F,A}\rhd B)
  \]
  for its canonical braiding \(\beta'\) and associator \(\alpha'\). By the bridge
  lemmas --- \ref{mk-lem:tensorBraiding_eq} (the hand-built braiding is the
  conjugate of \(\beta'\) by the bridge \(\Phi\) of \ref{mk-def:tensorObjIso}),
  \ref{mk-cor:sheafTensorObjAssoc} (the associator is the conjugate of \(\alpha'\)),
  and the whisker bridges \ref{mk-lem:tensorObjWhiskerLeftIso_eq},
  \ref{mk-lem:tensorObjWhiskerRightIso_eq} --- every project-level construct is
  replaced by its canonical counterpart conjugated by \(\Phi\). The bridges
  introduced by adjacent factors telescope and cancel in inverse pairs (the
  cancellation recipe of \ref{mk-lem:tensorPowAdd_zero_right}), leaving exactly the
  canonical hexagon, which discharges the identity. The braiding is genuine and is
  \emph{not} closed by the \(\mathrm{monoidal}\) coherence tactic (which handles
  only associators and unitors); the explicit canonical hexagon is what is
  consumed.
\end{proof}

\begin{lemma}[Order-reversing first-index successor recursion (brick 1$'$)]\label{mk-lem:tensorPowAdd_succ_left_braided}
  \lean{AlgebraicGeometry.Scheme.Modules.tensorPowAdd_succ_left_braided}
  \leanok
  \dcref{custom}

  \uses{mk-lem:sheafTensorPow_add, mk-def:isInvertible,
    mk-lem:braiding_eq_id_of_invertible, mk-def:tensorBraiding,
    mk-cor:sheafTensorObjAssoc, mk-lem:tensorBraiding_eq,
    mk-lem:tensorBraiding_hexagon_forward, mk-lem:tensorObjWhiskerLeftIso_eq,
    mk-lem:tensorObjWhiskerRightIso_eq, mk-lem:tensorPowAdd_zero_right,
    mk-lem:tensorPowAdd_zero_left}
  Let \(\mathcal{L}\) be an \emph{invertible} sheaf (\ref{mk-def:isInvertible}) and
  \(c, m \in \mathbb{N}\). The first-index successor comparison
  \[
    \mu_{c+1,m} : \mathcal{L}^{\otimes(c+1)} \otimes_{\mathcal{O}_X}
      \mathcal{L}^{\otimes m} \;\xrightarrow{\ \sim\ }\;
      \mathcal{L}^{\otimes((c+1)+m)}
  \]
  is recovered from the \emph{lower} comparison \(\mu_{c,m}\) by braiding the
  freshly-added left factor \(\mathcal{L}\) past the block \(\mathcal{L}^{\otimes
  m}\) and applying \(\mu_{c,m}\) on the right, framed by associators:
  \[
    \mu_{c+1,m}
      \;=\;
      (\text{reindex})\circ
      (\mu_{c,m} \rhd \mathcal{L})\circ
      \alpha^{-1}_{\mathcal{L}^{\otimes c},\,\mathcal{L}^{\otimes m},\,\mathcal{L}}\circ
      (\mathcal{L}^{\otimes c} \lhd \beta_{\mathcal{L},\,\mathcal{L}^{\otimes m}})\circ
      \alpha_{\mathcal{L}^{\otimes c},\,\mathcal{L},\,\mathcal{L}^{\otimes m}},
  \]
  where \(\alpha\) is the associator (\ref{mk-cor:sheafTensorObjAssoc}), \(\beta\)
  the braiding (\ref{mk-def:tensorBraiding}), \(\lhd,\rhd\) the canonical
  whiskerings, and the reindex is the definitional identification
  \(\mathcal{L}^{\otimes(c+m)}\otimes_{\mathcal{O}_X}\mathcal{L} =
  \mathcal{L}^{\otimes((c+m)+1)} = \mathcal{L}^{\otimes((c+1)+m)}\). The maps are
  displayed in \(\circ\) (right-to-left) order; the construction chains them
  left-to-right (diagrammatic order). In contrast with the
  order-\emph{preserving} \ref{mk-lem:tensorPowAdd_succ_left} --- which is
  braiding-free but supplies the wrong atom for the commutativity glue --- this
  order-\emph{reversing} form surfaces exactly the lower comparison \(\mu_{c,m}\)
  demanded by the inductive hypothesis of \ref{mk-lem:tensorPowAdd_comm}, at the
  cost of a genuine braiding \(\beta_{\mathcal{L},\mathcal{L}^{\otimes m}}\);
  hence the invertibility hypothesis enters here.
\end{lemma}

\begin{proof}
  \leanok
  \uses{mk-lem:sheafTensorPow_add, mk-def:isInvertible,
    mk-lem:braiding_eq_id_of_invertible, mk-def:tensorBraiding,
    mk-cor:sheafTensorObjAssoc, mk-lem:tensorBraiding_eq,
    mk-lem:tensorBraiding_hexagon_forward, mk-lem:tensorObjWhiskerLeftIso_eq,
    mk-lem:tensorObjWhiskerRightIso_eq, mk-lem:tensorPowAdd_zero_right,
    mk-lem:tensorPowAdd_zero_left}
    Induction on \(m\).

  \emph{Base \(m = 0\).} Here \(\mathcal{L}^{\otimes 0} = \mathbf{1}_X\), so the
  only braiding occurring is \(\beta_{\mathcal{L},\mathbf{1}_X}\), the unit
  coherence --- no self-braiding \(\beta_{\mathcal{L},\mathcal{L}}\) appears and
  no invertibility is needed. Reading \(\mu_{c,0}\) and \(\mu_{c+1,0}\) as the
  appropriate unit isomorphisms (\ref{mk-lem:tensorPowAdd_zero_right},
  \ref{mk-lem:tensorPowAdd_zero_left}) and passing the hand-built braiding and
  whiskerings to their canonical counterparts (\ref{mk-lem:tensorBraiding_eq},
  \ref{mk-lem:tensorObjWhiskerLeftIso_eq}, \ref{mk-lem:tensorObjWhiskerRightIso_eq}),
  the residual identity lies entirely in the associator/unitor
  (pentagon--triangle) fragment and is closed by the canonical coherence tactic.

  \emph{Successor \(m \to m+1\).} Expand both \(\mu_{c,m+1}\) and
  \(\mu_{c+1,m+1}\) by the second-index successor clause
  (\ref{mk-lem:sheafTensorPow_add}), \(\mu_{-,m+1} = \alpha^{-1}\circ(\mu_{-,m}
  \rhd\mathcal{L})\). The braiding on the factored target,
  \(\beta_{\mathcal{L},\,\mathcal{L}^{\otimes(m+1)}} =
  \beta_{\mathcal{L},\,\mathcal{L}^{\otimes m}\otimes_{\mathcal{O}_X}\mathcal{L}}\),
  splits over \(\mathcal{L}^{\otimes(m+1)} = \mathcal{L}^{\otimes m}
  \otimes_{\mathcal{O}_X}\mathcal{L}\) by the canonical right-tensor braiding
  identity (the forward hexagon \ref{mk-lem:tensorBraiding_hexagon_forward} read on
  the right factor) into \((\beta_{\mathcal{L},\mathcal{L}^{\otimes m}}\rhd
  \mathcal{L})\) and \((\mathcal{L}^{\otimes m}\lhd\beta_{\mathcal{L},\mathcal{L}})\),
  framed by associators. Rewrite \(\beta_{\mathcal{L},\mathcal{L}^{\otimes m}}\) by
  the inductive hypothesis and collapse the emergent self-braiding
  \(\beta_{\mathcal{L},\mathcal{L}}\) to the identity by
  \ref{mk-lem:braiding_eq_id_of_invertible} (this is the sole use of
  invertibility). After the hand-built whiskerings are rewritten to canonical ones
  (\ref{mk-lem:tensorObjWhiskerLeftIso_eq},
  \ref{mk-lem:tensorObjWhiskerRightIso_eq}), the remaining structural
  associator/unitor obligation is closed by the canonical coherence tactic. This
  is the braided analogue of the pentagon successor step of
  \ref{mk-lem:tensorPowAdd_assoc}.
\end{proof}

\begin{lemma}[Commutativity constraint for the tensor-power comparison]\label{mk-lem:tensorPowAdd_comm}
  \lean{AlgebraicGeometry.Scheme.Modules.tensorPowAdd_comm}
  \leanok
  \dcref{custom}

  \uses{mk-lem:sheafTensorPow_add, mk-def:isInvertible, mk-def:tensorBraiding,
    mk-lem:braiding_eq_id_of_invertible, mk-cor:sheafTensorObjAssoc,
    mk-lem:tensorBraiding_eq, mk-lem:tensorObjWhiskerLeftIso_eq,
    mk-lem:tensorObjWhiskerRightIso_eq, mk-lem:tensorPowAdd_zero_right,
    mk-lem:tensorPowAdd_succ_left_braided, mk-lem:tensorBraiding_hexagon_forward}
  Let \(\mathcal{L}\) be an \emph{invertible} sheaf (\ref{mk-def:isInvertible}) and
  \(m, m' \in \mathbb{N}\). Then the comparison family is symmetric:
  \[
    \mu_{m,m'} \;=\;
      \mu_{m',m} \circ
      \beta_{\mathcal{L}^{\otimes m}, \mathcal{L}^{\otimes m'}},
  \]
  i.e.\ \(\mu_{m,m'}\) factors through the braiding
  \(\beta_{\mathcal{L}^{\otimes m},\mathcal{L}^{\otimes m'}} :
  \mathcal{L}^{\otimes m} \otimes_{\mathcal{O}_X} \mathcal{L}^{\otimes m'}
  \xrightarrow{\sim}
  \mathcal{L}^{\otimes m'} \otimes_{\mathcal{O}_X} \mathcal{L}^{\otimes m}\)
  (\ref{mk-def:tensorBraiding}) followed by \(\mu_{m',m}\), after the index
  reindexing \(m' + m = m + m'\). This is the commutativity constraint of
  \ref{mk-lem:sheafTensorPow_add}. The invertibility hypothesis is essential: for a
  general sheaf of modules the section ring is the free tensor algebra on
  \(\Gamma(X,\mathcal{L})\), and \(\mu_{m,m'}\) and \(\mu_{m',m} \circ \beta\)
  realize \emph{different} permutations of the \(m + m'\) identical
  \(\mathcal{L}\)-factors; already at \(m = m' = 1\) the identity reduces to
  \(\beta_{\mathcal{L},\mathcal{L}} = \mathrm{id}\), which holds exactly when
  \(\mathcal{L}\) is invertible (\ref{mk-lem:braiding_eq_id_of_invertible}).
\end{lemma}

\begin{proof}
  \uses{mk-lem:sheafTensorPow_add, mk-def:tensorBraiding,
    mk-lem:braiding_eq_id_of_invertible, mk-cor:sheafTensorObjAssoc,
    mk-lem:tensorBraiding_eq, mk-lem:tensorObjWhiskerLeftIso_eq,
    mk-lem:tensorObjWhiskerRightIso_eq, mk-lem:tensorPowAdd_zero_right,
    mk-lem:tensorPowAdd_zero_left, mk-lem:tensorPowAdd_succ_left_braided,
    mk-lem:tensorBraiding_hexagon_forward, mk-cor:sheafModule_braided_symmetric}
  Induction mirroring the recursion of \(\mu\) (now on the second index), so the
  statement relates two distinct inductions. The base cases \(m = 0\) and
  \(m' = 0\) are the unit constraints: the right unit
  \ref{mk-lem:tensorPowAdd_zero_right} and the left unit
  \ref{mk-lem:tensorPowAdd_zero_left} read off the two sides.

  \emph{Successor \(m' = c+1\).} The successor step is the symmetric-monoidal
  hexagon, assembled from the bricks above. Unfold the left-hand side by the
  second-index successor clause (\ref{mk-lem:sheafTensorPow_add}), \(\mu_{m,c+1} =
  \alpha^{-1}\circ(\mu_{m,c}\rhd\mathcal{L})\), and apply the inductive hypothesis
  \(\mu_{m,c} = \beta_{\mathcal{L}^{\otimes m},\mathcal{L}^{\otimes c}}\circ
  \mu_{c,m}\circ(\text{reindex})\); this exposes the lower comparison
  \(\mu_{c,m}\) precomposed with the braiding
  \(\beta_{\mathcal{L}^{\otimes m},\mathcal{L}^{\otimes c}}\). On the right-hand
  side the comparison \(\mu_{c+1,m}\) is a \emph{first-index} successor that the
  inductive hypothesis cannot reach (it supplies \(\mu_{c,m}\), not
  \(\mu_{c+1,m}\)). Split the right-hand braiding
  \(\beta_{\mathcal{L}^{\otimes m},\,\mathcal{L}^{\otimes(c+1)}}\) on
  \(\mathcal{L}^{\otimes(c+1)}=\mathcal{L}^{\otimes c}\otimes_{\mathcal{O}_X}
  \mathcal{L}\) by the descended forward hexagon
  \ref{mk-lem:tensorBraiding_hexagon_forward} and distribute the whiskers; the
  shared prefix \(\alpha^{-1}\circ(\beta_{\mathcal{L}^{\otimes m},
  \mathcal{L}^{\otimes c}}\rhd\mathcal{L})\) cancels off both sides. The residual
  obligation is then \emph{exactly} the order-reversing recursion of
  \(\mu_{c+1,m}\), namely brick~1\('\) (\ref{mk-lem:tensorPowAdd_succ_left_braided}),
  whose right-hand factor \(\beta_{\mathcal{L},\mathcal{L}^{\otimes m}}\) meets the
  surviving \(\beta_{\mathcal{L}^{\otimes m},\mathcal{L}}\) coming from the
  hexagon. Substituting brick~1\('\) and passing the hand-built braiding and
  whiskerings to canonical maps (\ref{mk-lem:tensorBraiding_eq},
  \ref{mk-lem:tensorObjWhiskerLeftIso_eq},
  \ref{mk-lem:tensorObjWhiskerRightIso_eq}), the two opposed braidings annihilate by
  the symmetry relation \(\beta_{\mathcal{L}^{\otimes m},\mathcal{L}}\circ
  \beta_{\mathcal{L},\mathcal{L}^{\otimes m}} = \mathrm{id}\)
  (\ref{mk-cor:sheafModule_braided_symmetric}); the only remaining self-braiding
  \(\beta_{\mathcal{L},\mathcal{L}}\) on a doubled \(\mathcal{L}\)-factor is the
  identity by \ref{mk-lem:braiding_eq_id_of_invertible}. With both braidings
  discharged and the \(\mathrm{tensorObjIso}\) bridges cancelled pairwise
  (\ref{mk-lem:tensorPowAdd_zero_right}), the residue is a pure
  associator/unitor identity, closed by the canonical coherence
  (\(\mathrm{monoidal}\)); the reindex \(m'+m=m+m'\) matches the two ends.
\end{proof}

\begin{lemma}[Associativity constraint for the tensor-power comparison]\label{mk-lem:tensorPowAdd_assoc}
  \lean{AlgebraicGeometry.Scheme.Modules.tensorPowAdd_assoc}
  \leanok
  \dcref{custom}

  \uses{mk-lem:sheafTensorPow_add, mk-cor:sheafTensorObjAssoc,
    mk-def:sheafModule_monoidalStructure, mk-lem:tensorObjWhiskerRightIso_eq,
    mk-lem:tensorObjWhiskerLeftIso_eq, mk-lem:tensorPowAdd_zero_right}
  For an arbitrary \(\mathcal{L} : \mathbf{Mod}(\mathcal{O}_X)\) (no invertibility
  needed --- associativity holds for all \(\mathcal{L}\)) and \(m, m', m'' \in \mathbb{N}\)
  the two bracketings of
  \(\mathcal{L}^{\otimes m} \otimes_{\mathcal{O}_X} \mathcal{L}^{\otimes m'}
  \otimes_{\mathcal{O}_X} \mathcal{L}^{\otimes m''}\) into
  \(\mathcal{L}^{\otimes (m+m'+m'')}\) agree:
  \[
    \mu_{m+m', m''} \circ (\mu_{m,m'} \rhd \mathcal{L}^{\otimes m''})
      \;=\;
    \mu_{m, m'+m''} \circ (\mathcal{L}^{\otimes m} \lhd \mu_{m',m''})
      \circ \alpha_{\mathcal{L}^{\otimes m}, \mathcal{L}^{\otimes m'},
        \mathcal{L}^{\otimes m''}},
  \]
  where \(\alpha\) is the associator \ref{mk-cor:sheafTensorObjAssoc} and \(\lhd,
  \rhd\) the canonical whiskerings, modulo the reindexings
  \((m+m')+m'' = m+(m'+m'')\). This is the associativity constraint of
  \ref{mk-lem:sheafTensorPow_add}.
\end{lemma}

\begin{proof}
  \uses{mk-lem:sheafTensorPow_add, mk-cor:sheafTensorObjAssoc,
    mk-def:sheafModule_monoidalStructure, mk-lem:tensorObjWhiskerRightIso_eq,
    mk-lem:tensorObjWhiskerLeftIso_eq, mk-lem:tensorPowAdd_zero_right}
  Induction on the \emph{last} index \(m''\), mirroring the second-index recursion
  of \(\mu = \mathrm{tensorPowAdd}\) (\ref{mk-lem:sheafTensorPow_add}). Because that
  recursion is now braiding-free --- each clause is built from \(\alpha^{-1}\) and a
  single right-whisker --- both sides of the identity are composites of canonical
  associators and whiskerings only; \emph{no braiding occurs anywhere}, so the whole
  statement is an instance of Mac~Lane's pentagon coherence and is dischargeable by
  the \(\mathrm{monoidal}\) coherence tactic once the project's hand-built isos are
  rewritten to their canonical counterparts by \ref{mk-lem:tensorObjWhiskerRightIso_eq}
  (the associator is already canonical by \ref{mk-cor:sheafTensorObjAssoc}).

  \emph{Base case \(m'' = 0\).} Both \(\mu_{m+m',0}\) and \(\mu_{m',0}\),
  \(\mu_{m'+0,\dots}\) collapse to right unitors (\ref{mk-lem:tensorPowAdd_zero_right},
  which is now \(\mathrm{rfl}\)) and \(\mathcal{L}^{\otimes 0} = \mathbf{1}\), \(\,\cdot
  + 0 = \cdot\) hold definitionally; the two sides reduce to the unit triangle
  \(\rho_{A\otimes B} = (A \lhd \rho_B)\circ\dots\), closed by \(\mathrm{monoidal}\).

  \emph{Inductive step \(m'' = c+1\).} Unfold \(\mu_{-,\,c+1}\) on both sides by the
  second-index succ-clause \(\alpha^{-1}\circ(\mu_{-,c}\rhd\mathcal{L})\)
  (\ref{mk-lem:sheafTensorPow_add}). The outer right-whisker distributes over the
  composite and, after rewriting by \ref{mk-lem:tensorObjWhiskerRightIso_eq}, the
  inductive hypothesis (the \(m''=c\) instance, whiskered \(\rhd\mathcal{L}\)) folds
  into the central atom via one pure-associator slide
  \(\mathrm{associator\_inv\_naturality\_left}\) (\ref{mk-cor:sheafTensorObjAssoc}) ---
  NOT a braiding. The residual structural glue between the four folded comparison
  atoms and the associator \(\alpha_{\mathcal{L}^{\otimes m},\mathcal{L}^{\otimes
  m'},\mathcal{L}^{\otimes m''}}\) is a pure pentagon on canonical associators and
  whiskerings, closed by \(\mathrm{monoidal}\). The reindexings \((m+m')+c =
  m+(m'+c)\) etc.\ are \(\mathrm{Nat.add}\)-definitional except for one
  \(\mathrm{add\_assoc}\) transport in the frozen statement, dispatched by a
  \(\mathrm{subst}\) helper exactly as before.
\end{proof}

\section{Lax-monoidal global sections}
\label{mk-sec:sgr_lax_sections}

The graded multiplication is supplied by the lax-monoidal structure of the
global-sections functor \(\Gamma(X, -) : \ShMod(\mathcal{O}_X) \to
\mathrm{ModuleCat}(\Gamma(X, \mathcal{O}_X))\). It is only \emph{lax} because
global sections do not commute with sheafification: a global section of a
sheafified tensor product is not in general a finite sum of elementary tensors of
global sections.

\begin{definition}[Section multiplication]\label{mk-def:sectionMul}
  \lean{AlgebraicGeometry.Scheme.Modules.sectionsMul}
  \leanok
  \dcref{custom}

  \uses{mk-def:sheafTensorObj, mk-lem:presheafModule_monoidal_mathlib,
    mk-lem:presheafModule_sheafification_mathlib}
  Let \(\mathcal{F}, \mathcal{G}\) be sheaves of \(\mathcal{O}_X\)-modules. The
  \emph{section multiplication} is the \(\Gamma(X,\mathcal{O}_X)\)-bilinear map
  \[
    \Gamma(X, \mathcal{F}) \otimes_{\Gamma(X,\mathcal{O}_X)}
      \Gamma(X, \mathcal{G})
    \;\longrightarrow\;
    \Gamma(X, \mathcal{F} \otimes_{\mathcal{O}_X} \mathcal{G})
  \]
  defined as follows. A pair of global sections \((\sigma, \tau)\) determines,
  by the objectwise formula of \ref{mk-lem:presheafModule_monoidal_mathlib} at the
  top open, the elementary tensor \(\sigma \otimes \tau \in
  (\mathcal{F} \otimes_{p,\mathcal{O}_X} \mathcal{G})(X)\), a global section of
  the tensor product presheaf. Composing with the global-sections component of
  the sheafification unit
  \(\eta : \mathcal{F} \otimes_{p,\mathcal{O}_X} \mathcal{G} \to
  (\mathcal{F} \otimes_{p,\mathcal{O}_X} \mathcal{G})^{\#} =
  \mathcal{F} \otimes_{\mathcal{O}_X} \mathcal{G}\)
  (\ref{mk-lem:presheafModule_sheafification_mathlib},
  \ref{mk-def:sheafTensorObj}) yields a global section
  \(\sigma \cdot \tau \in
  \Gamma(X, \mathcal{F} \otimes_{\mathcal{O}_X} \mathcal{G})\). Bilinearity over
  \(\Gamma(X,\mathcal{O}_X)\) is inherited from bilinearity of the elementary
  tensor and \(\Gamma(X,\mathcal{O}_X)\)-linearity of \(\eta\).
\end{definition}

The coherence of the section multiplication is governed by how the canonical
coherence maps of \(X.\Modules\) interact with \(\mathrm{sectionsMul}\). Because
\(\mathrm{sectionsMul}\) is the image under \(\Gamma(X,-)\) of the sheafification
unit \(\eta\) evaluated at the top open, and \(\eta\) is natural, each canonical
coherence isomorphism --- the unitors, the braiding, the associator --- pushes
through \(\mathrm{sectionsMul}\) on the corresponding section factors. The four
naturality helpers below record these push-through identities at the level of
global sections; they are the section-level partners of the iso-level constraints
\ref{mk-lem:tensorPowAdd_zero_right}, \ref{mk-lem:tensorPowAdd_comm} and
\ref{mk-lem:tensorPowAdd_assoc}. The key local fact in each case is that the
\(\mathrm{MonoidalPresheaf}\) coherence map at the top open is the corresponding
\(\mathrm{ModuleCat}\) symmetric-monoidal structure map, computed on elementary
tensors of sections.

\begin{lemma}[Left unitor through the section multiplication]\label{mk-lem:tensorObjUnitIso_hom_sectionsMul}
  \lean{AlgebraicGeometry.Scheme.Modules.tensorObjUnitIso_hom_sectionsMul}
  \leanok
  \dcref{custom}

  \uses{mk-def:sectionMul, mk-def:tensorObjUnitIso}
  Let \(G\) be a sheaf of \(\mathcal{O}_X\)-modules and \(a \in \Gamma(X, G)\).
  Then
  \[
    \Gamma\bigl(\mathrm{tensorObjUnitIso}_G\bigr)
      \bigl(\mathrm{sectionsMul}_{\mathbf{1}_X, G}(1 \otimes a)\bigr) = a,
  \]
  i.e.\ applying global sections of the left unitor
  \(\mathbf{1}_X \otimes_{\mathcal{O}_X} G \xrightarrow{\sim} G\)
  (\ref{mk-def:tensorObjUnitIso}) to the section product of the unit section
  \(1 \in \Gamma(X,\mathcal{O}_X)\) and \(a\) (\ref{mk-def:sectionMul}) returns
  \(a\). This is the section-level reading of the left-unit law via
  \(\eta\)-naturality of the lax structure.
\end{lemma}

\begin{proof}
  \leanok
  unitor). Evaluating at the top open and using the \(\mathrm{ModuleCat}\)
  left-unitor formula \(r \otimes m \mapsto r \cdot m\) sends \(1 \otimes a\) to
  \(1 \cdot a = a\).
\end{proof}

\begin{lemma}[Right unitor through the section multiplication]\label{mk-lem:tensorObjRightUnitor_hom_sectionsMul}
  \lean{AlgebraicGeometry.Scheme.Modules.tensorObjRightUnitor_hom_sectionsMul}
  \leanok
  \dcref{custom}

  \uses{mk-def:sectionMul, mk-def:tensorObjRightUnitor}
  Let \(G\) be a sheaf of \(\mathcal{O}_X\)-modules and \(a \in \Gamma(X, G)\).
  Then
  \[
    \Gamma\bigl(\mathrm{tensorObjRightUnitor}_G\bigr)
      \bigl(\mathrm{sectionsMul}_{G, \mathbf{1}_X}(a \otimes 1)\bigr) = a,
  \]
  the right-unit analogue of \ref{mk-lem:tensorObjUnitIso_hom_sectionsMul}: global
  sections of the right unitor
  \(G \otimes_{\mathcal{O}_X} \mathbf{1}_X \xrightarrow{\sim} G\)
  (\ref{mk-def:tensorObjRightUnitor}) carry the section product of \(a\) with the
  unit section to \(a\), via the same \(\eta\)-naturality and right-triangle
  argument and the \(\mathrm{ModuleCat}\) right-unitor formula
  \(m \otimes r \mapsto r \cdot m\).
\end{lemma}

\begin{proof}
  \leanok
  \end{proof}

\begin{lemma}[Braiding through the section multiplication]\label{mk-lem:tensorBraiding_hom_sectionsMul}
  \lean{AlgebraicGeometry.Scheme.Modules.tensorBraiding_hom_sectionsMul}
  \leanok
  \dcref{custom}

  \uses{mk-def:sectionMul, mk-def:tensorBraiding}
  Let \(\mathcal{F}, \mathcal{G}\) be sheaves of \(\mathcal{O}_X\)-modules,
  \(a \in \Gamma(X,\mathcal{F})\) and \(b \in \Gamma(X,\mathcal{G})\). Then
  \[
    \Gamma\bigl(\mathrm{tensorBraiding}_{\mathcal{F},\mathcal{G}}\bigr)
      \bigl(\mathrm{sectionsMul}_{\mathcal{F},\mathcal{G}}(a \otimes b)\bigr)
      = \mathrm{sectionsMul}_{\mathcal{G},\mathcal{F}}(b \otimes a),
  \]
  i.e.\ global sections of the braiding (\ref{mk-def:tensorBraiding}) push through
  the section multiplication (\ref{mk-def:sectionMul}), swapping the two section
  factors. This is the section-level partner of the commutativity constraint
  \ref{mk-lem:tensorPowAdd_comm}.
\end{lemma}

\begin{proof}
  \leanok
  presheaf braiding at the top open. The
  \(\mathrm{MonoidalPresheaf}\) braiding at the top open is the
  \(\mathrm{ModuleCat}\) symmetric braiding (the tensor-swap
  \(a \otimes b \mapsto b \otimes a\)), so it sends \(a \otimes b\) to
  \(b \otimes a\), giving the stated equality.
\end{proof}

The associator push-through \ref{mk-lem:tensorObjAssoc_hom_sectionsMul} is \emph{not}
a one-step \(\eta\)-naturality like the braiding
\ref{mk-lem:tensorBraiding_hom_sectionsMul}. The braiding is the sheafification of a
single presheaf morphism, so a single naturality square moves it through
\(\mathrm{sectionsMul}\). The associator \ref{mk-cor:sheafTensorObjAssoc}, by
contrast, is the canonical Mac Lane associator \(\alpha\) of \(X.\Modules\)
conjugated by five tensor-agreement bridges \(\mathrm{tensorObjIso}\)
(\ref{mk-def:tensorObjIso}) acting through the localization comparison \(\mu\)
(\ref{mk-lem:monoidalLocalization_mathlib}) --- not the image of one presheaf map ---
so the one-step argument does not transfer. We therefore cut the proof into a chain
mirroring the bridge-telescoping construction of
\ref{mk-lem:tensorPowAdd_zero_right}: two whiskered-unit legs
(\ref{mk-lem:sectionsMul_whiskerRight_unit},
\ref{mk-lem:sectionsMul_whiskerLeft_unit}) identifying each iterated section product
with a sheafification unit on the raw presheaf tensor, the presheaf associator
computed at the top open (\ref{mk-lem:presheafAssociator_top_apply}), and a
presheaf-morphism factorization (\ref{mk-lem:tensorObjAssoc_eta_factor}) assembling
them; the target then evaluates that factorization on the iterated elementary
tensor. Throughout, \(\eta\) denotes the sheafification unit
(\ref{mk-lem:presheafModule_sheafification_mathlib}) and \(\otimes_{p,\mathcal{O}_X}\)
the presheaf monoidal product (\ref{mk-lem:presheafModule_monoidal_mathlib}).

\begin{lemma}[Right-whiskered-unit leg of the iterated section product]\label{mk-lem:sectionsMul_whiskerRight_unit}
  \lean{AlgebraicGeometry.Scheme.Modules.sectionsMul_whiskerRight_unit}
  \leanok
  \dcref{custom}

  \uses{mk-def:sectionMul, mk-lem:presheafModule_monoidal_mathlib,
    mk-lem:isIso_sheafification_whiskerRight_unit,
    mk-lem:presheafModule_sheafification_mathlib}
  Let \(A, B, C\) be sheaves of \(\mathcal{O}_X\)-modules. As morphisms of
  presheaves of modules
  \[
    (A \otimes_{p,\mathcal{O}_X} B) \otimes_{p,\mathcal{O}_X} C
      \;\longrightarrow\;
    (A \otimes_{\mathcal{O}_X} B) \otimes_{\mathcal{O}_X} C,
  \]
  the right-whiskered-unit leg of the iterated section product equals the
  sheafification unit after the right-whiskered unit:
  \[
    \eta_{(A \otimes_{\mathcal{O}_X} B) \otimes_{p,\mathcal{O}_X} C}
      \circ (\eta_{A \otimes_{p,\mathcal{O}_X} B} \rhd C),
  \]
  where the right-whiskered unit \(\eta_{A \otimes_{p,\mathcal{O}_X} B} \rhd C\) is
  the comparison whose sheafification is invertible by
  \ref{mk-lem:isIso_sheafification_whiskerRight_unit}. Its component at the top open is a \(\mathrm{ModuleCat}\)
  morphism; equivalently, evaluating that component on the elementary tensor
  \((a \otimes b) \otimes c\) (\(a \in \Gamma(X,A)\), \(b \in \Gamma(X,B)\),
  \(c \in \Gamma(X,C)\)) recovers the iterated section product over the
  already-sheafified first factor:
  \[
    \mathrm{sectionsMul}_{A \otimes_{\mathcal{O}_X} B,\, C}\bigl(
      \mathrm{sectionsMul}_{A,B}(a \otimes b) \otimes c\bigr)
    \;=\;
    \Gamma\bigl(
      \eta_{(A \otimes_{\mathcal{O}_X} B) \otimes_{p,\mathcal{O}_X} C}
        \circ (\eta_{A \otimes_{p,\mathcal{O}_X} B} \rhd C)
    \bigr)\bigl((a \otimes b) \otimes c\bigr).
  \]
\end{lemma}

\begin{proof}
  \leanok
  By the objectwise formula of \ref{mk-lem:presheafModule_monoidal_mathlib} the
  right-whiskering \(\eta_{A \otimes_{p,\mathcal{O}_X} B} \rhd C\) sends, at the top
  open, \((a \otimes b) \otimes c\) to
  \(\bigl(\eta_{A \otimes_{p,\mathcal{O}_X} B}(a \otimes b)\bigr) \otimes c\);
  composing the two units gives the stated equality. The whisker comparison
  \ref{mk-lem:isIso_sheafification_whiskerRight_unit} records that this leg is the
  invertible right-whiskered unit, the form in which it enters the bridge
  telescoping of \ref{mk-lem:tensorObjAssoc_eta_factor}.
\end{proof}

\begin{lemma}[Left-whiskered-unit leg of the iterated section product]\label{mk-lem:sectionsMul_whiskerLeft_unit}
  \lean{AlgebraicGeometry.Scheme.Modules.sectionsMul_whiskerLeft_unit}
  \leanok
  \dcref{custom}

  \uses{mk-def:sectionMul, mk-lem:presheafModule_monoidal_mathlib,
    mk-lem:isIso_sheafification_whiskerRight_unit,
    mk-lem:presheafModule_sheafification_mathlib}
  Let \(A, B, C\) be sheaves of \(\mathcal{O}_X\)-modules. As morphisms of
  presheaves of modules
  \[
    A \otimes_{p,\mathcal{O}_X} (B \otimes_{p,\mathcal{O}_X} C)
      \;\longrightarrow\;
    A \otimes_{\mathcal{O}_X} (B \otimes_{\mathcal{O}_X} C),
  \]
  the left-whiskered-unit leg of the iterated section product equals the
  sheafification unit after the left-whiskered unit:
  \[
    \eta_{A \otimes_{p,\mathcal{O}_X} (B \otimes_{\mathcal{O}_X} C)}
      \circ (A \lhd \eta_{B \otimes_{p,\mathcal{O}_X} C}),
  \]
  the left-whiskered analogue of \ref{mk-lem:sectionsMul_whiskerRight_unit}, with the
  left-whiskered unit \(A \lhd \eta_{B \otimes_{p,\mathcal{O}_X} C}\) (invertible
  after sheafification by \ref{mk-lem:isIso_sheafification_whiskerRight_unit} via the
  presheaf braiding) replacing the right-whiskered one. Its component at the
  top open is a \(\mathrm{ModuleCat}\) morphism; equivalently, evaluating that
  component on the elementary tensor \(a \otimes (b \otimes c)\)
  (\(a \in \Gamma(X,A)\), \(b \in \Gamma(X,B)\), \(c \in \Gamma(X,C)\)) recovers the
  iterated section product over the already-sheafified second factor:
  \[
    \mathrm{sectionsMul}_{A,\, B \otimes_{\mathcal{O}_X} C}\bigl(
      a \otimes \mathrm{sectionsMul}_{B,C}(b \otimes c)\bigr)
    \;=\;
    \Gamma\bigl(
      \eta_{A \otimes_{p,\mathcal{O}_X} (B \otimes_{\mathcal{O}_X} C)}
        \circ (A \lhd \eta_{B \otimes_{p,\mathcal{O}_X} C})
    \bigr)\bigl(a \otimes (b \otimes c)\bigr).
  \]
\end{lemma}

\begin{proof}
  \leanok
  \(A \lhd \eta_{B \otimes_{p,\mathcal{O}_X} C}\) at the top open as
  \(a \otimes (b \otimes c) \mapsto a \otimes \eta_{B \otimes_{p,\mathcal{O}_X} C}(b \otimes c)\);
  composing with the outer unit gives the equality. The left-whiskered form of
  \ref{mk-lem:isIso_sheafification_whiskerRight_unit} --- its proof handles left
  whiskering by conjugating with the presheaf braiding --- supplies the
  invertibility used in \ref{mk-lem:tensorObjAssoc_eta_factor}.
\end{proof}

\begin{lemma}[Right-whisker naturality of the section product]\label{mk-lem:sectionsMul_whiskerRight_natural}
  \lean{AlgebraicGeometry.Scheme.Modules.sectionsMul_whiskerRight_natural}
  \leanok
  \dcref{custom}

  \uses{mk-def:sectionMul, mk-lem:presheafModule_monoidal_mathlib,
    mk-lem:presheafModule_sheafification_mathlib}
  Let \(f : F \to F'\) be a morphism of sheaves of \(\mathcal{O}_X\)-modules and let
  \(G\) be a sheaf of modules. The section product \(\mathrm{sectionsMul}\) is
  \emph{natural in its first argument}: evaluating at the top open on
  \(x \in \Gamma(X,F)\), \(y \in \Gamma(X,G)\),
  \[
    \mathrm{sectionsMul}_{F',G}\bigl(\Gamma(f)(x) \otimes y\bigr)
      \;=\;
    \Gamma\bigl(f \rhd G\bigr)\bigl(\mathrm{sectionsMul}_{F,G}(x \otimes y)\bigr),
  \]
  where \(f \rhd G : F \otimes_{\mathcal{O}_X} G \to F' \otimes_{\mathcal{O}_X} G\) is
  the right-whiskering of \(f\) by \(G\) and \(\Gamma(-)\) denotes the top-open
  component of a sheaf morphism. This is the general-morphism analogue of
  \ref{mk-lem:sectionsMul_whiskerRight_unit} (which is the special case
  \(f = \eta_{F \otimes_p \mathbf{1}}\), the sheafification unit), and is the slide
  used in the associativity leg of \ref{mk-lem:sectionMul_coherent} to move an inner
  comparison \(\Gamma(\mu)\) out of the first tensor factor of an outer
  \(\mathrm{sectionsMul}\).
\end{lemma}

\begin{proof}
  \leanok
  By \ref{mk-def:sectionMul}, \(\mathrm{sectionsMul}_{F,G}\) is the top-open component of
  the sheafification unit \(\eta_{F \otimes_{p} G}\), itself the component of a natural
  transformation (\ref{mk-lem:presheafModule_sheafification_mathlib}). Naturality of
  \(\eta\) along the presheaf-of-modules right-whiskering \(f \rhd_p G\) reads, as
  presheaf morphisms,
  \[
    \eta_{F' \otimes_p G} \circ (f \rhd_p G)
      \;=\;
    (f \rhd G) \circ \eta_{F \otimes_p G}.
  \]
  The objectwise whiskering formula of \ref{mk-lem:presheafModule_monoidal_mathlib}
  identifies \(f \rhd_p G\) at the top open as
  \(x \otimes y \mapsto \Gamma(f)(x) \otimes y\). Evaluating the naturality square at
  the top open on the elementary tensor \(x \otimes y\) gives the displayed equality.
\end{proof}

\begin{lemma}[Left-whisker naturality of the section product]\label{mk-lem:sectionsMul_whiskerLeft_natural}
  \lean{AlgebraicGeometry.Scheme.Modules.sectionsMul_whiskerLeft_natural}
  \leanok
  \dcref{custom}

  \uses{mk-def:sectionMul, mk-lem:presheafModule_monoidal_mathlib,
    mk-lem:presheafModule_sheafification_mathlib}
  Let \(F\) be a sheaf of \(\mathcal{O}_X\)-modules and let \(g : G \to G'\) be a
  morphism of sheaves of modules. The section product \(\mathrm{sectionsMul}\) is
  \emph{natural in its second argument}: evaluating at the top open on
  \(x \in \Gamma(X,F)\), \(y \in \Gamma(X,G)\),
  \[
    \mathrm{sectionsMul}_{F,G'}\bigl(x \otimes \Gamma(g)(y)\bigr)
      \;=\;
    \Gamma\bigl(F \lhd g\bigr)\bigl(\mathrm{sectionsMul}_{F,G}(x \otimes y)\bigr),
  \]
  where \(F \lhd g : F \otimes_{\mathcal{O}_X} G \to F \otimes_{\mathcal{O}_X} G'\) is
  the left-whiskering of \(g\) by \(F\). This is the left-handed analogue of
  \ref{mk-lem:sectionsMul_whiskerRight_natural} (the general-morphism analogue of
  \ref{mk-lem:sectionsMul_whiskerLeft_unit}), and is the slide used in the associativity
  leg of \ref{mk-lem:sectionMul_coherent} to move an inner comparison \(\Gamma(\mu)\) out
  of the second tensor factor of an outer \(\mathrm{sectionsMul}\).
\end{lemma}

\begin{proof}
  \leanok
  Identical to \ref{mk-lem:sectionsMul_whiskerRight_natural} with the roles of the two
  tensor factors exchanged: naturality of the sheafification unit \(\eta\)
  (\ref{mk-lem:presheafModule_sheafification_mathlib}) along the left-whiskering
  \(F \lhd_p g\), whose top-open component is
  \(x \otimes y \mapsto x \otimes \Gamma(g)(y)\) by the objectwise formula of
  \ref{mk-lem:presheafModule_monoidal_mathlib}, evaluated on \(x \otimes y\).
\end{proof}

\begin{lemma}[Presheaf associator at the top open]\label{mk-lem:presheafAssociator_top_apply}
  \lean{AlgebraicGeometry.Scheme.Modules.presheafAssociator_top_apply}
  \leanok
  \dcref{custom}

  \uses{mk-lem:presheafModule_monoidal_mathlib}
  Let \(A, B, C\) be presheaves of \(\mathcal{O}_X\)-modules. The associator
  \(\alpha^{p}_{A,B,C}\) of the presheaf monoidal structure
  (\ref{mk-lem:presheafModule_monoidal_mathlib}) evaluated at the top open is the
  \(\mathrm{ModuleCat}\) associator on the module sections, sending the iterated
  elementary tensor to its reassociation:
  \[
    \alpha^{p}_{A,B,C}\bigm|_{\top}\bigl((a \otimes b) \otimes c\bigr)
      = a \otimes (b \otimes c),
    \qquad a \in A(\top),\ b \in B(\top),\ c \in C(\top).
  \]
\end{lemma}

\begin{proof}
  \leanok\end{proof}

\begin{lemma}[Presheaf-morphism factorization of the associator]\label{mk-lem:tensorObjAssoc_eta_factor}
  \lean{AlgebraicGeometry.Scheme.Modules.tensorObjAssoc_eta_factor}
  \leanok
  \dcref{custom}

  \uses{mk-cor:sheafTensorObjAssoc, mk-def:tensorObjIso,
    mk-lem:isIso_sheafification_whiskerRight_unit, mk-lem:monoidalLocalization_mathlib,
    mk-lem:presheafModule_sheafification_mathlib, mk-lem:presheafModule_monoidal_mathlib,
    mk-lem:tensorObjAssoc_eta_factor_sheaf}
  Let \(A, B, C\) be sheaves of \(\mathcal{O}_X\)-modules. Then, as morphisms of
  presheaves of modules
  \((A \otimes_{p,\mathcal{O}_X} B) \otimes_{p,\mathcal{O}_X} C
    \;\longrightarrow\; A \otimes_{\mathcal{O}_X} (B \otimes_{\mathcal{O}_X} C)\),
  \[
    \mathrm{tensorObjAssoc}_{A,B,C}
      \circ \eta_{(A \otimes_{\mathcal{O}_X} B) \otimes_{p,\mathcal{O}_X} C}
      \circ (\eta_{A \otimes_{p,\mathcal{O}_X} B} \rhd C)
    \;=\;
    \eta_{A \otimes_{p,\mathcal{O}_X} (B \otimes_{\mathcal{O}_X} C)}
      \circ (A \lhd \eta_{B \otimes_{p,\mathcal{O}_X} C})
      \circ \alpha^{p}_{A,B,C},
  \]
  where \(\alpha^{p}\) is the presheaf associator
  (\ref{mk-lem:presheafModule_monoidal_mathlib}). This is the
  \(\eta\)-naturality-plus-bridge-telescoping identity that makes
  \(\Gamma(\mathrm{tensorObjAssoc})\) push the presheaf associator through the
  iterated \(\mathrm{sectionsMul}\). It is stated as an equality of presheaf
  morphisms (holding component-wise at every open), the form in which the
  bridge-telescoping is carried out; evaluating both sides at the top open on the
  iterated elementary tensor is what feeds the associator push-through
  \ref{mk-lem:tensorObjAssoc_hom_sectionsMul}.
\end{lemma}

\begin{proof}
  \leanok
  obtain it by transporting an identity of sheaf morphisms across the naturality of
  the sheafification unit \(\eta\) (\ref{mk-lem:presheafModule_sheafification_mathlib}).
  Each of the three constituent presheaf maps --- the right-whiskered unit
  \(\eta_{A \otimes_{p} B} \rhd C\), the presheaf associator \(\alpha^{p}_{A,B,C}\)
  (\ref{mk-lem:presheafModule_monoidal_mathlib}), and the left-whiskered unit
  \(A \lhd \eta_{B \otimes_{p} C}\) --- fits into a naturality square for \(\eta\):
  composing it with the unit at its codomain equals the unit at its domain composed
  with \(T(L(\cdot))\), the underlying presheaf map of its sheafification. Applying
  these three naturality squares rewrites both sides of the claim as
  \(\eta_{(A \otimes_{p} B) \otimes_{p} C}\) post-composed with \(T\) of a single
  sheaf morphism, so the claim reduces to the equality of those two sheaf morphisms
  inside \(X.\Modules\),
  \[
    L(\eta_{A \otimes_{p} B} \rhd C) \;\circ\; \mathrm{tensorObjAssoc}_{A,B,C}
      \;=\; L(\alpha^{p}_{A,B,C}) \;\circ\; L(A \lhd \eta_{B \otimes_{p} C}),
  \]
  which is the sheaf-level core \ref{mk-lem:tensorObjAssoc_eta_factor_sheaf}. (The
  shared middle object is presented two definitionally-equal but not syntactically
  equal ways --- \(\eta\)'s codomain carries the right-adjoint's
  \(\mathrm{restrictScalars}(\mathbf{1})\) decoration absent from
  \(\mathrm{tensorObjAssoc}\)'s domain --- so the three naturality rewrites are
  performed up to this definitional equality.)
\end{proof}

\begin{lemma}[Sheafification unit on the left triangle]\label{mk-lem:sheafification_map_unit_eq}
  \lean{AlgebraicGeometry.Scheme.Modules.sheafification_map_unit_eq}
  \leanok
  \dcref{custom}

  \uses{mk-lem:presheafModule_sheafification_mathlib}
  Let \(L\) be the sheafification functor on presheaves of \(\mathcal{O}_X\)-modules,
  with unit \(\eta\) and counit \(\varepsilon\) of the sheafification adjunction
  (\ref{mk-lem:presheafModule_sheafification_mathlib}). For a presheaf of modules
  \(P\),
  \[
    L(\eta_P) \;=\; \varepsilon_{L P}^{-1},
  \]
  i.e.\ the sheafification of the unit at \(P\) is the inverse of the counit
  isomorphism at \(LP\).
\end{lemma}

\begin{proof}
  \leanok\end{proof}

\begin{lemma}[Sheaf-level associator factorization (core of \ref{mk-lem:tensorObjAssoc_eta_factor})]\label{mk-lem:tensorObjAssoc_eta_factor_sheaf}
  \lean{AlgebraicGeometry.Scheme.Modules.tensorObjAssoc_eta_factor_sheaf}
  \leanok
  \dcref{custom}

  \uses{mk-cor:sheafTensorObjAssoc, mk-def:tensorObjIso,
    mk-lem:isIso_sheafification_whiskerRight_unit, mk-lem:monoidalLocalization_mathlib,
    mk-lem:sheafification_map_unit_eq, mk-lem:tensorPowAdd_zero_right}
  Let \(A, B, C\) be sheaves of \(\mathcal{O}_X\)-modules and \(L\) the sheafification
  functor. As morphisms of sheaves of \(\mathcal{O}_X\)-modules
  \[
    L\bigl(\eta_{A \otimes_{p,\mathcal{O}_X} B} \rhd C\bigr)
      \;\circ\; \mathrm{tensorObjAssoc}_{A,B,C}
    \;=\;
    L\bigl(\alpha^{p}_{A,B,C}\bigr)
      \;\circ\; L\bigl(A \lhd \eta_{B \otimes_{p,\mathcal{O}_X} C}\bigr),
  \]
  where \(\eta\) is the sheafification unit and \(\alpha^{p}\) the presheaf
  associator (\ref{mk-lem:presheafModule_monoidal_mathlib}). Equivalently,
  \(\mathrm{tensorObjAssoc}\) (the canonical-\(\alpha\) form
  \ref{mk-cor:sheafTensorObjAssoc}) coincides with the original \(\eta\)-conjugated
  associator \(\bigl(\mathrm{asIso}\,L(\eta \rhd C)\bigr)^{-1}
  \circ L(\alpha^{p}) \circ \mathrm{asIso}\,L(A \lhd \eta)\).
\end{lemma}
\begin{proof}
  associator \(\alpha\) of \(X.\Modules\) flanked by the four tensor-agreement
  bridges \(\Phi = \mathrm{tensorObjIso}\) (\ref{mk-def:tensorObjIso}) whiskered on the
  outer and inner factors,
  \[
    \mathrm{tensorObjAssoc} = \Phi^{-1}_{A\otimes B,\,C}
      \circ (\Phi^{-1}_{A,B} \rhd C) \circ \alpha_{A,B,C}
      \circ (A \lhd \Phi_{B,C}) \circ \Phi_{A,\,B\otimes C}.
  \]
  Each bridge \(\Phi\) is built from the localization comparison \(\mu\)
  (\ref{mk-lem:monoidalLocalization_mathlib}); the canonical associator \(\alpha\)
  transported through the \(\mu\)'s is the sheafification \(L(\alpha^{p})\) of the
  presheaf associator. Substitute \(L(\eta_P) = \varepsilon_{LP}^{-1}\)
  (\ref{mk-lem:sheafification_map_unit_eq}) so each whiskered sheafification unit on the
  two sides becomes a whiskered counit, and rewrite the three localization-comparison
  segments via the strong-monoidality whisker comparison
  \ref{mk-lem:isIso_sheafification_whiskerRight_unit} (in its right- and, via the
  presheaf braiding, left-whiskered forms) together with the naturality of the
  localization comparison \(\mu\) (\ref{mk-lem:monoidalLocalization_mathlib}) in each of
  its two arguments and the comparison identifying the transported canonical associator
  with \(L(\alpha^{p})\). The \(\mathrm{tensorObjIso}\) bridges
  then telescope in adjacent pairs --- the bridge-cancellation recipe already carried
  out for \ref{mk-lem:tensorPowAdd_zero_right} --- leaving exactly
  \(L(\eta \rhd C)\) on the left and \(L(\alpha^{p}) \circ L(A \lhd \eta)\) on the
  right, which is the stated equality. (As in \ref{mk-lem:tensorObjAssoc_eta_factor},
  every \(\mu\)-conjugation step bridges the
  \(\mathrm{restrictScalars}(\mathbf{1})\) definitional decoration on \(\eta\)'s
  codomain against \(\mathrm{tensorObjIso}\)'s undecorated factor, so the rewrites are
  carried out up to that definitional equality.)
\end{proof}

\begin{lemma}[Associator through the section multiplication]\label{mk-lem:tensorObjAssoc_hom_sectionsMul}
  \lean{AlgebraicGeometry.Scheme.Modules.tensorObjAssoc_hom_sectionsMul}
  \leanok
  \dcref{custom}

  \uses{mk-def:sectionMul, mk-cor:sheafTensorObjAssoc,
    mk-lem:sectionsMul_whiskerRight_unit, mk-lem:sectionsMul_whiskerLeft_unit,
    mk-lem:presheafAssociator_top_apply, mk-lem:tensorObjAssoc_eta_factor}
  Let \(\mathcal{A}, \mathcal{B}, \mathcal{C}\) be sheaves of
  \(\mathcal{O}_X\)-modules and \(a \in \Gamma(X,\mathcal{A})\),
  \(b \in \Gamma(X,\mathcal{B})\), \(c \in \Gamma(X,\mathcal{C})\). Applying global
  sections of the associator (\ref{mk-cor:sheafTensorObjAssoc}) to the iterated
  section product reassociates the three factors:
  \[
    \Gamma\bigl(\mathrm{tensorObjAssoc}_{\mathcal{A},\mathcal{B},\mathcal{C}}\bigr)
      \Bigl(\mathrm{sectionsMul}\bigl(
        \mathrm{sectionsMul}(a \otimes b) \otimes c\bigr)\Bigr)
      = \mathrm{sectionsMul}\bigl(a \otimes \mathrm{sectionsMul}(b \otimes c)\bigr),
  \]
  i.e.\ \(\Gamma(\alpha)\) sends \((a \otimes b) \otimes c\) to
  \(a \otimes (b \otimes c)\) at the level of section multiplications. This is the
  \(\eta\)-naturality of the lax structure pushing the associator through
  \(\mathrm{sectionsMul}\), the section-level partner of the associativity
  constraint \ref{mk-lem:tensorPowAdd_assoc}.
\end{lemma}

\begin{proof}
  \leanok
  \ref{mk-lem:sectionsMul_whiskerRight_unit}; its right side is the RHS of the claim
  by \ref{mk-lem:presheafAssociator_top_apply} (the presheaf associator sends
  \((a \otimes b) \otimes c \mapsto a \otimes (b \otimes c)\)) followed by
  \ref{mk-lem:sectionsMul_whiskerLeft_unit}.
\end{proof}

The graded-monoid axioms below are not raw equalities between sections of two
different tensor powers --- on \(\mathbb{N}\) the indices \((i+j)+k\) and
\(i+(j+k)\), or \(0+n\) and \(n\), are equal but not definitionally so, so the two
sides a priori inhabit distinct section modules. Following
\texttt{Mathlib.LinearAlgebra.TensorPower.Basic}, we phrase each coherence as an
honest equality in a single section module, obtained by transporting one side
along an index-equality isomorphism, and provide a bridge that repackages such a
transported equality as an equality of dependent pairs.

\begin{definition}[Section component in degree \(m\)]\label{mk-def:sectionDeg}
  \lean{AlgebraicGeometry.Scheme.Modules.sectionDeg}
  \leanok
  \dcref{custom}

  \uses{mk-def:sheafTensorPow}
  Let \(\mathcal{L}\) be a sheaf of \(\mathcal{O}_X\)-modules. The
  \emph{degree-\(m\) section component} is the \(\Gamma(X,\mathcal{O}_X)\)-module of
  global sections of the \(m\)th tensor power,
  \[
    (\mathrm{sectionDeg}\ \mathcal{L})_m
      \;:=\; \Gamma(X, \mathcal{L}^{\otimes m})
      \;=\; \mathcal{L}^{\otimes m}(X)
        \qquad(\ref{mk-def:sheafTensorPow}).
  \]
  It inherits its additive-group and \(\Gamma(X,\mathcal{O}_X)\)-module structure
  from the underlying module object \(\mathcal{L}^{\otimes m}(X)\). These are the
  carriers \(m \mapsto (\mathrm{sectionDeg}\ \mathcal{L})_m\) of the section graded
  ring \(R(X_s, L_s)\) (\ref{mk-lem:sectionGradedRing_gcommSemiring}) and the
  domains of the index transport \ref{mk-def:sectionsCast}.
\end{definition}

\begin{definition}[Index-equality transport of section components]\label{mk-def:sectionsCast}
  \lean{AlgebraicGeometry.Scheme.Modules.sectionsCast}
  \leanok
  \dcref{custom}

  \uses{mk-def:sheafTensorPow}
  Let \(\mathcal{L}\) be an invertible sheaf on \(X\) and let \(h : i = j\) be an
  equality of natural numbers. Applying the global-sections functor
  \(\Gamma(X,-)\) to the canonical isomorphism
  \(\mathcal{L}^{\otimes i} \xrightarrow{\sim} \mathcal{L}^{\otimes j}\) induced by
  \(h\) under the tensor-power functor (\ref{mk-def:sheafTensorPow}) yields the
  \(\Gamma(X,\mathcal{O}_X)\)-linear isomorphism
  \[
    \mathrm{sectionsCast}(h) :
    \Gamma(X, \mathcal{L}^{\otimes i}) \;\xrightarrow{\ \sim\ }\;
    \Gamma(X, \mathcal{L}^{\otimes j}),
  \]
  the \emph{section-component transport} along \(h\). It is the section-level
  analogue of the tensor-power index cast \(\mathrm{TensorPower.cast}\) of
  \texttt{Mathlib.LinearAlgebra.TensorPower.Basic}, which transports
  \(\bigotimes^{i} M\) to \(\bigotimes^{j} M\) along \(i = j\). Because the
  isomorphism induced by the reflexive equality is the identity and
  \(\Gamma(X,-)\) preserves identities, the transport along a reflexive index
  equality is the identity map (the reflexivity simplification
  \(\mathrm{sectionsCast\_refl}\), \ref{mk-lem:sectionsCast_refl}).
\end{definition}

\begin{lemma}[Reflexivity of the section transport]\label{mk-lem:sectionsCast_refl}
  \lean{AlgebraicGeometry.Scheme.Modules.sectionsCast_refl}
  \leanok
  \dcref{custom}

  \uses{mk-def:sectionsCast}
  For any \(i\), the transport along the reflexive equality \(i = i\)
  (\ref{mk-def:sectionsCast}) equals the identity automorphism of
  \(\Gamma(X, \mathcal{L}^{\otimes i})\). This is the section-level counterpart of
  \(\mathrm{TensorPower.cast\_refl}\), and serves as a simplification lemma
  collapsing the transport along a reflexive index equality.
\end{lemma}

\begin{proof}
  \leanok\end{proof}

\begin{lemma}[Cast-mediated equality in the graded monoid]\label{mk-lem:gradedMonoid_eq_of_cast}
  \lean{AlgebraicGeometry.Scheme.Modules.gradedMonoid_eq_of_cast}
  \leanok
  \dcref{custom}

  \uses{mk-def:sectionsCast, mk-lem:sectionsCast_refl}
  Write the carrier of the graded monoid as the dependent sum
  \(\coprod_{n} \Gamma(X, \mathcal{L}^{\otimes n})\), whose elements are pairs
  \((n, s)\) with \(s \in \Gamma(X, \mathcal{L}^{\otimes n})\). Let
  \(a = (i, x)\) and \(b = (j, y)\) be two such pairs. If \(h : i = j\) and the
  transported component agrees, \(\mathrm{sectionsCast}(h)\,x = y\), then
  \(a = b\). That is, a component-level equality mediated by the index transport
  \ref{mk-def:sectionsCast} repackages into a genuine equality of dependent pairs.
  This is the section-level analogue of the bridge
  \(\mathrm{gradedMonoid\_eq\_of\_cast}\) of
  \texttt{Mathlib.LinearAlgebra.TensorPower.Basic} (line~123 there).
\end{lemma}

\begin{proof}
  \leanok
  \end{proof}

\begin{definition}[Graded multiplication on section components]
  \label{mk-def:sectionGradedGMul}
  \lean{AlgebraicGeometry.Scheme.Modules.instGMulNatSectionDeg}
  \dcref{custom}

  \uses{mk-def:sectionDeg, mk-def:sectionMul, mk-lem:sheafTensorPow_add}
  The section components \ref{mk-def:sectionDeg} carry a graded multiplication
  \[
    (\mathrm{sectionDeg}\ \mathcal{L})_i \times
      (\mathrm{sectionDeg}\ \mathcal{L})_j
    \;\longrightarrow\; (\mathrm{sectionDeg}\ \mathcal{L})_{i+j},
    \qquad (a, b) \longmapsto a \cdot b,
  \]
  landing in the index-sum component (a \(\mathrm{GradedMonoid.GMul}\) on
  \(m \mapsto (\mathrm{sectionDeg}\ \mathcal{L})_m\)). It is the section
  multiplication \(a \otimes b \mapsto a \cdot b\) (\ref{mk-def:sectionMul}) on
  \(\mathcal{L}^{\otimes i}\) and \(\mathcal{L}^{\otimes j}\), followed by global
  sections of the tensor-power comparison isomorphism \(\mu_{i,j}\)
  (\ref{mk-lem:sheafTensorPow_add}),
  \[
    \Gamma(X, \mathcal{L}^{\otimes i}) \otimes_{\Gamma(X,\mathcal{O}_X)}
      \Gamma(X, \mathcal{L}^{\otimes j})
    \xrightarrow{\ \cdot\ }
    \Gamma(X, \mathcal{L}^{\otimes i} \otimes_{\mathcal{O}_X}
      \mathcal{L}^{\otimes j})
    \xrightarrow{\ \Gamma(\mu_{i,j})\ }
    \Gamma(X, \mathcal{L}^{\otimes (i+j)}).
  \]
\end{definition}

\begin{definition}[Graded unit in degree zero]
  \label{mk-def:sectionGradedGOne}
  \lean{AlgebraicGeometry.Scheme.Modules.instGOneNatSectionDeg}
  \dcref{custom}

  \uses{mk-def:sectionDeg, mk-lem:moduleUnit_mathlib}
  The graded unit (a \(\mathrm{GradedMonoid.GOne}\) on
  \(m \mapsto (\mathrm{sectionDeg}\ \mathcal{L})_m\)) is the element
  \[
    1 \;\in\; (\mathrm{sectionDeg}\ \mathcal{L})_0
      = \Gamma(X, \mathcal{L}^{\otimes 0}) = \Gamma(X, \mathcal{O}_X),
  \]
  the image of \(1 \in \Gamma(X, \mathcal{O}_X)\) under the canonical
  \(\Gamma(X,\mathcal{O}_X)\)-module identification
  \(\Gamma(X, \mathcal{O}_X) = \Gamma(X, \mathcal{L}^{\otimes 0})\) coming from the
  unit module \(\mathcal{L}^{\otimes 0} = \mathbf{1}_X\)
  (\ref{mk-lem:moduleUnit_mathlib}).
\end{definition}

\begin{lemma}[Coherence of the section multiplication]
  \label{mk-lem:sectionMul_coherent}
  \lean{AlgebraicGeometry.Scheme.Modules.sectionsMul_one_mul,
    AlgebraicGeometry.Scheme.Modules.sectionsMul_mul_one,
    AlgebraicGeometry.Scheme.Modules.sectionsMul_mul_assoc,
    AlgebraicGeometry.Scheme.Modules.sectionsMul_mul_comm}
  \dcref{custom}

  \uses{mk-def:sectionMul, mk-def:sectionsCast, mk-def:sectionGradedGMul,
    mk-def:sectionGradedGOne, mk-def:isInvertible, mk-lem:moduleUnit_mathlib,
    mk-lem:sheafTensorPow_add,
    mk-lem:tensorPowAdd_zero_right, mk-lem:tensorPowAdd_zero_left, mk-lem:tensorPowAdd_comm,
    mk-lem:tensorPowAdd_assoc,
    mk-lem:tensorObjUnitIso_hom_sectionsMul, mk-lem:tensorObjRightUnitor_hom_sectionsMul,
    mk-lem:tensorBraiding_hom_sectionsMul, mk-lem:tensorObjAssoc_hom_sectionsMul,
    mk-lem:sectionsMul_whiskerRight_natural, mk-lem:sectionsMul_whiskerLeft_natural}
  Write \(a \cdot b\) for the degreewise section multiplication on tensor powers:
  for \(a \in \Gamma(X, \mathcal{L}^{\otimes i})\) and
  \(b \in \Gamma(X, \mathcal{L}^{\otimes j})\) it is the elementary section
  multiplication (\ref{mk-def:sectionMul}) followed by the tensor-power comparison
  isomorphism \(\mu_{i,j}\) (\ref{mk-lem:sheafTensorPow_add}), landing in
  \(\Gamma(X, \mathcal{L}^{\otimes (i+j)})\); and let \(1\) denote the unit of
  degree \(0\), the image of \(1 \in \Gamma(X,\mathcal{O}_X)\) under the
  identification \(\Gamma(X,\mathcal{O}_X) = \Gamma(X, \mathcal{L}^{\otimes 0})\)
  (\ref{mk-lem:moduleUnit_mathlib}). Mirroring
  \(\mathrm{TensorPower.\{one\_mul, mul\_one, mul\_assoc\}}\) together with the
  commutativity lemma of \texttt{Mathlib.LinearAlgebra.TensorPower.Basic}, the
  multiplication satisfies the following four component-level identities. Each is
  an honest equality in a single section module, obtained by transporting one side
  along the relevant index equality via \ref{mk-def:sectionsCast}; the subscript on
  \(\mathrm{sectionsCast}\) records that index equality. The first three ---
  left and right unitality and associativity --- hold for an \emph{arbitrary}
  sheaf of modules \(\mathcal{L}\), and make the section components a
  non-commutative graded semiring. The fourth --- commutativity --- requires
  \(\mathcal{L}\) to be \emph{invertible} (\ref{mk-def:isInvertible}); it is the
  upgrade to a graded \emph{commutative} semiring and is false for general
  \(\mathcal{L}\), the section ring then being the free tensor algebra on
  \(\Gamma(X,\mathcal{L})\).
  \begin{itemize}
    \item \emph{Left unitality.} For \(a \in \Gamma(X, \mathcal{L}^{\otimes n})\),
      \[
        \mathrm{sectionsCast}_{\,0+n=n}\,(1 \cdot a) = a .
      \]
    \item \emph{Right unitality.} For \(a \in \Gamma(X, \mathcal{L}^{\otimes n})\),
      \[
        \mathrm{sectionsCast}_{\,n+0=n}\,(a \cdot 1) = a .
      \]
    \item \emph{Associativity.} For
      \(a \in \Gamma(X, \mathcal{L}^{\otimes n_a})\),
      \(b \in \Gamma(X, \mathcal{L}^{\otimes n_b})\),
      \(c \in \Gamma(X, \mathcal{L}^{\otimes n_c})\),
      \[
        \mathrm{sectionsCast}_{\,(n_a+n_b)+n_c = n_a+(n_b+n_c)}\,
          ((a \cdot b) \cdot c) = a \cdot (b \cdot c) .
      \]
    \item \emph{Commutativity (\(\mathcal{L}\) invertible,
      \ref{mk-def:isInvertible}).} For
      \(a \in \Gamma(X, \mathcal{L}^{\otimes n_a})\) and
      \(b \in \Gamma(X, \mathcal{L}^{\otimes n_b})\),
      \[
        \mathrm{sectionsCast}_{\,n_a+n_b = n_b+n_a}\,(a \cdot b) = b \cdot a .
      \]
  \end{itemize}
  The transport is indispensable: the indices on the two sides of each identity
  are equal but not definitionally so, hence a priori live in distinct section
  modules, and \ref{mk-def:sectionsCast} moves one side into the other's module to
  produce a genuine equality there.
\end{lemma}

\begin{proof}  \uses{mk-def:sectionMul, mk-def:sectionsCast, mk-def:sectionGradedGMul,
    mk-lem:sheafTensorPow_add, mk-lem:moduleUnit_mathlib,
    mk-lem:tensorPowAdd_zero_right, mk-lem:tensorPowAdd_zero_left, mk-lem:tensorPowAdd_comm,
    mk-lem:tensorPowAdd_assoc,
    mk-lem:tensorObjUnitIso_hom_sectionsMul, mk-lem:tensorObjRightUnitor_hom_sectionsMul,
    mk-lem:tensorBraiding_hom_sectionsMul, mk-lem:tensorObjAssoc_hom_sectionsMul,
    mk-lem:sectionsMul_whiskerRight_natural, mk-lem:sectionsMul_whiskerLeft_natural}
  Each identity is obtained by unfolding the degreewise multiplication
  \(a \cdot b = \Gamma(\mu_{i,j})\bigl(\mathrm{sectionsMul}(a \otimes b)\bigr)\)
  (\ref{mk-def:sectionGradedGMul}, \ref{mk-def:sectionMul},
  \ref{mk-lem:sheafTensorPow_add}) and combining two ingredients: an
  \emph{iso-level constraint} on the comparison family \(\{\mu_{m,m'}\}\), and the
  \emph{section-level naturality} that pushes the corresponding canonical coherence
  map through \(\mathrm{sectionsMul}\). The index transport \ref{mk-def:sectionsCast}
  --- the image under \(\Gamma(X,-)\) of the reindexing isomorphism of tensor
  powers --- matches the two sides, which a priori inhabit the distinct section
  modules of two equal-but-not-definitionally-equal indices, into a common module.

  \emph{Left unitality.} Unfolding \(1 \cdot a\) and using the base case
  \(\mu_{0,n} = \lambda\) of \ref{mk-lem:sheafTensorPow_add}, the inner reindexing
  cast pairs with the outer \(\mathrm{sectionsCast}_{\,0+n=n}\) and the two cancel;
  \ref{mk-lem:tensorObjUnitIso_hom_sectionsMul} then closes the leg.

  \emph{Right unitality.} Unfolding \(a \cdot 1\) and rewriting the comparison by
  the right-unit constraint \(\mu_{n,0} = \mathrm{tensorObjRightUnitor}\)
  (\ref{mk-lem:tensorPowAdd_zero_right}),
  \ref{mk-lem:tensorObjRightUnitor_hom_sectionsMul} closes the leg after the residual
  reflexive transport collapses.

  \emph{Associativity.} Unfold both bracketings to composites
  \(\Gamma(\mu) \circ \mathrm{sectionsMul}\). Unlike the unit legs --- where the
  inner factor is the literal degree-\(0\) unit and carries no comparison --- here
  each bracketing hides an \emph{inner} comparison nested inside one tensor factor of
  the outer \(\mathrm{sectionsMul}\), and these must first be slid out. On the left,
  \[
    (a \cdot b) \cdot c
      = \Gamma(\mu_{n_a+n_b,\,n_c})\Bigl(
          \mathrm{sectionsMul}_{\mathcal{L}^{n_a+n_b},\mathcal{L}^{n_c}}\bigl(
            \Gamma(\mu_{n_a,n_b})(\mathrm{sectionsMul}(a \otimes b)) \otimes c\bigr)\Bigr),
  \]
  and the right-whisker naturality \ref{mk-lem:sectionsMul_whiskerRight_natural}
  (applied to \(f = \mu_{n_a,n_b}\), \(G = \mathcal{L}^{n_c}\)) slides
  \(\Gamma(\mu_{n_a,n_b})\) out of the first factor, giving
  \(\Gamma\bigl(\mu_{n_a+n_b,\,n_c} \circ (\mu_{n_a,n_b} \rhd \mathcal{L}^{n_c})\bigr)\)
  applied to the bare iterated product
  \(\mathrm{sectionsMul}(\mathrm{sectionsMul}(a \otimes b) \otimes c)\). Symmetrically,
  on the right, the left-whisker naturality
  \ref{mk-lem:sectionsMul_whiskerLeft_natural} (applied to \(g = \mu_{n_b,n_c}\),
  \(F = \mathcal{L}^{n_a}\)) slides \(\Gamma(\mu_{n_b,n_c})\) out of the second factor,
  giving \(\Gamma\bigl(\mu_{n_a,\,n_b+n_c} \circ (\mathcal{L}^{n_a} \lhd \mu_{n_b,n_c})\bigr)\)
  applied to \(\mathrm{sectionsMul}(a \otimes \mathrm{sectionsMul}(b \otimes c))\). Both
  sides are now \(\Gamma(\text{comparison composite}) \circ (\text{iterated }
  \mathrm{sectionsMul})\). The iso-level associativity (pentagon) constraint
  \ref{mk-lem:tensorPowAdd_assoc} equates the two comparison composites along the
  associator \(\alpha\) modulo the reindexing
  \((n_a+n_b)+n_c = n_a+(n_b+n_c)\):
  \[
    \mu_{n_a+n_b,\,n_c} \circ (\mu_{n_a,n_b} \rhd \mathcal{L}^{n_c})
      \;=\;
    \mathrm{(reindex)} \circ \mu_{n_a,\,n_b+n_c}
      \circ (\mathcal{L}^{n_a} \lhd \mu_{n_b,n_c}) \circ \alpha,
  \]
  and the section-level associator naturality
  \ref{mk-lem:tensorObjAssoc_hom_sectionsMul} pushes the residual \(\Gamma(\alpha)\)
  through the iterated \(\mathrm{sectionsMul}\), reassociating the three section
  factors; the transport
  \(\mathrm{sectionsCast}_{\,(n_a+n_b)+n_c = n_a+(n_b+n_c)}\) absorbs the reindexing
  and identifies the two sides in a single module. This is the right-unit leg with the
  associator triple
  (\ref{mk-lem:tensorPowAdd_assoc} + \ref{mk-lem:tensorObjAssoc_hom_sectionsMul} +
  the two slides \ref{mk-lem:sectionsMul_whiskerRight_natural},
  \ref{mk-lem:sectionsMul_whiskerLeft_natural}) in place of the right-unitor pair.

  \emph{Commutativity.} Symmetrically, the iso-level commutativity constraint
  \ref{mk-lem:tensorPowAdd_comm} factors \(\mu_{n_a,n_b}\) through the braiding into
  \(\mu_{n_b,n_a}\), and the section-level naturality
  \ref{mk-lem:tensorBraiding_hom_sectionsMul} pushes \(\Gamma(\beta)\) through
  \(\mathrm{sectionsMul}\), swapping the two section factors; the transport
  \(\mathrm{sectionsCast}_{\,n_a+n_b = n_b+n_a}\) identifies the two sides. This is
  again the right-unit leg with the braiding pair (\ref{mk-lem:tensorPowAdd_comm} +
  \ref{mk-lem:tensorBraiding_hom_sectionsMul}) in place of the right-unitor pair.
\end{proof}

\section{Graded-ring and graded-module assembly}
\label{mk-sec:sgr_graded_assembly}

The components and multiplications of the previous two sections are assembled
into the graded ring and graded module through Mathlib's
\emph{external-direct-sum} graded API: the relevant inputs are families of
abelian groups indexed by \(\mathbb{N}\) together with multiplication maps
landing in the index-sum component, not submodules of a pre-existing ring.

\begin{lemma}[External-direct-sum graded semiring]
  \label{mk-lem:directSum_gsemiring_mathlib}
  \lean{DirectSum.GSemiring}
  \mathlibok
  \dcref{custom}

  Let \(A : \mathbb{N} \to \mathrm{Type}\) be a family of additive commutative
  monoids equipped with a graded multiplication \(A_i \times A_j \to A_{i+j}\) and
  a unit \(1 \in A_0\) satisfying graded associativity, two-sided unitality and
  distributivity --- but \emph{not} necessarily graded commutativity (the data of
  a \(\mathrm{GSemiring}\) on \(A\)). Then the external direct sum
  \(\bigoplus_i A_i\) is a (not necessarily commutative) semiring, with
  multiplication the bilinear extension of the graded multiplication and identity
  the image of \(1 \in A_0\). This is the general assembly; the commutative
  upgrade is \ref{mk-lem:directSum_gcommSemiring_mathlib}, which adds graded
  commutativity.
\end{lemma}

\begin{lemma}[External-direct-sum graded commutative semiring]
  \label{mk-lem:directSum_gcommSemiring_mathlib}
  \lean{DirectSum.GCommSemiring}
  \mathlibok
  \dcref{custom}

  Let \(A : \mathbb{N} \to \mathrm{Type}\) be a family of additive commutative
  monoids equipped with a graded multiplication
  \(A_i \times A_j \to A_{i+j}\) and a unit \(1 \in A_0\) satisfying graded
  associativity, two-sided unitality, distributivity, and graded commutativity
  (the data of a \(\mathrm{GCommSemiring}\) on \(A\)). Then the external direct
  sum \(\bigoplus_i A_i\) is a commutative semiring, with multiplication the
  bilinear extension of the graded multiplication and identity the image of
  \(1 \in A_0\).
\end{lemma}

\begin{lemma}[External-direct-sum graded module]
  \label{mk-lem:directSum_gmodule_mathlib}
  \lean{DirectSum.Gmodule}
  \mathlibok
  \dcref{custom}

  Let \(A : \mathbb{N} \to \mathrm{Type}\) carry a graded semiring structure as
  in \ref{mk-lem:directSum_gsemiring_mathlib} and let
  \(M : \mathbb{N} \to \mathrm{Type}\) be a family of additive commutative
  monoids with a graded scalar action \(A_i \times M_j \to M_{i+j}\) satisfying
  the graded module axioms (the data of a \(\mathrm{Gmodule}\) of \(M\) over
  \(A\)). Then the external direct sum \(\bigoplus_j M_j\) is a module over the
  semiring \(\bigoplus_i A_i\), with scalar multiplication the bilinear
  extension of the graded action.
\end{lemma}

\begin{lemma}[Graded monoid structure on the section components]\label{mk-lem:sectionGradedRing_gmonoid}
  \lean{AlgebraicGeometry.Scheme.Modules.sectionGradedRing_gmonoid}
  \leanok
  \dcref{custom}

  \uses{mk-lem:sectionMul_coherent, mk-lem:gradedMonoid_eq_of_cast,
    mk-def:sectionGradedGMul, mk-def:sectionGradedGOne}
  Let \(\mathcal{L}\) be an arbitrary sheaf of \(\mathcal{O}_X\)-modules. The degree
  family \(\mathrm{sectionDeg}\,\mathcal{L}\), \(m \mapsto
  \Gamma(X,\mathcal{L}^{\otimes m})\), equipped with the graded multiplication
  \(\mathrm{gMul}\) (\ref{mk-def:sectionGradedGMul}) and graded unit \(\mathrm{gOne}\)
  (\ref{mk-def:sectionGradedGOne}), is a graded monoid (a
  \(\mathrm{GradedMonoid.GMonoid}\) on \(\mathrm{sectionDeg}\,\mathcal{L}\)).
\end{lemma}

\begin{proof}
  \leanok
  \uses{mk-lem:sectionMul_coherent, mk-lem:gradedMonoid_eq_of_cast,
    mk-def:sectionGradedGMul, mk-def:sectionGradedGOne}
  The three graded-monoid axioms --- left unitality, right unitality and
  associativity of \(\mathrm{gMul}\) about \(\mathrm{gOne}\) --- are the first three
  \emph{hypothesis-free} clauses of \ref{mk-lem:sectionMul_coherent}, stated as
  transport-mediated equalities of sections. Each is converted into the required
  equality of dependent \(\mathrm{GradedMonoid}\) pairs by the bridge
  \ref{mk-lem:gradedMonoid_eq_of_cast}, which turns a transport-mediated equality into an
  equality of pairs. This assembles the \(\mathrm{GMonoid}\) on
  \(\mathrm{sectionDeg}\,\mathcal{L}\).
\end{proof}

\begin{lemma}[Graded semiring structure on the section components (general)]
\label{mk-lem:sectionGradedRing_gsemiring}
  \lean{AlgebraicGeometry.Scheme.Modules.sectionGradedRing_gsemiring}
  \leanok
  \dcref{custom}

  \uses{mk-def:sheafTensorPow, mk-lem:sheafTensorPow_add, mk-def:sectionMul,
    mk-lem:sectionMul_coherent, mk-def:sectionsCast, mk-lem:gradedMonoid_eq_of_cast,
    mk-lem:directSum_gsemiring_mathlib, mk-lem:moduleUnit_mathlib}
  Let \(\mathcal{L}\) be an \emph{arbitrary} sheaf of \(\mathcal{O}_X\)-modules on
  \(X\). The family of \(\Gamma(X,\mathcal{O}_X)\)-modules
  \(m \mapsto \Gamma(X, \mathcal{L}^{\otimes m})\) carries a graded
  \emph{semiring} structure (a \(\mathrm{GSemiring}\),
  \ref{mk-lem:directSum_gsemiring_mathlib}) whose degree-\((m,m')\) multiplication is
  the composite
  \[
    \Gamma(X, \mathcal{L}^{\otimes m}) \otimes_{\Gamma(X,\mathcal{O}_X)}
      \Gamma(X, \mathcal{L}^{\otimes m'})
    \xrightarrow{\ \cdot\ }
    \Gamma\bigl(X, \mathcal{L}^{\otimes m} \otimes_{\mathcal{O}_X}
      \mathcal{L}^{\otimes m'}\bigr)
    \xrightarrow{\ \Gamma(\mu_{m,m'})\ }
    \Gamma\bigl(X, \mathcal{L}^{\otimes (m+m')}\bigr),
  \]
  the section multiplication (\ref{mk-def:sectionMul}) followed by the tensor-power
  comparison isomorphism (\ref{mk-lem:sheafTensorPow_add}), with unit
  \(1 \in \Gamma(X, \mathcal{O}_X) = \Gamma(X, \mathcal{L}^{\otimes 0})\)
  (\ref{mk-lem:moduleUnit_mathlib}). Consequently
  \(\bigoplus_{m \ge 0} \Gamma(X, \mathcal{L}^{\otimes m})\) is a semiring. This
  structure exists for \emph{every} \(\mathcal{L}\) and is in general
  \emph{non-commutative} --- it is the free tensor algebra on
  \(\Gamma(X,\mathcal{L})\) --- because it uses only the unit and associativity
  clauses of \ref{mk-lem:sectionMul_coherent}, which hold without hypotheses; graded
  commutativity (the invertible upgrade \ref{mk-lem:sectionGradedRing_gcommSemiring})
  is \emph{not} part of this layer.
\end{lemma}

\begin{proof}
  \leanok
  \uses{mk-lem:sectionGradedRing_gmonoid, mk-lem:gradedMonoid_eq_of_cast, mk-def:sectionMul,
    mk-lem:directSum_gsemiring_mathlib}
  The structure is assembled field-for-field as \(\mathrm{TensorPower}\) builds its
  graded semiring in \texttt{Mathlib.LinearAlgebra.TensorPower.Basic}. The
  degreewise multiplication \(\Gamma(\mu_{m,m'}) \circ (\,\cdot\,)\) of the
  statement and the degree-\(0\) unit \(1 \in \Gamma(X, \mathcal{L}^{\otimes 0})\)
  provide the graded multiplication and unit. Their graded-monoid axioms --- left
  and right unitality and associativity --- are the first three (hypothesis-free)
  component-level identities of \ref{mk-lem:sectionMul_coherent}, each passed through
  the bridge \ref{mk-lem:gradedMonoid_eq_of_cast}, which converts a
  transport-mediated equality (\ref{mk-def:sectionsCast}) into the required equality
  of dependent pairs; the graded power operation is the default iterated product
  and contributes no further data. This yields the graded-monoid layer.

  The remaining semiring axioms are the bilinearity (left and right distributivity
  and annihilation by \(0\)) of the graded multiplication and the behaviour of the
  natural-number coercion. Bilinearity is immediate: the degreewise multiplication
  \(\Gamma(\mu_{m,m'}) \circ (\,\cdot\,)\) is \(\Gamma(X,\mathcal{O}_X)\)-bilinear,
  being the composite of the bilinear section multiplication (\ref{mk-def:sectionMul})
  with the linear comparison \(\Gamma(\mu_{m,m'})\)
  (\ref{mk-lem:sheafTensorPow_add}), so each side vanishes on \(0\) and distributes
  over sums. The natural-number coercion sends \(n\) to \(n \cdot 1\), the
  \(n\)-fold sum of the degree-\(0\) unit. These are precisely the data of a
  \(\mathrm{GSemiring}\) on \(m \mapsto \Gamma(X, \mathcal{L}^{\otimes m})\), so
  \ref{mk-lem:directSum_gsemiring_mathlib} endows
  \(\bigoplus_{m \ge 0} \Gamma(X, \mathcal{L}^{\otimes m})\) with its semiring
  structure. No commutativity clause enters.
\end{proof}

\begin{lemma}[A semiring structure on the section direct sum exists]\label{mk-lem:sectionGradedRing_semiring_nonempty}
  \lean{AlgebraicGeometry.Scheme.Modules.sectionGradedRing_semiring_nonempty}
  \leanok
  \dcref{custom}

  \uses{mk-lem:sectionGradedRing_gsemiring, mk-lem:directSum_gsemiring_mathlib}
  Let \(\mathcal{L}\) be an arbitrary sheaf of \(\mathcal{O}_X\)-modules. Then the
  external direct sum \(\bigoplus_{m \ge 0} \Gamma(X,\mathcal{L}^{\otimes m})\)
  admits a semiring structure; equivalently the type of semiring structures on it
  is nonempty.
\end{lemma}

\begin{proof}
  \leanok
  \uses{mk-lem:sectionGradedRing_gsemiring, mk-lem:directSum_gsemiring_mathlib}
  By \ref{mk-lem:sectionGradedRing_gsemiring} the family \(m \mapsto
  \Gamma(X,\mathcal{L}^{\otimes m})\) carries a \(\mathrm{GSemiring}\). Mathlib's
  external-direct-sum assembly (\ref{mk-lem:directSum_gsemiring_mathlib}, via
  \(\mathrm{DirectSum.toSemiring}\)) turns this graded data into a genuine semiring
  on \(\bigoplus_{m \ge 0} \Gamma(X,\mathcal{L}^{\otimes m})\); exhibiting this one
  structure witnesses the nonemptiness.
\end{proof}

\begin{lemma}[A commutative semiring structure on the section direct sum exists]\label{mk-lem:sectionGradedRing_commSemiring_nonempty}
  \lean{AlgebraicGeometry.Scheme.Modules.sectionGradedRing_commSemiring_nonempty}
  \leanok
  \dcref{custom}

  \uses{mk-lem:sectionGradedRing_gcommSemiring, mk-lem:directSum_gcommSemiring_mathlib}
  Let \(L\) be an \emph{invertible} sheaf of \(\mathcal{O}_X\)-modules. Then the
  external direct sum \(\bigoplus_{m \ge 0} \Gamma(X, L^{\otimes m})\) admits a
  commutative semiring structure.
\end{lemma}

\begin{proof}
  \leanok
  \uses{mk-lem:sectionGradedRing_gcommSemiring, mk-lem:directSum_gcommSemiring_mathlib}
  By \ref{mk-lem:sectionGradedRing_gcommSemiring} the family \(m \mapsto
  \Gamma(X, L^{\otimes m})\) carries a \(\mathrm{GCommSemiring}\) when \(L\) is
  invertible. The external-direct-sum assembly
  \ref{mk-lem:directSum_gcommSemiring_mathlib} turns this graded data into a genuine
  commutative semiring on \(\bigoplus_{m \ge 0} \Gamma(X, L^{\otimes m})\).
\end{proof}

\begin{lemma}[Graded commutative semiring for an invertible sheaf]\label{mk-lem:sectionGradedRing_gcommSemiring}
  \lean{AlgebraicGeometry.Scheme.Modules.sectionGradedRing_gcommSemiring}
  \leanok
  \dcref{01CV,custom}

  \uses{mk-def:isInvertible, mk-lem:sectionGradedRing_gsemiring, mk-lem:sheafTensorPow_add,
    mk-def:sectionMul, mk-lem:sectionMul_coherent, mk-def:sectionsCast,
    mk-lem:gradedMonoid_eq_of_cast, mk-lem:directSum_gcommSemiring_mathlib,
    mk-lem:moduleUnit_mathlib}
  \textit{Source: \cite{stacks-project}, "Sheaves of Modules", Tag 01CV (the graded
  ring structure on \(\bigoplus \Gamma(X,\mathcal{L}^{\otimes n})\)).}
  Let \(L_s\) be an \emph{invertible} sheaf --- a line bundle,
  \ref{mk-def:isInvertible} --- on \(X_s\). Then the graded semiring of
  \ref{mk-lem:sectionGradedRing_gsemiring} on
  \(m \mapsto \Gamma(X_s, L_s^{\otimes m})\) is graded \emph{commutative} (a
  \(\mathrm{GCommSemiring}\), \ref{mk-lem:directSum_gcommSemiring_mathlib}), with the
  same degree-\((m,m')\) multiplication
  \[
    \Gamma(X_s, L_s^{\otimes m}) \otimes_{\Gamma(X_s,\mathcal{O}_{X_s})}
      \Gamma(X_s, L_s^{\otimes m'})
    \xrightarrow{\ \cdot\ }
    \Gamma\bigl(X_s, L_s^{\otimes m} \otimes L_s^{\otimes m'}\bigr)
    \xrightarrow{\ \Gamma(\mu_{m,m'})\ }
    \Gamma\bigl(X_s, L_s^{\otimes (m+m')}\bigr)
  \]
  and unit \(1 \in \Gamma(X_s, \mathcal{O}_{X_s}) = \Gamma(X_s, L_s^{\otimes 0})\)
  (\ref{mk-lem:moduleUnit_mathlib}). Consequently
  \(R(X_s, L_s) = \bigoplus_{m \ge 0} \Gamma(X_s, L_s^{\otimes m})\) is a
  commutative semiring; this is the graded-ring structure underlying
  \ref{mk-def:sectionGradedRing}. Invertibility of \(L_s\) enters only here, through
  the commutativity clause; for a general sheaf only the non-commutative semiring
  \ref{mk-lem:sectionGradedRing_gsemiring} is available.
\end{lemma}

\begin{proof}
  \leanok
  \uses{mk-def:isInvertible, mk-lem:sectionGradedRing_gsemiring, mk-lem:sectionMul_coherent,
    mk-lem:gradedMonoid_eq_of_cast, mk-def:sectionsCast,
    mk-lem:directSum_gcommSemiring_mathlib}
    By \ref{mk-lem:sectionGradedRing_gsemiring} the family
  \(m \mapsto \Gamma(X_s, L_s^{\otimes m})\) is already a graded semiring, so it
  remains only to supply graded commutativity. Since \(L_s\) is invertible, the
  commutativity clause of \ref{mk-lem:sectionMul_coherent} holds; passing it through
  the bridge \ref{mk-lem:gradedMonoid_eq_of_cast} --- which converts the
  transport-mediated equality (\ref{mk-def:sectionsCast}) into an equality of
  dependent pairs --- yields the graded-commutativity axiom. Adjoining it to the
  graded-semiring data produces a \(\mathrm{GCommSemiring}\) on
  \(m \mapsto \Gamma(X_s, L_s^{\otimes m})\), so
  \ref{mk-lem:directSum_gcommSemiring_mathlib} endows
  \(\bigoplus_{m \ge 0} \Gamma(X_s, L_s^{\otimes m})\) with its commutative-semiring
  structure.
\end{proof}

\begin{definition}[Action comparison isomorphism for the twisted module]\label{mk-def:moduleTensorPowAdd}
  \lean{AlgebraicGeometry.Scheme.Modules.moduleTensorPowAdd}
  \leanok
  \dcref{custom}

  \uses{mk-cor:sheafTensorObjAssoc, mk-def:tensorBraiding, mk-def:sheafTensorObj,
    mk-def:sheafTensorPow, mk-lem:sheafTensorPow_add}
  Let \(\mathcal{F}\) be a sheaf of \(\mathcal{O}_X\)-modules and \(\mathcal{L}\) a
  line bundle. For \(i, j \in \mathbb{N}\) the \emph{action comparison
  isomorphism} is the isomorphism of sheaves of \(\mathcal{O}_X\)-modules
  \[
    \mathcal{L}^{\otimes i} \otimes_{\mathcal{O}_X}
      \bigl(\mathcal{F} \otimes_{\mathcal{O}_X} \mathcal{L}^{\otimes j}\bigr)
    \;\xrightarrow{\ \cong\ }\;
    \mathcal{F} \otimes_{\mathcal{O}_X} \mathcal{L}^{\otimes (i+j)}
  \]
  relating the left graded action of the section ring on the twisted module
  \(\mathcal{F}(j)\) (\ref{mk-def:sheafModuleTwist}) with the merge of the
  line-bundle factors. It is assembled from the sheaf-level associator
  (\ref{mk-cor:sheafTensorObjAssoc}), the braiding
  \(\beta_{\mathcal{L}^{\otimes i},\mathcal{F}}\) (\ref{mk-def:tensorBraiding})
  exchanging the two \emph{distinct} factors \(\mathcal{L}^{\otimes i}\) and
  \(\mathcal{F}\) --- a braiding of distinct objects, requiring \emph{no}
  invertibility hypothesis --- the whiskering isomorphisms of \(\otimes\)
  (\ref{mk-def:sheafTensorObj}), and the tensor-power comparison \(\mu_{i,j}\)
  merging \(\mathcal{L}^{\otimes i} \otimes_{\mathcal{O}_X} \mathcal{L}^{\otimes j}
  \cong \mathcal{L}^{\otimes (i+j)}\) (\ref{mk-lem:sheafTensorPow_add},
  \ref{mk-def:sheafTensorPow}). It underlies the graded-module
  structure of \ref{mk-lem:sectionGradedModule_gmodule}.
\end{definition}

\begin{lemma}[Associativity (hexagon) constraint for the action comparison]\label{mk-lem:moduleTensorPowAdd_assoc}
  \lean{AlgebraicGeometry.Scheme.Modules.moduleTensorPowAdd_assoc}
  \leanok
  \dcref{custom}

  \uses{mk-def:moduleTensorPowAdd, mk-cor:sheafTensorObjAssoc, mk-def:tensorBraiding,
    mk-lem:tensorPowAdd_assoc, mk-lem:sheafTensorPow_add}
  Let \(\mathcal{F}\) be a sheaf of \(\mathcal{O}_X\)-modules, \(\mathcal{L}\) a line
  bundle, and \(i, j, k \in \mathbb{N}\). Write \(a_{p,q}\) for the action comparison
  isomorphism (\ref{mk-def:moduleTensorPowAdd})
  \(\mathcal{L}^{\otimes p}\otimes(\mathcal{F}\otimes\mathcal{L}^{\otimes q})
  \xrightarrow{\cong}\mathcal{F}\otimes\mathcal{L}^{\otimes(p+q)}\). The two ways of
  contracting
  \(\mathcal{L}^{\otimes i}\otimes\bigl(\mathcal{L}^{\otimes j}\otimes
  (\mathcal{F}\otimes\mathcal{L}^{\otimes k})\bigr)\) down to
  \(\mathcal{F}\otimes\mathcal{L}^{\otimes(i+j+k)}\) agree:
  \[
    a_{i+j,\,k}\circ\bigl(\mu_{i,j}\rhd(\mathcal{F}\otimes\mathcal{L}^{\otimes k})\bigr)
      \circ \alpha^{-1}_{\mathcal{L}^{\otimes i},\,\mathcal{L}^{\otimes j},\,
        \mathcal{F}\otimes\mathcal{L}^{\otimes k}}
    \;=\;
    \mathrm{(reindex)}\circ a_{i,\,j+k}\circ
      \bigl(\mathcal{L}^{\otimes i}\lhd a_{j,k}\bigr),
  \]
  where \(\mu\) is the tensor-power comparison (\ref{mk-lem:sheafTensorPow_add}),
  \(\alpha\) the associator (\ref{mk-cor:sheafTensorObjAssoc}), \(\rhd, \lhd\) the
  canonical whiskerings, and the reindexing is \((i+j)+k = i+(j+k)\). One side merges
  the two line-bundle factors first and then acts once; the other acts twice in
  succession. This is the module analogue of the associativity constraint
  \ref{mk-lem:tensorPowAdd_assoc}, with the action comparison \(a\) in place of the
  ring comparison \(\mu\) on the two legs that touch \(\mathcal{F}\).
\end{lemma}

\begin{proof}
  \leanok
  \uses{mk-def:moduleTensorPowAdd, mk-cor:sheafTensorObjAssoc, mk-def:tensorBraiding,
    mk-lem:tensorPowAdd_assoc, mk-lem:sheafTensorPow_add}
    Unlike \ref{mk-lem:tensorPowAdd_assoc}, this identity is \emph{not} braiding-free:
  the action comparison \(a_{p,q}\) (\ref{mk-def:moduleTensorPowAdd}) contains the
  braiding \(\beta_{\mathcal{L}^{\otimes p},\mathcal{F}}\) (\ref{mk-def:tensorBraiding})
  sliding the line-bundle block \(\mathcal{L}^{\otimes p}\) past the module factor
  \(\mathcal{F}\). Since \(\mathcal{F}\) is \emph{not} assumed invertible the braiding
  does not collapse to the identity, so the statement is a genuine
  symmetric-monoidal \emph{hexagon} (Mac Lane), not a pentagon: on each side the
  block \(\mathcal{L}^{\otimes i}\) (and then \(\mathcal{L}^{\otimes j}\)) must be
  braided across \(\mathcal{F}\), and the identity records the two orders in which
  the merge \(\mu_{i,j}\) and the two braidings may be performed.

  Expand \(a_{p,q}\) on both sides into its constituents --- the associator
  (\ref{mk-cor:sheafTensorObjAssoc}), the braiding
  \(\beta_{\mathcal{L}^{\otimes p},\mathcal{F}}\) (\ref{mk-def:tensorBraiding}), the
  whiskerings, and the comparison \(\mu_{p,q}\) (\ref{mk-lem:sheafTensorPow_add}). The
  line-bundle-merge halves of the two sides are equated by the braiding-free pentagon
  \ref{mk-lem:tensorPowAdd_assoc} applied to the three factors
  \(\mathcal{L}^{\otimes i},\mathcal{L}^{\otimes j},\mathcal{L}^{\otimes k}\); the
  residual structural glue is a hexagon among canonical associators, whiskerings and
  the braiding \(\beta\) carrying \(\mathcal{L}^{\otimes i}\) (then
  \(\mathcal{L}^{\otimes j}\)) across \(\mathcal{F}\), discharged by coherence for
  symmetric monoidal categories.
\end{proof}

\begin{definition}[Twisted-section component in degree \(m\)]\label{mk-def:moduleSectionDeg}
  \lean{AlgebraicGeometry.Scheme.Modules.moduleSectionDeg}
  \leanok
  \dcref{custom}

  \uses{mk-def:sheafModuleTwist, mk-def:sheafTensorPow}
  Let \(\mathcal{F}\) be a sheaf of \(\mathcal{O}_X\)-modules and \(\mathcal{L}\) a
  sheaf of \(\mathcal{O}_X\)-modules. The \emph{degree-\(m\) twisted-section
  component} is the \(\Gamma(X,\mathcal{O}_X)\)-module of global sections of the
  \(m\)th twist of \(\mathcal{F}\),
  \[
    (\mathrm{moduleSectionDeg}\ \mathcal{F}\ \mathcal{L})_m
      \;:=\; \Gamma(X, \mathcal{F} \otimes_{\mathcal{O}_X} \mathcal{L}^{\otimes m})
      \;=\; (\mathcal{F} \otimes_{\mathcal{O}_X} \mathcal{L}^{\otimes m})(X)
        \qquad(\ref{mk-def:sheafModuleTwist},\ \ref{mk-def:sheafTensorPow}).
  \]
  It inherits its additive-group and \(\Gamma(X,\mathcal{O}_X)\)-module structure
  from the underlying module object
  \((\mathcal{F} \otimes_{\mathcal{O}_X} \mathcal{L}^{\otimes m})(X)\). These are the
  carriers \(m \mapsto (\mathrm{moduleSectionDeg}\ \mathcal{F}\ \mathcal{L})_m\) of
  the graded section module \(M(X_s, \mathcal{F}_s, L_s)\)
  (\ref{mk-lem:sectionGradedModule_gmodule}) and the domains of the index transport
  \ref{mk-def:moduleSectionsCast}. It is the module analogue of the degree-\(m\)
  section component \ref{mk-def:sectionDeg}, with the twisted module
  \(\mathcal{F} \otimes \mathcal{L}^{\otimes m}\) in place of the pure power
  \(\mathcal{L}^{\otimes m}\).
\end{definition}

\begin{definition}[Index-equality transport of twisted-section components]\label{mk-def:moduleSectionsCast}
  \lean{AlgebraicGeometry.Scheme.Modules.moduleSectionsCast}
  \leanok
  \dcref{custom}

  \uses{mk-def:moduleSectionDeg}
  Let \(\mathcal{F}\) be a sheaf of \(\mathcal{O}_X\)-modules, \(\mathcal{L}\) a sheaf
  of \(\mathcal{O}_X\)-modules, and let \(h : i = j\) be an equality of natural
  numbers. Applying the global-sections functor \(\Gamma(X,-)\) to the canonical
  isomorphism
  \(\mathcal{F} \otimes_{\mathcal{O}_X} \mathcal{L}^{\otimes i}
    \xrightarrow{\sim}
    \mathcal{F} \otimes_{\mathcal{O}_X} \mathcal{L}^{\otimes j}\)
  induced by \(h\) under the twist functor
  \(m \mapsto \mathcal{F} \otimes_{\mathcal{O}_X} \mathcal{L}^{\otimes m}\)
  (\ref{mk-def:moduleSectionDeg}) yields the \(\Gamma(X,\mathcal{O}_X)\)-linear
  isomorphism
  \[
    \mathrm{moduleSectionsCast}(h) :
    \Gamma(X, \mathcal{F} \otimes_{\mathcal{O}_X} \mathcal{L}^{\otimes i})
      \;\xrightarrow{\ \sim\ }\;
    \Gamma(X, \mathcal{F} \otimes_{\mathcal{O}_X} \mathcal{L}^{\otimes j}),
  \]
  the \emph{twisted-section-component transport} along \(h\). It is the
  twisted-module analogue of the pure section-component transport
  \ref{mk-def:sectionsCast}. Because the isomorphism induced by the reflexive equality
  is the identity and \(\Gamma(X,-)\) preserves identities, the transport along a
  reflexive index equality is the identity map (the reflexivity simplification
  \(\mathrm{moduleSectionsCast\_refl}\), \ref{mk-lem:moduleSectionsCast_refl}).
\end{definition}

\begin{lemma}[Reflexivity of the twisted-section transport]\label{mk-lem:moduleSectionsCast_refl}
  \lean{AlgebraicGeometry.Scheme.Modules.moduleSectionsCast_refl}
  \leanok
  \dcref{custom}

  \uses{mk-def:moduleSectionsCast}
  For any \(i\), the transport along the reflexive equality \(i = i\)
  (\ref{mk-def:moduleSectionsCast}) equals the identity automorphism of
  \(\Gamma(X, \mathcal{F} \otimes_{\mathcal{O}_X} \mathcal{L}^{\otimes i})\). This is
  the twisted-module counterpart of \ref{mk-lem:sectionsCast_refl}, and serves as a
  simplification lemma collapsing the transport along a reflexive index equality.
\end{lemma}

\begin{proof}
  \leanok\end{proof}

\begin{lemma}[Cast-mediated equality in the graded module]\label{mk-lem:moduleGradedMonoid_eq_of_cast}
  \lean{AlgebraicGeometry.Scheme.Modules.moduleGradedMonoid_eq_of_cast}
  \leanok
  \dcref{custom}

  \uses{mk-def:moduleSectionsCast, mk-lem:moduleSectionsCast_refl}
  Write the carrier of the graded module as the dependent sum
  \(\coprod_{n} \Gamma(X, \mathcal{F} \otimes_{\mathcal{O}_X} \mathcal{L}^{\otimes n})\),
  whose elements are pairs \((n, s)\) with
  \(s \in \Gamma(X, \mathcal{F} \otimes_{\mathcal{O}_X} \mathcal{L}^{\otimes n})\). Let
  \(a = (i, x)\) and \(b = (j, y)\) be two such pairs. If \(h : i = j\) and the
  transported component agrees, \(\mathrm{moduleSectionsCast}(h)\,x = y\), then
  \(a = b\). That is, a component-level equality mediated by the index transport
  \ref{mk-def:moduleSectionsCast} repackages into a genuine equality of dependent
  pairs. This is the twisted-module analogue of the ring-side bridge
  \ref{mk-lem:gradedMonoid_eq_of_cast}; it is the bridge that converts each
  transport-mediated clause of \ref{mk-lem:moduleSectionAction_coherent} into the
  dependent-pair form required by the \(\mathrm{Gmodule}\) axioms.
\end{lemma}

\begin{proof}
  \leanok

\end{proof}

\begin{definition}[Graded scalar action on twisted-section components]
  \label{mk-def:moduleSectionGradedGSMul}
  \lean{AlgebraicGeometry.Scheme.Modules.instGSMulNatSectionDegModuleSectionDeg}
  \dcref{custom}

  \uses{mk-def:moduleSectionDeg, mk-def:sectionDeg, mk-def:sectionMul, mk-def:moduleTensorPowAdd}
  The section components \ref{mk-def:sectionDeg} act on the twisted-section components
  \ref{mk-def:moduleSectionDeg} by a graded scalar multiplication
  \[
    (\mathrm{sectionDeg}\ \mathcal{L})_i \times
      (\mathrm{moduleSectionDeg}\ \mathcal{F}\ \mathcal{L})_j
    \;\longrightarrow\; (\mathrm{moduleSectionDeg}\ \mathcal{F}\ \mathcal{L})_{i+j},
    \qquad (r, x) \longmapsto r \star x,
  \]
  landing in the index-sum component (a \(\mathrm{GradedMonoid.GSMul}\) of
  \(\mathrm{moduleSectionDeg}\ \mathcal{F}\ \mathcal{L}\) over
  \(\mathrm{sectionDeg}\ \mathcal{L}\); the action degree is \(i+j\)). It is the section multiplication
  \(r \otimes x \mapsto r \cdot x\) (\ref{mk-def:sectionMul}) of the line-bundle power
  \(\mathcal{L}^{\otimes i}\) against the twisted module
  \(\mathcal{F} \otimes \mathcal{L}^{\otimes j}\), followed by global sections of the
  action comparison isomorphism \(a_{i,j}\) (\ref{mk-def:moduleTensorPowAdd}),
  \[
    \Gamma(X, \mathcal{L}^{\otimes i}) \otimes_{\Gamma(X,\mathcal{O}_X)}
      \Gamma(X, \mathcal{F} \otimes \mathcal{L}^{\otimes j})
    \xrightarrow{\ \cdot\ }
    \Gamma\bigl(X, \mathcal{L}^{\otimes i} \otimes
      (\mathcal{F} \otimes \mathcal{L}^{\otimes j})\bigr)
    \xrightarrow{\ \Gamma(a_{i,j})\ }
    \Gamma(X, \mathcal{F} \otimes \mathcal{L}^{\otimes (i+j)}).
  \]
  It is the module analogue of the graded multiplication \ref{mk-def:sectionGradedGMul},
  with the action comparison \(a_{i,j}\) in place of the ring comparison
  \(\mu_{i,j}\) and the twisted module in the right-hand factor.
\end{definition}

\begin{lemma}[Left-unit (base-case) constraint for the action comparison]\label{mk-lem:moduleTensorPowAdd_zero_left}
  \lean{AlgebraicGeometry.Scheme.Modules.moduleTensorPowAdd_zero_left}
  \leanok
  \dcref{custom}

  \uses{mk-def:moduleTensorPowAdd, mk-lem:tensorPowAdd_zero_left, mk-def:tensorBraiding,
    mk-def:tensorObjUnitIso, mk-lem:tensorObjUnitIso_eq, mk-cor:sheafTensorObjAssoc}
  Let \(\mathcal{F}\) be a sheaf of \(\mathcal{O}_X\)-modules, \(\mathcal{L}\) a line
  bundle, and \(k \in \mathbb{N}\). The degree-\((0,k)\) action comparison
  isomorphism \(a_{0,k}\) (\ref{mk-def:moduleTensorPowAdd}) is the left unitor,
  reindexed along \(0 + k = k\):
  \[
    a_{0,k} \;=\;
      \mathrm{tensorObjUnitIso}_{\mathcal{F}\otimes\mathcal{L}^{\otimes k}}
      \;:\; \mathbf{1}_X \otimes_{\mathcal{O}_X}
        (\mathcal{F}\otimes_{\mathcal{O}_X}\mathcal{L}^{\otimes k})
      \;\xrightarrow{\ \sim\ }\;
      \mathcal{F}\otimes_{\mathcal{O}_X}\mathcal{L}^{\otimes k}
  \]
  (\ref{mk-def:tensorObjUnitIso}, transport of the left unitor \(\lambda\)). This is
  the base case of the action comparison, the module analogue of the left-unit
  constraint \ref{mk-lem:tensorPowAdd_zero_left}, and is the key input to the
  Unitality clause of \ref{mk-lem:moduleSectionAction_coherent}.
\end{lemma}

\begin{proof}
  \leanok
  \uses{mk-def:moduleTensorPowAdd, mk-lem:tensorPowAdd_zero_left, mk-def:tensorBraiding,
    mk-def:tensorObjUnitIso, mk-lem:tensorObjUnitIso_eq, mk-cor:sheafTensorObjAssoc}
    Expand \(a_{0,k}\) into its constituents (\ref{mk-def:moduleTensorPowAdd}): the
  sheaf-level associator (\ref{mk-cor:sheafTensorObjAssoc}), the braiding
  \(\beta_{\mathcal{L}^{\otimes 0},\mathcal{F}}\) (\ref{mk-def:tensorBraiding}), the
  whiskerings, and the line-bundle merge \(\mu_{0,k}\). In degree \(p=0\) the
  line-bundle block collapses, \(\mathcal{L}^{\otimes 0}=\mathbf{1}_X\), so the merge
  reduces to the left-unit constraint \(\mu_{0,k}=\lambda\) of
  \ref{mk-lem:tensorPowAdd_zero_left}, and the braiding
  \(\beta_{\mathcal{L}^{\otimes 0},\mathcal{F}}=\beta_{\mathbf{1}_X,\mathcal{F}}\)
  becomes the unit-braiding, the composite \(\lambda_{\mathcal{F}} \circ
  \rho_{\mathcal{F}}^{-1}\) of the left unitor with the inverse right unitor on the
  unit object.

  The verification requires the unit-braiding coherence: the monoidal unit object of
  \(X.\Modules\) is the structure sheaf \(\mathbf{1}_X = \mathcal{O}_X\), and the
  unit-braiding identity \(\beta_{\mathbf{1},\mathcal{F}} = \lambda_{\mathcal{F}} \circ
  \rho_{\mathcal{F}}^{-1}\) is established over the formal monoidal unit \(\mathbf{1}\)
  by symmetric-monoidal coherence (Mac Lane), using the equality of unitors on the
  unit object (\(\mathrm{MonoidalCategory.unitors\_equal}\), descended through
  sheafification as in \ref{mk-lem:tensorPowAdd_zero_left}), and then applied at
  \(\mathbf{1}_X\) via the definitional equality of the two unit presentations.
  The bridge \ref{mk-lem:tensorObjUnitIso_eq} then identifies the left-unitor with
  \(\mathrm{tensorObjUnitIso}\), and the reindexing cast along \(0+k=k\) absorbs
  into it, yielding
  \(a_{0,k}=\mathrm{tensorObjUnitIso}_{\mathcal{F}\otimes\mathcal{L}^{\otimes k}}\).
\end{proof}

\begin{lemma}[Coherence of the module action]
  \label{mk-lem:moduleSectionAction_coherent}
  \lean{}
  \lean{AlgebraicGeometry.Scheme.Modules.moduleSectionAction_one_smul,
    AlgebraicGeometry.Scheme.Modules.moduleSectionAction_mul_smul}
  \dcref{custom}

  \uses{mk-lem:moduleTensorPowAdd_assoc, mk-lem:moduleTensorPowAdd_zero_left,
    mk-lem:tensorBraiding_hom_sectionsMul,
    mk-lem:tensorObjAssoc_hom_sectionsMul, mk-lem:sectionsMul_whiskerRight_natural,
    mk-lem:sectionsMul_whiskerLeft_natural, mk-lem:tensorObjUnitIso_hom_sectionsMul,
    mk-def:sectionMul, mk-def:moduleTensorPowAdd, mk-def:sectionsCast, mk-lem:sheafTensorPow_add}
  Let \(\mathcal{F}\) be a sheaf of \(\mathcal{O}_X\)-modules and \(\mathcal{L}\) a
  line bundle. For \(r\in\Gamma(X,\mathcal{L}^{\otimes i})\) and
  \(x\in\Gamma(X,\mathcal{F}\otimes\mathcal{L}^{\otimes k})\) write
  \[
    r\star x \;=\; \Gamma(a_{i,k})\bigl(\mathrm{sectionsMul}(r\otimes x)\bigr)
      \;\in\;\Gamma(X,\mathcal{F}\otimes\mathcal{L}^{\otimes(i+k)})
  \]
  for the degreewise action: the section multiplication (\ref{mk-def:sectionMul})
  followed by global sections of the action comparison \(a_{i,k}\)
  (\ref{mk-def:moduleTensorPowAdd}). Then, as transport-mediated equalities of sections
  (the subscript on \(\mathrm{sectionsCast}\), \ref{mk-def:sectionsCast}, recording the
  index equality that matches the two a-priori-distinct section modules):
  \begin{itemize}
    \item \emph{Unitality.} For \(x\in\Gamma(X,\mathcal{F}\otimes
      \mathcal{L}^{\otimes k})\) and \(1\) the degree-\(0\) unit of the section ring,
      \[ \mathrm{sectionsCast}_{\,0+k=k}\,(1\star x) = x. \]
    \item \emph{Compatibility.} For \(r\in\Gamma(X,\mathcal{L}^{\otimes i})\),
      \(r'\in\Gamma(X,\mathcal{L}^{\otimes j})\) and
      \(x\in\Gamma(X,\mathcal{F}\otimes\mathcal{L}^{\otimes k})\),
      \[ \mathrm{sectionsCast}_{\,(i+j)+k=i+(j+k)}\,\bigl((r\cdot r')\star x\bigr)
        = r\star(r'\star x), \]
      where \(r\cdot r'\) is the section ring multiplication (\ref{mk-def:sectionMul}
      followed by \(\Gamma(\mu_{i,j})\)).
  \end{itemize}
  This is the module analogue of the unit and associativity clauses of
  \ref{mk-lem:sectionMul_coherent}, with the action comparison \(a\) replacing the
  ring comparison \(\mu\) on the legs touching \(\mathcal{F}\).
\end{lemma}

\begin{proof}
  \uses{mk-lem:moduleTensorPowAdd_assoc, mk-lem:moduleTensorPowAdd_zero_left,
    mk-lem:tensorBraiding_hom_sectionsMul,
    mk-lem:tensorObjAssoc_hom_sectionsMul, mk-lem:sectionsMul_whiskerRight_natural,
    mk-lem:sectionsMul_whiskerLeft_natural, mk-lem:tensorObjUnitIso_hom_sectionsMul,
    mk-def:sectionsCast, mk-lem:sheafTensorPow_add}
  Both identities are the section-level shadow of the iso-level constraints, pushed
  through \(\mathrm{sectionsMul}\) by the naturality helpers --- exactly as in
  \ref{mk-lem:sectionMul_coherent}, but with the action comparison \(a\) in place of
  the ring comparison \(\mu\) on the legs that touch \(\mathcal{F}\).

  \emph{Unitality.} Unfold \(1\star x\). The base case \(a_{0,k}\) of the action
  comparison (\ref{mk-def:moduleTensorPowAdd}) is the left unitor
  \(\lambda_{\mathcal{F}\otimes\mathcal{L}^{\otimes k}}\) reindexed along \(0+k=k\),
  by \ref{mk-lem:moduleTensorPowAdd_zero_left} (where the unit-braiding coherence is
  established; see that proof). Granting that identification, the inner
  reindexing cast pairs with the outer \(\mathrm{sectionsCast}_{\,0+k=k}\) and the two
  cancel; \ref{mk-lem:tensorObjUnitIso_hom_sectionsMul} then closes the leg.

  \emph{Compatibility.} Unfold both sides to composites
  \(\Gamma(a)\circ\mathrm{sectionsMul}\). On the left, \((r\cdot r')\star x\) hides
  the ring comparison \(\mu_{i,j}\) inside the first factor of the outer
  \(\mathrm{sectionsMul}\); the right-whisker naturality
  \ref{mk-lem:sectionsMul_whiskerRight_natural} (applied to \(f=\mu_{i,j}\),
  \(G=\mathcal{F}\otimes\mathcal{L}^{\otimes k}\)) slides \(\Gamma(\mu_{i,j})\) out.
  On the right, \(r\star(r'\star x)\) hides the action comparison \(a_{j,k}\) inside
  the second factor; the left-whisker naturality
  \ref{mk-lem:sectionsMul_whiskerLeft_natural} (applied to \(g=a_{j,k}\),
  \(F=\mathcal{L}^{\otimes i}\)) slides \(\Gamma(a_{j,k})\) out. Both sides are now
  \(\Gamma(\text{comparison composite})\) applied to a bare iterated product,
  \(\mathrm{sectionsMul}(\mathrm{sectionsMul}(r\otimes r')\otimes x)\) on the left and
  \(\mathrm{sectionsMul}(r\otimes\mathrm{sectionsMul}(r'\otimes x))\) on the right.
  The iso-level hexagon \ref{mk-lem:moduleTensorPowAdd_assoc} equates the two
  comparison composites modulo the associator \(\alpha\) and the braiding \(\beta\).
  The residual \(\Gamma(\alpha)\) is pushed through the iterated
  \(\mathrm{sectionsMul}\) by \ref{mk-lem:tensorObjAssoc_hom_sectionsMul},
  reassociating the three section factors, and the residual
  \(\Gamma(\beta_{\mathcal{L}^{\otimes i},\mathcal{F}})\) is pushed through by
  \ref{mk-lem:tensorBraiding_hom_sectionsMul} --- stated precisely for the two
  \emph{distinct} factors \(\mathcal{L}^{\otimes i}\) and \(\mathcal{F}\) --- swapping
  those two section factors. The transport
  \(\mathrm{sectionsCast}_{\,(i+j)+k=i+(j+k)}\) absorbs the reindexing and identifies
  the two sides in a single module. This is the associativity leg of
  \ref{mk-lem:sectionMul_coherent} with the action hexagon and braiding slide in place
  of the pentagon and the right-unitor pair.
\end{proof}

\begin{lemma}[Graded module structure on the twisted-section components]\label{mk-lem:sectionGradedModule_gmodule}
  \lean{AlgebraicGeometry.Scheme.Modules.sectionGradedModule_gmodule}
  \leanok
  \dcref{01CV,custom}

  \uses{mk-def:sheafModuleTwist, mk-lem:sheafTensorPow_add, mk-def:sectionMul,
    mk-lem:moduleSectionAction_coherent, mk-lem:moduleTensorPowAdd_assoc,
    mk-def:moduleSectionDeg, mk-def:moduleSectionsCast, mk-def:moduleSectionGradedGSMul,
    mk-def:sectionsCast, mk-lem:moduleGradedMonoid_eq_of_cast,
    mk-lem:directSum_gmodule_mathlib, mk-lem:sectionGradedRing_gsemiring,
    mk-def:moduleTensorPowAdd, mk-def:tensorBraiding}
  \textit{Source: \cite{stacks-project}, "Sheaves of Modules", Tag 01CV (the graded
  module \(\Gamma_*(X,\mathcal{L},\mathcal{F})\)).}
  This lemma establishes the \(\mathbb{N}\)-graded (\(n\ge 0\)) analogue of Tag 01CV: the
  Stacks source assembles the \(\mathbb{Z}\)-graded module over an \emph{invertible}
  \(\mathcal{L}\) (using \(\mathcal{L}^\vee\) for negative degrees), whereas the
  action comparison \(a_{i,j}\) (\ref{mk-def:moduleTensorPowAdd}) exists for an
  arbitrary sheaf of modules \(\mathcal{F}\) and an arbitrary \(\mathcal{L}\), yielding
  the \(n\ge 0\) truncation \(\bigoplus_{m\ge 0}\Gamma(X,\mathcal{F}\otimes
  \mathcal{L}^{\otimes m})\) without any invertibility hypothesis.
  Let \(\mathcal{F}_s\) be a coherent sheaf and \(L_s\) a line bundle on
  \(X_s\). The family
  \(m \mapsto \Gamma(X_s, \mathcal{F}_s \otimes L_s^{\otimes m})\)
  (\ref{mk-def:sheafModuleTwist}) carries a graded module structure (a
  \(\mathrm{Gmodule}\), \ref{mk-lem:directSum_gmodule_mathlib}) over the graded
  semiring of \ref{mk-lem:sectionGradedRing_gsemiring}, with degree-\((m,m')\)
  action
  \[
    \Gamma(X_s, L_s^{\otimes m}) \otimes_{\Gamma(X_s,\mathcal{O}_{X_s})}
      \Gamma(X_s, \mathcal{F}_s \otimes L_s^{\otimes m'})
    \xrightarrow{\ \cdot\ }
    \Gamma\bigl(X_s, L_s^{\otimes m} \otimes
      (\mathcal{F}_s \otimes L_s^{\otimes m'})\bigr)
    \xrightarrow{\ \sim\ }
    \Gamma\bigl(X_s, \mathcal{F}_s \otimes L_s^{\otimes (m+m')}\bigr),
  \]
  the section multiplication followed by the comparison isomorphism that
  reassociates and merges the line-bundle factors (using
  \ref{mk-lem:sheafTensorPow_add} and the symmetry of the tensor product).
  Consequently \(M(X_s, \mathcal{F}_s, L_s) = \bigoplus_{m \ge 0}
  \Gamma(X_s, \mathcal{F}_s \otimes L_s^{\otimes m})\) is a graded module over
  \(R(X_s, L_s)\); this is the graded-module structure underlying
  \ref{mk-def:sectionGradedModule}.
\end{lemma}

\begin{proof}
  \leanok
  \uses{mk-lem:moduleSectionAction_coherent, mk-lem:moduleGradedMonoid_eq_of_cast,
    mk-def:moduleSectionDeg, mk-def:moduleSectionsCast, mk-def:moduleSectionGradedGSMul,
    mk-lem:directSum_gmodule_mathlib, mk-lem:sectionGradedRing_gsemiring,
    mk-lem:sectionGradedRing_gmonoid}
    Assemble the \(\mathrm{Gmodule}\) field-for-field over the graded semiring of
  \ref{mk-lem:sectionGradedRing_gsemiring}, exactly as
  \ref{mk-lem:sectionGradedRing_gmonoid} assembles the graded monoid. The action grades
  by ordinary addition on \(\mathbb{N}\) (degree-\((m,m')\) action landing in degree
  \(m+m'\)), with the degreewise action \(r\star x\) of
  \ref{mk-lem:moduleSectionAction_coherent}. Its two \(\mathrm{Gmodule}\) axioms ---
  unitality \(1 \cdot x = x\) and compatibility \(r \cdot (r' \cdot x) = (rr')\cdot x\)
  with the ring multiplication --- are the unitality and compatibility clauses of
  \ref{mk-lem:moduleSectionAction_coherent}, each turned from its transport-mediated
  form into the required equality of dependent pairs by the module bridge
  \ref{mk-lem:moduleGradedMonoid_eq_of_cast}. Bilinearity (additivity and annihilation by
  \(0\) in each variable) is free, as in \ref{mk-lem:sectionGradedRing_gsemiring}. By
  \ref{mk-lem:directSum_gmodule_mathlib} this produces the module structure of
  \(M(X_s,\mathcal{F}_s,L_s)\) over \(R(X_s,L_s)\).
\end{proof}

\section{P0.5: the bridge from local triviality to invertible-graded sheaves}
\label{mk-sec:sgr_invertibleGr_bridge}

\ref{mk-def:isInvertible} defines invertibility of a sheaf of modules through
a \emph{trivializing basis}; the line bundles that occur in practice (the
Serre twists \(\mathcal{O}(m)\) of \ref{mk-chap:Picard_QuotScheme}) are given
instead as \emph{locally trivial} sheaves in the geometric sense of
\(\mathrm{IsLocallyTrivial}\) (\ref{mk-def:IsLocallyTrivial}: every point has
an affine neighbourhood on which the sheaf restricts to the structure
sheaf). This section supplies the missing bridge between the two notions,
unlocking the graded-commutative section ring
\ref{mk-lem:sectionGradedRing_gcommSemiring} for every locally trivial line
bundle.

\begin{lemma}
[Section-level invertibility from a chart trivialisation]
  \label{mk-lem:isInvertible_of_restrict_iso}
  \lean{AlgebraicGeometry.Scheme.LineBundle.isInvertible_of_restrict_iso}
  \leanok
  \dcref{custom}
  \uses{mk-def:IsLocallyTrivial}
  Let \(M\) be a sheaf of \(\mathcal{O}_X\)-modules and \(V \subseteq X\)
  an open with a trivialisation \(M|_V \cong \mathcal{O}_V\). Then
  \(\Gamma(M, V)\) is an invertible \(\Gamma(X, V)\)-module.
\end{lemma}
\begin{proof}
  \leanok
  Evaluating the trivialisation at the top open of \(V\) gives a
  \(\Gamma(X,V)\)-linear isomorphism \(\Gamma(M, V) \cong \Gamma(X, V)\)
  (using \(\Gamma(\mathcal{O}_V, \top) = \Gamma(X, V)\)), exhibiting
  \(\Gamma(M,V)\) as a free rank-one, hence invertible, \(\Gamma(X,V)\)-module.
\end{proof}

\begin{theorem}
[P0.5: local triviality implies invertible-graded]
  \label{mk-thm:isLocallyTrivial_isInvertibleGr}
  \lean{AlgebraicGeometry.Scheme.LineBundle.IsLocallyTrivial.isInvertibleGr}
  \leanok
  \source{stacks-modules}
  \uses{mk-def:IsLocallyTrivial, mk-def:isInvertible, mk-lem:isInvertible_of_restrict_iso,
    mk-cor:IsLocallyTrivial_exists_affine_trivializing_le}
  \dcref{custom}
  If \(M\) is locally trivial (\ref{mk-def:IsLocallyTrivial}), then \(M\)
  is invertible in the sense of \ref{mk-def:isInvertible}.
\end{theorem}
\begin{proof}
  \leanok
  The trivialising opens \(\{V : M|_V \cong \mathcal{O}_V\}\) form a
  basis of the topology of \(X\): every neighbourhood of every point
  contains an affine trivialising chart
  (\ref{mk-cor:IsLocallyTrivial_exists_affine_trivializing_le}). On each
  such \(V\), \(\Gamma(M,V)\) is an invertible \(\Gamma(X,V)\)-module by
  \ref{mk-lem:isInvertible_of_restrict_iso}. Together these are exactly
  the trivializing-basis data required by \ref{mk-def:isInvertible}.
\end{proof}

\begin{corollary}
[The Serre twist family is invertible-graded]
  \label{mk-cor:serreTwist_isInvertibleGr}
  \lean{AlgebraicGeometry.ProjTwist.serreTwistGlued_isInvertibleGr,
    AlgebraicGeometry.ProjTwist.serreTwist_isInvertibleGr,
    AlgebraicGeometry.ProjTwist.twistingSheaf_isInvertibleGr}
  \leanok
  \dcref{custom}
  \uses{mk-thm:isLocallyTrivial_isInvertibleGr, mk-lem:serreTwist_isLocallyTrivial}
  Every member of the Serre twist family --- the twist on the glued
  model, on the integral model \(\operatorname{Proj}\mathbb{Z}[X_i]\),
  and on the relative projective space \(\mathbb{P}(n; S)\)
  (\ref{mk-lem:serreTwist_isLocallyTrivial}: each is locally trivial) ---
  is invertible in the sense of \ref{mk-def:isInvertible}. Consequently
  \ref{mk-lem:sectionGradedRing_gcommSemiring} applies unconditionally to
  each: its section ring \(\bigoplus_{m \ge 0} \Gamma(-, \mathcal{O}(m))\)
  is graded-commutative.
\end{corollary}
\begin{proof}
  \leanok
  Immediate instances of \ref{mk-thm:isLocallyTrivial_isInvertibleGr}
  applied to \ref{mk-lem:serreTwist_isLocallyTrivial}.
\end{proof}

\section{D5: finite-dimensional graded pieces from finite generation}
\label{mk-sec:sgr_graded_pieces_finite}

The Hilbert-series rationality induction \(\ref{mk-lem:graded_subquotient_isRatHilb}\)
(Stacks Tag 00K1) assumes each graded piece \(\mathcal{M}_n\) of the ambient
graded module is finite-dimensional over the base field \(\kappa\). This
section discharges that hypothesis from finite generation alone, for a
module graded by a finite family of commuting degree-raising endomorphisms
(\ref{mk-def:graded_raisesDegree}) --- the abstraction, used throughout
\ref{mk-chap:Picard_QuotScheme}, of the degree-one generators \(X_1,
\dots, X_r\) of a polynomial grading ring acting on \(M\).

\begin{theorem}
[D5: finite generation forces finite-dimensional graded pieces]
  \label{mk-thm:graded_finiteDimensional_of_polyModule_finite}
  \lean{AlgebraicGeometry.GradedModule.finiteDimensional_of_polyModule_finite}
  \leanok
  \source{hilbert-serre-algebra}
  \uses{mk-def:graded_polyModule, mk-def:graded_raisesDegree}
  \dcref{custom}
  Let \(M = \bigoplus_n \mathcal{M}_n\) be an \(\mathbb{N}\)-graded
  \(\kappa\)-vector space and \(t \colon \mathrm{Fin}\,r \to
  \mathrm{End}_\kappa(M)\) a finite family of pairwise commuting
  endomorphisms, each raising degree by one
  (\ref{mk-def:graded_raisesDegree}), making \(M\) a module over
  \(\kappa[X_1, \dots, X_r]\) with \(X_i\) acting as \(t_i\)
  (\ref{mk-def:graded_polyModule}). If \(M\) is finitely generated over
  this polynomial ring, then every graded piece \(\mathcal{M}_n\) is
  finite-dimensional over \(\kappa\).
\end{theorem}
\begin{proof}
  \leanok
  A monomial \(X^\alpha\) of total degree \(|\alpha|\) shifts a vector
  of \(\mathcal{M}_d\) into \(\mathcal{M}_{d+|\alpha|}\) (induction on
  \(|\alpha|\) via the degree-raising hypothesis on each \(t_i\)); by
  linearity, a homogeneous polynomial of degree \(e\) shifts
  \(\mathcal{M}_d\) into \(\mathcal{M}_{d+e}\). Fix a finite generating
  set of \(M\) over \(\kappa[X_1,\dots,X_r]\) and replace each generator
  by its (finitely many) homogeneous components \(h \in \mathcal{M}_{d}\);
  these homogeneous pieces still generate \(M\). For \(p \in
  \kappa[X_1,\dots,X_r]\) and a homogeneous generator \(h \in
  \mathcal{M}_d\), decomposing \(p\) into its homogeneous components
  shows that the degree-\(n\) part of \(p \cdot h\) is
  \((\mathrm{hc}_{n-d}\,p) \cdot h\), the action of the degree-\((n-d)\)
  homogeneous component of \(p\); since a homogeneous polynomial of
  degree \(\le n\) ranges over a \emph{finite-dimensional} space of
  coefficients (in the finitely many variables \(X_1,\dots,X_r\) of
  degree \(\le n\)), the degree-\(n\) part of \(p \cdot h\) lies in the
  image of a fixed finite-dimensional space under \(\kappa\)-linear
  action of \(h\). Every \(x \in \mathcal{M}_n\) is (via the finite
  generating set and the graded-piece projection) a finite sum of such
  contributions, one for each of the finitely many homogeneous
  generators; hence \(\mathcal{M}_n\) is contained in a finite sum of
  finite-dimensional subspaces, so is itself finite-dimensional.
\end{proof}

\begin{remark}
[Feeding the Hilbert-series induction]
  \label{mk-rem:graded_pieces_finite_consumer}
  \dcref{custom}
  \uses{mk-thm:graded_finiteDimensional_of_polyModule_finite, mk-lem:graded_subquotient_isRatHilb}
  Registering \(\ref{mk-thm:graded_finiteDimensional_of_polyModule_finite}\)'s
  conclusion as the instance \([\forall n,
  \mathrm{FiniteDimensional}\,\kappa\,(\mathcal{M}_n)]\) discharges the
  finiteness hypothesis of the Hilbert-series rationality induction
  \(\ref{mk-lem:graded_subquotient_isRatHilb}\) from finite generation
  alone, with no further geometric input.
\end{remark}
